\documentclass[11pt,oneside,openany]{book}
\usepackage[T1]{fontenc}
\usepackage[utf8]{inputenc}
\usepackage{lmodern}
\usepackage[a4paper,inner=28mm,outer=25mm,top=26mm,bottom=27mm,headheight=15pt]{geometry}
\usepackage{amsmath,amssymb,amsthm,mathtools,bm,mathrsfs}
\usepackage{microtype,booktabs,longtable,array,tabularx,enumitem,xcolor}
\usepackage{fancyhdr}
\usepackage[breaklinks=true,colorlinks=true,linkcolor=black,citecolor=black,urlcolor=black,pdfusetitle]{hyperref}
\usepackage{aliascnt}
\usepackage[nameinlink,noabbrev]{cleveref}
\setlength{\emergencystretch}{2em}
\setlength{\parskip}{3pt}
\setlength{\parindent}{1.2em}
\setcounter{tocdepth}{1}
\setcounter{secnumdepth}{2}
\raggedbottom
\allowdisplaybreaks[2]
\pagestyle{fancy}
\fancyhf{}
\fancyhead[L]{\small\nouppercase{\leftmark}}
\fancyfoot[C]{\thepage}
\renewcommand{\headrulewidth}{0.3pt}
\renewcommand{\chaptermark}[1]{\markboth{\thechapter.\ #1}{}}
\fancypagestyle{plain}{\fancyhf{}\fancyfoot[C]{\thepage}\renewcommand{\headrulewidth}{0pt}}
\newtheorem{theorem}{Theorem}[chapter]
\newaliascnt{proposition}{theorem}
\newtheorem{proposition}[proposition]{Proposition}
\aliascntresetthe{proposition}
\newaliascnt{lemma}{theorem}
\newtheorem{lemma}[lemma]{Lemma}
\aliascntresetthe{lemma}
\newaliascnt{corollary}{theorem}
\newtheorem{corollary}[corollary]{Corollary}
\aliascntresetthe{corollary}
\newaliascnt{criterion}{theorem}
\newtheorem{criterion}[criterion]{Conditional criterion}
\aliascntresetthe{criterion}
\theoremstyle{definition}
\newaliascnt{definition}{theorem}
\newtheorem{definition}[definition]{Definition}
\aliascntresetthe{definition}
\newaliascnt{certificate}{theorem}
\newtheorem{certificate}[certificate]{Reported finite certificate}
\aliascntresetthe{certificate}
\newaliascnt{hypothesis}{theorem}
\newtheorem{hypothesis}[hypothesis]{Hypothesis}
\aliascntresetthe{hypothesis}
\newaliascnt{problem}{theorem}
\newtheorem{problem}[problem]{Remaining problem}
\aliascntresetthe{problem}
\theoremstyle{remark}
\newaliascnt{remark}{theorem}
\newtheorem{remark}[remark]{Remark}
\aliascntresetthe{remark}
\newaliascnt{scope}{theorem}
\newtheorem{scope}[scope]{Scope}
\aliascntresetthe{scope}
\crefname{theorem}{theorem}{theorems}
\crefname{proposition}{proposition}{propositions}
\crefname{lemma}{lemma}{lemmas}
\crefname{corollary}{corollary}{corollaries}
\crefname{criterion}{criterion}{criteria}
\crefname{definition}{definition}{definitions}
\crefname{hypothesis}{hypothesis}{hypotheses}
\crefname{certificate}{finite certificate}{finite certificates}
\Crefname{certificate}{Finite certificate}{Finite certificates}
\crefname{problem}{remaining problem}{remaining problems}
\crefname{remark}{remark}{remarks}
\crefname{scope}{scope statement}{scope statements}
\newcommand{\R}{\mathbb R}
\newcommand{\C}{\mathbb C}
\newcommand{\N}{\mathbb N}
\newcommand{\Z}{\mathbb Z}
\newcommand{\T}{\mathbb T}
\newcommand{\E}{\mathbb E}
\newcommand{\cH}{\mathcal H}
\newcommand{\cA}{\mathcal A}
\newcommand{\cG}{\mathcal G}
\newcommand{\cK}{\mathcal K}
\newcommand{\one}{\mathbf1}
\newcommand{\dd}{\mathrm d}
\newcommand{\ii}{\mathrm i}
\newcommand{\Tr}{\operatorname{Tr}}
\newcommand{\tr}{\operatorname{tr}}
\newcommand{\diag}{\operatorname{diag}}
\newcommand{\Ran}{\operatorname{Ran}}
\newcommand{\Ker}{\operatorname{Ker}}
\newcommand{\Var}{\operatorname{Var}}
\newcommand{\Cov}{\operatorname{Cov}}
\newcommand{\supp}{\operatorname{supp}}
\newcommand{\dist}{\operatorname{dist}}
\newcommand{\osc}{\operatorname{osc}}
\newcommand{\gap}{\operatorname{gap}}
\newcommand{\SU}{\mathrm{SU}}
\newcommand{\OS}{\mathrm{OS}}
\newcommand{\op}{\mathrm{op}}
\newcommand{\rank}{\operatorname{rank}}
\newcommand{\norm}[1]{\left\lVert#1\right\rVert}
\newcommand{\abs}[1]{\left\lvert#1\right\rvert}
\newcommand{\ip}[2]{\left\langle#1,#2\right\rangle}
\newcommand{\avg}[1]{\left\langle#1\right\rangle_{\T^4}}
\newcommand{\src}[2]{\cite[#2]{#1}}
\newcommand{\sourcebasis}[1]{\par\smallskip\noindent #1\par}
\newcolumntype{L}[1]{>{\raggedright\arraybackslash}p{#1}}
\newcolumntype{Y}{>{\raggedright\arraybackslash}X}
\title{\textbf{Renewal Yang--Mills Theory}\\[0.55em]\Large Finite Construction, Renormalization,\\and Conditional Mass-Gap Closure}
\author{A.~P\'elissier}
\date{Specialized monograph, Version 4\\13 September 2026}
\hypersetup{
  pdftitle={Renewal Yang--Mills Theory: Finite Construction, Renormalization, and Conditional Mass-Gap Closure},
  pdfauthor={A. Pélissier},
  pdfsubject={Finite Yang--Mills construction, gauge-compatible Gaussian shells, and conditional continuum mass-gap closure}
}
\begin{document}
\frontmatter
\begin{titlepage}
\centering
\vspace*{25mm}
{\small THE RENEWAL PROGRAMME\par}
\vspace{15mm}
{\Huge\bfseries Renewal\\[3mm]Yang--Mills Theory\par}
\vspace{12mm}
{\Large Finite Construction, Renormalization,\\[2mm]and Conditional Mass-Gap Closure\par}
\vspace{20mm}
{\large A.~P\'elissier\par}
\vfill
{\normalsize Specialized monograph\quad\textbar\quad Version 4\par}
\vspace{3mm}
{\normalsize 13 September 2026\par}
\end{titlepage}
\chapter*{Abstract}
\addcontentsline{toc}{chapter}{Abstract}
The Yang--Mills sector of the Renewal Grand Tensor admits a finite,
source-resolved formulation in which gauge reduction, conditional
elimination, response, and reflected transfer arise from a common
occurrence law. This monograph develops the conditional route from that
construction to a nontrivial continuum theory with a physical mass gap.
It combines exact Schur and conditional-variance identities, normalized
connected-source calculus, measured field redefinitions, degree-faithful
analytic forests, calibrated renormalization, and complete local
spectral criteria. An explicit nested block-average Gaussian tower has
positive conditional shells, quantitative transverse nondegeneracy, and
a fixed-point difference decaying at rate $4^{-j}$. Its Wilson response
retains the distinction between fluctuation representative, background
prescription, and physical normalization, including the correct
punctured-torus regularity of the representative.

The nonlinear root-path block has a smooth compact completion with
uniformly invertible pivots and an exact inverse-Haar density. Its
moving constrained minimum gives a measured one-loop approximation with
a volume-uniform $O(\beta^{-1/2})$ Hessian remainder on the original
controlled local box. A complete fourth-jet formula and a differentiated
small-momentum integral retain the constraint, pivot acceleration, and
measure terms. An all-momentum flat Bloch identity proves the exact Haar
cancellation and identifies finite-volume constant discrepancies as
Fourier aliases. These local and coefficient-level results are not
identifications of the full compact complement or the physical anomaly.

For the actual compact $\SU(2)$ parent, reported stationary-orbit
certificates imply both an explicit local tube Poincar\'e inequality
and metastability of the full hidden diffusion. The higher well's
localized contribution to the retained effective action nevertheless
admits a normalized derivative comparison without hidden mixing. A
global same-law resolvent inequality transfers visible negative-curvature
certificates to original-clock interface heat-bath repair floors.
Non-wrapping shells, local coarse-gauge transport, noncentral matrix
tests, and genuinely active pivot sources give complementary certified
repair regions. A weighted conditional-patch criterion makes the
additional all-exterior estimate for unrestricted interface control
quantitative.

The physical spectral endpoint remains separate from these auxiliary
inequalities. Exact two-level transfer keeps both retained temporal
coercivity and conditional detail control visible; vacuum-corrected
Feshbach and complete local Osterwalder--Schrader Gram criteria provide
alternative sufficient mechanisms. The terminal theorem gives a positive
Hamiltonian mass bound from a compatible nontrivial continuum law,
noncollapsed local probes, and contraction on the complete centered
local reflected algebra at one fixed physical time. The remaining
requirements are explicit: consistently normalized all-level response,
complete compact control, cofinal interacting construction, and physical
contraction on the same calibrated branch. No finite fibre floor, local
repair event, or Gaussian coefficient is substituted for that packet.

\chapter*{Scope and logical conventions}
\addcontentsline{toc}{chapter}{Scope and logical conventions}
The exposition is ordered by mathematical dependency. Finite algebraic
identities, estimates on declared analytic domains, and continuum closure
criteria are distinguished throughout. Alternative sufficient spectral
mechanisms are not treated as simultaneous requirements. Repeated uses of
an identity refer to the same law, source space, quotient, and physical
normalization unless an explicit transport is given.

An \emph{exact identity} is an algebraic, probabilistic, or operator
identity on its stated domain. An \emph{estimate} retains its gauge group,
regulator, analytic norm, and uniformity hypotheses. A \emph{reported
finite certificate} is a computational identity or inequality documented
in a cited reference; it is not presented as an independently rerun
calculation. A \emph{conditional criterion} is a proved implication whose
input estimates must still be verified for the proposed four-dimensional
Yang--Mills family.

Three distinctions govern the Gaussian-to-interacting interface. The
gauge quotient does not determine the representative in which individual
vertex legs are evaluated. The fluctuation representative does not choose
the background with respect to which a response is differentiated. A
fixed-background calculation does not identify the physical anomaly
without the corresponding determinant and background comparison. These
distinctions are implemented in \cref{ch:tower} and carried into the
closure hypotheses of \cref{ch:closure}.

Two further distinctions govern the nonlinear continuation. Equality of
local block germs does not identify complete compact pushforwards, and
a one-loop remainder on a restricted box is not a derivative bound on
its compact complement. The conditional construction in
\cref{ch:compact-response} retains both distinctions. In
\cref{ch:compact-repair}, a local Neumann gap, a killed exit floor, an
unrestricted auxiliary gap, and a physical Hamiltonian gap are separate
objects. The counterexample for the hidden diffusion is kept alongside
the local repair theorems rather than superseded by them.

The mathematical dependencies and precise result locations are collected
in \cref{app:concordance}. The finite construction inherited from the
Grand Tensor is recalled only to the extent needed to fix the common
Yang--Mills law, source algebra, and physical clock. No independent
verification of the large computational certificates, or external
literature audit, is asserted.

\tableofcontents
\mainmatter
\chapter{The target and the conditional architecture}
\label{ch:architecture}
\section{What has to be constructed}
The target is a nontrivial four-dimensional Yang--Mills quantum theory in a specified gauge group and vacuum sector, obtained as a compatible continuum limit of the finite renewal loading. Its local gauge-invariant Euclidean correlations must support Osterwalder--Schrader reconstruction, and the reconstructed Hamiltonian must satisfy
\begin{equation}\label{eq:target}
H\Omega=0,\qquad \Ker H=\C\Omega,\qquad
H|_{\Omega^\perp}\succeq mI,\qquad m>0.
\end{equation}
The adjective \emph{physical} is essential: the lower bound is measured in a declared time unit, not in an adjustable update counter or an auxiliary heat-bath time. The algebra must also be declared. A gap on bounded-support local observables is not automatically a gap in every wrapping sector of every finite torus.

Most finite constructions are formulated for a compact gauge group. The earlier colour-traced response banks use the declared $\SU(3)$ normalization, whereas the actual one-loop checks and compact parent, orbit, and shell certificates in $\SU(2)$ have their own Hilbert--Schmidt and action conventions. The smooth pivot construction is formulated for the stated compact matrix-group setting. None of these group-specific quantitative results is silently promoted to every compact simple group; the closure inputs must be supplied for the group under consideration.

\section{The dependency chain}
The principal construction has the following order:
\begin{equation}\label{eq:main-chain}
\begin{aligned}
&\text{common finite history law and gauge-invariant sources}\\
&\quad\longrightarrow\text{exact compact reduction and blocking}\\
&\quad\longrightarrow\text{Gaussian principal tower and controlled interactions}\\
&\quad\longrightarrow\text{calibrated marginal trajectory and local continuum law}\\
&\quad\longrightarrow\text{nontrivial reflected local Hilbert space}\\
&\quad\longrightarrow\text{uniform fixed-physical-time contraction}\\
&\quad\longrightarrow\text{positive Hamiltonian mass gap}.
\end{aligned}
\end{equation}
Each arrow means a theorem with a specific input interface. In particular, the final arrow is a short spectral argument; the preceding contraction estimate is not obtained merely by renaming a finite-dimensional Hessian gap.

The finite steps in this chain have explicit interfaces. The exact
blocking differential retains a unit pullback sector and a conditionally
centered nonlinear remainder. Its measured redundant quotient is
preserved through the pushed conditional current. The third current
jet is tied to four-entry cumulants by the normalized score hierarchy;
the degree-faithful forest and corrected contour estimate provide a
sufficient marked analytic route without replacing its uniformity
hypotheses.

There are alternative ways to supply the spectral input: a complete
effective-action Dobrushin bound, a source-resolved physical
Poincar\'e estimate with a proved same-generator incidence, a two-level
physical form comparison, an orbit-coercive Feshbach argument, or a direct
reflected Gram contraction. These alternatives are not extra conjunctive
requirements. The exact two-level identity identifies the retained
infrared estimate that a detail-fibre gap cannot supply. A successful
route need not solve every historical sufficient condition.

\begin{definition}[Three logically different closure tasks]
\label{def:three-tasks}
\emph{Construction closure} supplies compatible interacting measures, local limits, the required Euclidean axioms, and nontriviality. \emph{Renormalization closure} controls the stable coordinates and identifies the physical marginal trajectory, including its anomaly coefficient and clock. \emph{Spectral closure} gives a uniform contraction or equivalent physical coercivity on the complete centered local sector. None of the three is substituted for either of the others.
\end{definition}

\section{The Gaussian endpoint and its response interface}
The Gaussian construction is an actual nested block-average covariance
tower, rather than a repeatedly reset endpoint chart. Its positive shells,
fixed point, convergence rate, and transverse inverse bounds are explicit.
The full fine-edge Wilson vertices can be coupled to those shells only
after a matching fluctuation representative is chosen. The tree-gauge map
in \cref{sec:tree-gauge} supplies that finite algebraic interface.

The next response requirement is not simply a scalar evaluation with
unspecified vertices. It is a complete full-torus calculation with common
normalization, momentum-derivative control, and a declared background.
The fixed fine-background prescription uses the existing vertex banks;
the harmonic prescription follows the minimizing blocked action and
contains nonzero background aliases. A calculation in the former supplies
the physical Gaussian anomaly only after the averaged comparison in
\cref{crit:background-ADMC}, or an equivalent direct identification, has
been proved. The finite map and corner identities do not establish that
comparison.

The nonlinear compact continuation now has an explicit one-step carrier:
\cref{cor:compact-product} gives the solved compact pivots and their
positive density, and \cref{thm:uniform-local-loop} gives a measured
volume-uniform local one-loop remainder. Their combination does not
identify the full compact complement or establish a compatible
all-level trajectory. The degree-faithful forests, owner decomposition,
and reciprocal-coupling shielding remain the interfaces for that
continuation. Likewise, the actual visible repair floors in
\cref{ch:compact-repair} need either complete coverage and assembly or a
populated conditional-patch criterion, followed by physical form
transport. Even a correctly matched ultraviolet trajectory does not
replace that infrared requirement.

\begin{scope}[Meaning of the word closure]
Throughout this monograph, a \emph{conditional closure theorem} means that the conclusion follows from the displayed packet. It does not mean that the complete packet has already been populated. The results established here do not yet imply \cref{eq:target} without the remaining continuum and physical-contraction hypotheses.
\end{scope}
\sourcebasis{Local closure and physical source criteria are established in \src{F1}{V1--V4}; the nested Gaussian construction and representative-dependent response in \src{F11}{V31--V41}. The common finite source and physical-calibration framework is developed in \cite{GT}.}

\chapter{A short reconstruction of the finite Yang--Mills loading}
\label{ch:finite}
\section{From histories to the readable finite state}
The Grand Tensor begins with a typed event grammar and actual occurrence records. Two past histories are predictively equivalent when every admitted record-conditioned future has the same law. The quotient is the coarsest persistent predictive carrier. A finite ledger is a compiled realization of this quotient, not an independently postulated hidden state space.

Prediction alone does not supply spatial routing, orientation, or physical time. The readable enrichment retains the same-cut realization, protected orientation, root and scheduler marks, calibration, determining-field data, and action records needed by the downstream construction. Forgetting those added readable marks recovers the original Store cylinders. Across cutoffs, old--detail polarization tables carry the refinement information. In the Yang--Mills specialization, this common interface fixes the finite gauge variables, their action, the admitted gauge-invariant writers, and the reflected time transfer. It is not necessary to repeat the Standard Model or spacetime classifications to formulate the Yang--Mills closure problem. \cite{GT}

\section{One law, its reductions, and its sources}
At a finite regulator let $X$ denote the field configuration, $\mu$ its probability law, and $\mathfrak A$ the admitted gauge-invariant source algebra. The essential discipline is that a source, a conditional expectation, a Schur complement, and a response derivative refer to this same law or to an explicitly specified pushforward of it.

A subgroup restriction, conditioning on a retained field, integrating out a detail field, and pushing a measure to a quotient are different operations. Gauge invariance alone does not make their resulting laws identical. In particular, a reduced Wilson fibre at fixed coarse background is not automatically the physical Perron ground law of a transfer operator.

\begin{proposition}[Exact Gaussian elimination]\label{prop:schur}
Let a quadratic form on retained and detail variables have matrix
\[
L=\begin{pmatrix}A&B\\B^*&C\end{pmatrix},\qquad C\succ0.
\]
The minimizing detail is $y=-C^{-1}B^*x$, and
\begin{equation}\label{eq:schur}
\ip{(x,y)}{L(x,y)}
=\ip{x}{(A-BC^{-1}B^*)x}
 +\ip{y+C^{-1}B^*x}{C(y+C^{-1}B^*x)}.
\end{equation}
Gaussian integration retains both the Schur form $A-BC^{-1}B^*$ and the determinant contribution $\frac12\log\det C$ to the effective action, up to background-independent normalization.
\end{proposition}
\begin{proof}
Expand the last quadratic term in \cref{eq:schur}. Its cross terms recover those of $L$, and its retained square cancels the subtracted Schur term. Integration of the centered detail Gaussian gives $(\det C)^{-1/2}$ times the fixed Gaussian normalization.
\end{proof}

This identity fixes two independent pieces of data: the retained response and the normalization. Keeping the first while discarding a background-dependent determinant changes the effective action. Repeating an exact elimination also requires the retained law from the preceding step, rather than an unrelated reference with the same quadratic Taylor jet.

\section{Half-slab amplitudes and the physical Doob law}
Let $F(y,x)$ be the amplitude for a half-slab, and define
\[
(Af)(y)=\int F(y,x)f(x)\,\dd\mu_0(x),\qquad T=A^*A.
\]
For $F\in L^2$, $A$ is Hilbert--Schmidt and
\begin{equation}\label{eq:HS}
\Tr T=\norm{A}_{\rm HS}^2=\int|F|^2,\qquad
\Tr(\wedge^rT)=\norm{\wedge^rA}_{\rm HS}^2\ge0.
\end{equation}
Thus the exterior-power trace coefficients are genuine positive finite data. They do not alone control the ratio of the first two singular values uniformly in the regulator.

Assume now that $F$ is strictly positive and continuous on compact spaces with full-support reference measures. Write its largest singular value as $s_1$, with positive normalized singular functions $u_0,v_0$. The joint probability law
\begin{equation}\label{eq:half-law}
\dd\Lambda(x,y)=s_1^{-1}u_0(x)F(y,x)v_0(y)
\,\dd\mu_0(x)\dd\mu_{1/2}(y)
\end{equation}
has marginals $u_0^2\mu_0$ and $v_0^2\mu_{1/2}$. Conditional expectation between the two marginals is unitarily equivalent to $s_1^{-1}A$.

\begin{theorem}[Half-slab correlation and full transfer]\label{thm:half}
Under these finite-regulator hypotheses, the centered half-slab maximal correlation is
\[
\rho_{1/2}=\frac{s_2}{s_1}.
\]
The centered normalized full transfer has norm $q=\rho_{1/2}^2$. For a full slab of physical duration $t$, its spectral mass is
\begin{equation}\label{eq:half-mass}
m_t=-\frac{2}{t}\log\rho_{1/2}.
\end{equation}
\end{theorem}
\begin{proof}
Multiplication by the positive Perron functions identifies the two marginal $L^2$ spaces with the original reference spaces. In these coordinates the conditional channel is $A/s_1$, the constant functions correspond to the top singular vectors, and its centered singular norm is $s_2/s_1$. Composing the half channels gives $A^*A/s_1^2$. Functional calculus gives \cref{eq:half-mass}.
\end{proof}

If a common measurable factor survives on both boundaries, the relevant constant projection is replaced by the projection $P_{\rm com}$ onto that common factor. With boundary projections $P_0,P_1$,
\[
\rho_{\rm rel}=\norm{(P_1-P_{\rm com})(P_0-P_{\rm com})}.
\]
The two-boundary averaging chain has relative gap $(1-\rho_{\rm rel})/2$. This is a gap relative to retained common labels. It does not justify deleting protected sectors and calling the remainder the full theory.

\subsection{The exact conditional-variance form}
The half-slab description has a useful variational formulation which
retains the physical Perron factors. Let $\mathsf Kf(y)=\E_\Lambda[f(X)\mid
Y=y]$ and let $P_{\rm slab}=\mathsf K^*\mathsf K$ on the initial marginal
$\pi=u_0^2\mu_0$. Write $\pi_{1/2}=v_0^2\mu_{1/2}$ for the
other marginal. The ground-transform identification above gives the normalized
full transfer, not an auxiliary conditional sampler.

\begin{proposition}[Half-slab variance identity]
\label{prop:half-variance}
For every $f\in L^2(\pi)$,
\begin{equation}\label{eq:half-variance}
\begin{aligned}
\ip{f}{(I-P_{\rm slab})f}_{L^2(\pi)}
&=\E_\Lambda\!\left[\Var_\Lambda(f(X)\mid Y)\right]\\
&=\inf_{h\in L^2(\pi_{1/2})}
  \E_\Lambda|f(X)-h(Y)|^2.
\end{aligned}
\end{equation}
In particular, a bound by $\gamma\Var_\pi(f)$ from below is equivalent
to $\norm{P_{\rm slab}|_{\one^\perp}}\le1-\gamma$. For a full physical
slab of duration $t$, $0<\gamma<1$ therefore gives
\begin{equation}\label{eq:half-variance-mass}
 m_t\ge-\frac1t\log(1-\gamma).
\end{equation}
\end{proposition}
\begin{proof}
Conditional expectation is the orthogonal projection of $f(X)$ onto
functions of $Y$. Its squared projection norm is $\norm{\mathsf Kf}^2$;
subtracting this from $\norm f^2$ proves both expressions in
\cref{eq:half-variance}. Positivity and self-adjointness of
$P_{\rm slab}$ identify the best centered form bound with one minus
its centered operator norm. Functional calculus proves
\cref{eq:half-variance-mass}. \src{F6}{V58}
\end{proof}

\begin{lemma}[Same-carrier density comparison]
\label{lem:half-density-comparison}
Let $\Lambda_0$ be a reference probability law on the same boundary
spaces and suppose
\[
0<m_0\le\frac{\dd\Lambda}{\dd\Lambda_0}\le M_0<\infty.
\]
If the reference conditional-variance form is at least
$\gamma_0\Var_{\pi_0}(f)$, then the physical form in
\cref{eq:half-variance} is at least
\begin{equation}\label{eq:half-density-comparison}
\frac{m_0}{M_0}\gamma_0\Var_\pi(f).
\end{equation}
\end{lemma}
\begin{proof}
The lower density bound applies before the infimum over $h$ in
\cref{eq:half-variance}. The upper marginal density bound gives
$\Var_\pi(f)\le M_0\Var_{\pi_0}(f)$ by minimizing over constants.
Combining the two estimates proves the assertion. \src{F6}{V58}
\end{proof}

The comparison must include the singular ground profiles in
\cref{eq:half-law}. For a Gibbs density ratio, $M_0/m_0$ is controlled
by the oscillation of the complete density difference, which can grow
with volume even when its local summands are small. The lemma gives an
explicit law-transport cost; it does not make that cost uniform. Local
conditional comparison or a signed form estimate is needed whenever a
global ratio loses the required physical margin.

\section{Physical update versus auxiliary dynamics}
For a normalized positive transfer $T\Omega=\Omega$, the Doob operator
\[
Pf=\Omega^{-1}T(\Omega f),\qquad \dd\pi=|\Omega|^2\dd\nu,
\]
is a self-adjoint Markov contraction. In its stationary two-sided path realization, let $U$ be the shift and $J$ the time-zero embedding. Then
\begin{equation}\label{eq:koopman}
J^*U^nJ=P^n,\qquad
J^*(U-I)^*(U-I)J=2(I-P).
\end{equation}
These are identities for the actual transfer. A spatial update chain, a conditional Gibbs sampler, or a comparison Witten operator can be useful upstream, but is not made into $-\log T/t$ by notation.

\begin{remark}[Why the finite recap is enough]
The rest of the monograph uses only the admitted gauge-invariant source algebra, the actual compact action and its exact pushforwards, the same-law reflected transfer, and the calibrated time. The upstream predictive and readable construction explains why these objects belong to one interface; it does not remove the analytic obligations of continuum Yang--Mills theory.
\end{remark}
\sourcebasis{Grand Tensor Yang--Mills loading, half-slab and Doob results \cite{GT}; Perron-law and same-generator analysis \src{F1}{V7--V11, V21}; exact half-slab variance and same-carrier density comparison \src{F6}{V58}.}

\chapter{Source-resolved coercivity and the law-matching problem}
\label{ch:source}
\section{The physical ground law cannot be guessed}
Let a symmetric transfer have the form
\[
T=e^{-G/2}\,\mathcal C\,e^{-G/2},\qquad Th=\lambda h,
\]
where $h>0$. Its one-slice law is proportional to $h^2\dd\mathfrak m$. A candidate density $Z^{-1}e^{-F}\dd\mathfrak m$ agrees with that law precisely when
\begin{equation}\label{eq:law-incidence}
\mathcal C(e^{-(G+F)/2})=\lambda e^{(G-F)/2}.
\end{equation}

\begin{proposition}[No naive same-exponent identification]\label{prop:law-no-go}
If $\mathcal C1=1$ and $\mathcal C$ is injective, then the candidate with $F=G$ is the Perron law if and only if $G$ is constant almost everywhere.
\end{proposition}
\begin{proof}
With $F=G$, \cref{eq:law-incidence} gives $\mathcal C(e^{-G})=\lambda1$. Hence $\mathcal C(e^{-G}-\lambda)=0$. Injectivity forces $e^{-G}=\lambda$. The converse is immediate.
\end{proof}

The temporal Wilson convolution in the stated finite setting has strictly positive character multipliers, so this obstruction applies to its full configuration-space law. A conditional identity at fixed coarse field is a different statement and must be checked by disintegration. The valid options are to prove that conditional incidence, transport through an explicit density defect with a controlled loss, or rebuild the source calculus directly on the Perron law. \src{F1}{V21}

\section{Scores, curvature, and the complete signed source}
For a normalized conditional density proportional to $e^{-S(a,x)}$, parameter differentiation gives
\begin{equation}\label{eq:score}
\partial_a\E_a f=\E_a(\partial_af)-\Cov_a(f,\partial_aS),
\end{equation}
and the effective action $\mathcal F(a)=-\log\int e^{-S(a,x)}\dd x$ satisfies
\begin{equation}\label{eq:free-hess}
\partial_b\partial_a\mathcal F
=\E_a(\partial_b\partial_aS)-\Cov_a(\partial_aS,\partial_bS).
\end{equation}
The covariance is a signed part of the complete curvature source. Bounding a direct term and later charging the same covariance again changes the operator. The Perron and sequential-marginal analogues retain the terminal Perron weight and its derivatives; the unweighted conditional fibre is not a replacement for them.

For the moving conditional projection $D_i=I-\E_i$, with normalized conditional density $\rho_i$, the source identity is
\begin{equation}\label{eq:moving-score}
X_aD_if=D_iX_af-\Cov_i(f,X_a\log\rho_i).
\end{equation}
Different sign conventions for the density require the corresponding score sign. The identity, not a memorized sign, determines the response.

\subsection{Higher normalized responses}
The first-order score identity extends exactly to the higher responses
needed by blocking. Work on a fixed finite conditional chart with
normalized law
\[
\dd\nu_a(x)=Z(a)^{-1}e^{-S(a,x)}\dd m(x).
\]
Assume the indicated derivatives may pass under the integral. A moving
reference density is included in $S$ after transport to this chart; it
cannot be omitted from the score. For a finite set $I$ of labelled,
commuting parameter directions, write $\partial_I=\prod_{i\in I}\partial_i$,
$F_J=\partial_JF$, and $S_B=\partial_BS$. Denote joint cumulants under
$\nu_a$ by $\kappa_a$, with $\kappa_a(G)=\E_aG$, and let $\Pi(A)$ be
the set of partitions of $A$.

\begin{theorem}[Complete connected score hierarchy]
\label{thm:connected-score-hierarchy}
With the empty-partition convention,
\begin{equation}\label{eq:connected-score-hierarchy}
\partial_I\E_aF
=\sum_{J\subseteq I}\ \sum_{\pi\in\Pi(I\setminus J)}
 (-1)^{|\pi|}\kappa_a\bigl(F_J,(S_B)_{B\in\pi}\bigr).
\end{equation}
In particular, a third derivative of a normalized current can contain a
four-entry connected cumulant. All denominator derivatives are already
included in this identity.
\end{theorem}
\begin{proof}
Differentiate the logarithm of the joint moment-generating function to
obtain
\[
\partial_i\kappa_a(G_1,\ldots,G_k)
=\sum_{\ell=1}^k\kappa_a(G_1,\ldots,\partial_iG_\ell,\ldots,G_k)
 -\kappa_a(G_1,\ldots,G_k,S_{\{i\}}).
\]
For $k=1$ this is \cref{eq:score}. Induct on $|I|$: a new direction
joins $J$, joins an existing score block, or creates a singleton score
block. These possibilities enumerate each term in the new partition
sum exactly once, with the sign displayed in
\cref{eq:connected-score-hierarchy}. \src{F8}{V60}
\end{proof}

For example, writing subscripts for parameter derivatives, the second
response is
\begin{equation}\label{eq:score-second-explicit}
\begin{aligned}
\partial_b\partial_a\E F
={}&\E F_{ab}-\kappa(F_a,S_b)-\kappa(F_b,S_a)\\
 &-\kappa(F,S_{ab})+\kappa(F,S_a,S_b).
\end{aligned}
\end{equation}
At third order the term with three singleton score blocks is
$-\kappa(F,S_a,S_b,S_c)$. Pair-covariance control cannot pay for this
term. The four-mark criterion in \cref{thm:four-mark-clustering} supplies
a sufficient connected estimate once its complete source-dependent
activity hypothesis is verified.

\begin{scope}[An exact hierarchy is not an all-order norm estimate]
Formula \cref{eq:connected-score-hierarchy} holds at every fixed order
under its differentiation hypotheses. Bounding its partitions separately
does not prove a uniform exponential-grade analytic estimate. The
factorial restrictions in \cref{ch:trees} and the radius budgets in
\cref{ch:forests} continue to apply.
\end{scope}

\section{A physical response-matrix compiler}
Suppose the physical ground-state representation has Dirichlet form
\begin{equation}\label{eq:groundform}
\ip{\Omega f}{(H-E_0)\Omega f}
=\kappa_H\int|\nabla f|^2\dd\pi.
\end{equation}
For the sequential filtration $\mathcal F_i$, let $M_i f=\E[f\mid\mathcal F_i]$ and $D_i=M_i-M_{i-1}$. The exact variance ledger is
\[
\Var_\pi(f)=\sum_i\norm{D_if}_{L^2(\pi)}^2.
\]

\begin{criterion}[Sequential source response]\label{thm:response-matrix}
Assume every sequential single-coordinate conditional has Poincar\'e floor at least $\rho_{\rm site}>0$, and a nonnegative matrix $\mathsf A$ satisfies
\[
\norm{\nabla_iM_if}_{L^2(\pi)}
\le\sum_bA_{ib}\norm{\nabla_bf}_{L^2(\pi)}
\]
for every admitted smooth gauge-invariant $f$. Then
\begin{equation}\label{eq:response-gap}
\Var_\pi(f)\le\frac{\norm{\mathsf A}_{2\to2}^2}{\rho_{\rm site}}
\int|\nabla f|^2\dd\pi,\qquad
\gap(H)\ge\frac{\kappa_H\rho_{\rm site}}{\norm{\mathsf A}_{2\to2}^2}.
\end{equation}
\end{criterion}
\begin{proof}
Conditional Poincar\'e gives $\norm{D_if}^2\le\rho_{\rm site}^{-1}\norm{\nabla_iM_if}^2$. Apply the matrix norm to the vector of gradient norms, sum the martingale differences, and use \cref{eq:groundform}.
\end{proof}

A simpler sufficient packet is $\norm{D_if}^2\le\sum_bw_{ib}\int|\nabla_bf|^2\dd\pi$ with bounded column congestion $C_{\rm col}=\sup_b\sum_iw_{ib}$. It gives $\gap(H)\ge\kappa_H/C_{\rm col}$. Conditional site gaps alone are insufficient: their responses to the previously exposed field must also be controlled. The source-resolved Green-column construction is designed to populate these response matrices on the actual signed source range, rather than by the norm of a full inverse. \src{F1}{V7, V14--V16}

\section{Comparison-minus-defect and physical landing}
On a declared one-form source carrier, write
\[
L=\mathcal A-\mathcal V\mathcal V^*,\qquad \mathcal A\succ0,
\]
and define
\begin{equation}\label{eq:woodbury-objects}
\mathcal K=\mathcal V^*\mathcal A^{-1}\mathcal V,
\quad \mathcal B=\mathcal A^{-1}\mathcal V,
\quad \mathcal M=\mathcal B^*\mathcal B.
\end{equation}
When the inverse is defined and $\norm{\mathcal K}<1$, Woodbury gives
\[
L^{-1}=\mathcal A^{-1}
+\mathcal A^{-1}\mathcal V(I-\mathcal K)^{-1}\mathcal V^*\mathcal A^{-1}.
\]
The exact identities are more informative than this sufficient Neumann condition.

\begin{theorem}[Channel-to-physical landing]\label{thm:landing}
For $\omega_z=\mathcal Bz$,
\begin{equation}\label{eq:landing}
\begin{aligned}
\norm{\omega_z}^2&=\ip{z}{\mathcal Mz},\\
\ip{\omega_z}{\mathcal A\omega_z}&=\ip{z}{\mathcal Kz},\\
\ip{\omega_z}{L\omega_z}&=\ip{z}{\mathcal K(I-\mathcal K)z}.
\end{aligned}
\end{equation}
On the closure of $\Ran\mathcal B$ in the comparison-form norm,
\begin{equation}\label{eq:relative-edge}
1-\norm{\mathcal K}
=\inf_{\omega\ne0}
\frac{\ip{\omega}{L\omega}}{\ip{\omega}{\mathcal A\omega}}.
\end{equation}
\end{theorem}
\begin{proof}
The first two equalities are direct substitutions. Also $\mathcal V^*\omega_z=\mathcal Kz$, giving the third. The supremum of $\ip{z}{\mathcal K^2z}/\ip{z}{\mathcal Kz}$ is $\norm{\mathcal K}$ by spectral calculus, which proves \cref{eq:relative-edge}.
\end{proof}

A near-unit channel sequence therefore need not be a physical infrared sequence. Its landing norm can degenerate while its comparison energy escapes to high frequencies. A uniform landing frame or a bounded comparison-energy window is needed before a channel survivor gives a normalized physical soft vector. Conversely, a uniform bound on the complete defect synthesis can force the required landing window. The source-carrier Combes--Thomas estimate then uses a genuine spectral edge and a controlled conjugation modulus to obtain spatially summable response columns. It does not manufacture the edge. \src{F1}{V16--V20}

\section{Bad-set capacity follows source incidence}
Let $B_i$ be a predictable bad event for martingale increment $D_i$. Define $Z_i(f)=\norm{\one_{B_i}D_if}$ and $X_b(f)=\norm{\nabla_bf}$. Suppose an assembled source vector $R(f)$ obeys
\[
Z(f)\le\mathsf I R(f)+\mathsf E X(f),\qquad R(f)\le\mathsf QX(f)
\]
entrywise, with nonnegative matrices. Then
\begin{equation}\label{eq:bad-capacity}
\sum_i Z_i(f)^2\le
\norm{\mathsf I\mathsf Q+\mathsf E}_{2\to2}^2\int|\nabla f|^2\dd\pi.
\end{equation}
With good-site floor $\rho_g$ and good response matrix $\mathsf A$, the hybrid physical gap is bounded below by
\begin{equation}\label{eq:hybrid-gap}
\gap(H)\ge
\frac{\kappa_H}{\rho_g^{-1}\norm{\mathsf A}_{2\to2}^2+
\norm{\mathsf I\mathsf Q+\mathsf E}_{2\to2}^2}.
\end{equation}
The first matrix is source incidence; the second is capacity routing. A small probability of $B_i$ is neither one. This distinction is the reason that a Fisher-information average or a Chebyshev estimate cannot replace a restricted operator estimate. \src{F1}{V19--V20}

\section{What these reductions actually close}
The source calculus closes the algebraic relation between physical response, signed curvature, comparison return, and the martingale gap. It also identifies exact failure witnesses. It does not yet supply a regulator-uniform response matrix or bad-capacity stock for the weak-coupling continuum trajectory.

The finite temporal Wilson-to-Hamiltonian construction is another retained result: calibrated character multipliers converge to their Casimir symbols on fixed representation banks, heat-kernel domination controls high representations, and a diagonal growing schedule gives the fixed-spatial-volume differential-Hamiltonian limit. Strong semigroup convergence at fixed volume is not a uniform thermodynamic mass theorem. The order of limits and the physical coefficient in \cref{eq:groundform} remain part of the application. \src{F1}{V10--V11}
\sourcebasis{The actual-law, sequential-response, source-carrier, and capacity results are developed in \src{F1}{V6--V22}; the moving conditional calculus is continued in \src{F6}{terminal developments} and \src{F7}{opening conditional analysis}.}

\chapter{Harmonic normalization and connected-source estimates}
\label{ch:trees}
\section{Peter--Weyl bookkeeping before inequalities}
The earliest response programme studies products of representation-valued sources and the cancellation of disconnected partitions. A recurring obstruction is incorrect normalization: dimensions, raw Fourier coefficients, representation probabilities, and marked-coordinate probabilities cannot be interchanged.

Let
\[
\bigotimes_{i=1}^m H_{R_i}
=\bigoplus_U W_U(H_U\otimes M_U),\qquad
X_U=W_U^*\left(\bigotimes_i A_i\right)W_U,
\]
where $W_U$ denotes the isometric block embedding. With the Fourier and kernel conventions of the source, the exact coefficient is
\begin{equation}\label{eq:PW}
\widehat h(U)=\frac{d_{\rm in}}{d_U}
\Tr_{M_U}\bigl[X_U(I\otimes K_U)\bigr],\qquad
 d_{\rm in}=\prod_i d_{R_i}.
\end{equation}
The trace-class pinching inequality is
\begin{equation}\label{eq:pinching}
\sum_U\norm{X_U}_1\le\prod_i\norm{A_i}_1.
\end{equation}
The output factor $d_U$ cancels only against its prescribed Fourier weight. A dual-singlet test retains raw mass $1/d_R$; replacing it by a probability inferred from representation dimensions changes the estimate. These elementary tests are retained because later tree estimates depend on their exact grade. \src{F1}{harmonic source programme}; \src{F2}{opening normalization audit}

For nonnegative coordinate stocks $a_{i,e}$ define
\begin{equation}\label{eq:marks}
\Xi_i=\sum_e a_{i,e},\quad
H_{ij}=\sum_e a_{i,e}a_{j,e},\quad
\theta_{ij}=\frac{H_{ij}}{\Xi_i\Xi_j}\le1.
\end{equation}
Zero-stock sources are removed before normalization. The $\theta_{ij}$ are genuine probability-mark overlaps. A tree numerator in the unnormalized $H_{ij}$ is not yet the normalized marked-tree statement in the $\theta_{ij}$.

\section{First connectivity and the quartic closure}
A connected coordinate word has a unique first coordinate at which the current row partition becomes connected. Before that coordinate, the word factors over a proper partition. At that coordinate, the complete joining cumulant acts. After it, the suffix is a complete moment. The final resolvent is paid once after this assembly, rather than separately in every predecessor block.

The quartic problem is resolved only after the normalization and suffix corrections. Its fourteen proper predecessor partitions consist of four of type $3+1$, three of type $2+2$, six of type $2+1+1$, and one all-singleton partition. Certain literal local weighted estimates fail because a suffix can rescue a channel that is small in its prefix. The successful bounds sum the rescued square globally over coordinates before taking the final norm.

\begin{theorem}[Completed fixed-arity quartic marked-tree estimate]\label{thm:quartic}
For the declared normalized quartic source problem, with its single final resolvent and probability-mark convention, the complete connected coefficient satisfies
\begin{equation}\label{eq:quartic}
\kappa_{\otimes}^{\rm abs}(\mathcal C^{\rm pay}_4)
\le C_{G,s}\sum_{T\in\mathcal T_4}\prod_{\{i,j\}\in T}\theta_{ij}.
\end{equation}
Here $\kappa_{\otimes}^{\rm abs}$ is the source-defined absolute product-state coefficient capacity, and $C_{G,s}$ has the group and regulator dependence specified in that construction.
\end{theorem}
\begin{proof}[Proof structure]
Split at first connectivity. Use the corrected normalized $3+1$ atlas and the globally coordinate-summable rescued squares for the $3+1$ and $2+1+1$ cases; treat $2+2$ and the singleton case by their completed collision bounds. Summing the fourteen proper partitions and retaining one final payer yields \cref{eq:quartic}. The conclusion is the normalized coefficient-capacity estimate, not an unnormalized operator-tree numerator; the globally summed suffix is essential. \src{F3}{V89--V96}; \src{F4}{inherited quartic ledger}
\end{proof}

At arity five there are fifty-one proper predecessor partitions and $5^3=125$ labeled trees. The local cumulant and projected-tail development closes substantial fixed-arity estimates, including the Gaussian tangent degree filtration. It also produces counterexamples: one Casimir payment per input is not sufficient in the stated coherent $\SU(3)$ packet; a whole $\SU(2)$ heat slot retains singlet capacity of order $(1+C_2)^{-1/2}$; and endpoint leakage cannot be removed by separately bounding merge differences. These results constrain the all-order extension rather than establishing it. \src{F3}{V97--V100}; \src{F4}{quintic and projected-tail developments}

\section{The factorial-grade correction}
The all-order question is not settled by proving each fixed arity with a large constant.

\begin{definition}[Strong and factorial-graded tree bounds]\label{def:factorial}
A strong marked-tree bound has, for each fixed labeled tree,
\[
\norm{\mathcal K_{m,T}}\le K_0^{m-1}\mathcal W_T.
\]
A factorial-graded bound instead has
\[
\norm{\mathcal K_{m,T}}\le m!K^m\mathcal W_T.
\]
The distinction is uniform in $m$, not a change in the numerical constant at one arity.
\end{definition}

\begin{proposition}[The extra factorial cannot be hidden]\label{prop:factorial}
The all-order $\SU(2)$ estimates of \src{F5}{V1--V11} have the factorial-graded form. They do not imply the strong marked-tree estimate or its uniform field-cumulant consequence.
\end{proposition}
\begin{proof}
The exact predecessor-partition sum is
\begin{equation}\label{eq:partition-factorial}
\sum_{\rho\in\Pi([m])}|\rho|!\prod_{B\in\rho}|B|!
=m!2^{m-1}.
\end{equation}
Rooted labeled-tree enumeration already uses the allowed factorial,
\[
m^{m-1}\le e^{m-1}m!.
\]
A factorial per fixed tree therefore produces order $(m!)^2K_1^m$ after spatial compilation. No fixed exponential base absorbs the extra $m!$. The coefficient $1/m!$ in an exponential generating series cancels only one factorial; the remaining ordinary series is still of factorial grade. This is the factorial-grade obstruction; it leaves the finite inequalities valid but rules out the stronger uniform analytic inference. \src{F5}{V12}
\end{proof}

The correction leaves the displayed all-order inequalities valid. It removes only their overstrong analytic interpretation. It is therefore inappropriate either to discard the finite estimates or to describe the strong all-order source gate as closed.

\section{Exact Haar assembly recovers one exponential-grade component}
The $\SU(2)$ exponential-chart Jacobian, in the source convention, is
\begin{equation}\label{eq:Haar-series}
J_\epsilon(y)=\left(\frac{\sin(\sqrt\epsilon|y|)}{\sqrt\epsilon|y|}\right)^2
=\sum_{n\ge0}\frac{(-1)^n2^{2n+1}}{(2n+2)!}\epsilon^n|y|^{2n}.
\end{equation}
Its logarithmic vertices have large isolated high derivatives; bounding them before recombination loses the factorial denominators of the full Jacobian.

\begin{theorem}[Integrated exact-Haar derivative bound]\label{thm:Haar}
There are $C<\infty$ and $\epsilon_0>0$ such that for every centered Gaussian $Y\in\R^3$ with covariance bounded by $I$, every $0<\epsilon\le\epsilon_0$, and every $v\ge1$,
\begin{equation}\label{eq:Haar-bound}
\E\norm{\nabla^vJ_\epsilon(Y)}_{\rm row}
\le C^v\epsilon^{v/2}.
\end{equation}
The zeroth-order expectation is bounded by a fixed constant.
\end{theorem}
\begin{proof}[Proof outline]
Differentiate \cref{eq:Haar-series} term by term. The radial derivative row is at most
\[
C^v\frac{(2n)!}{(2n-v)!}(1+|y|)^{2n-v}.
\]
Gaussian moments give $\E(1+|Y|)^q\le C_0^q q^{q/2}$. The denominator $(2n+2)!$ cancels the dangerous derivative factorial, while $q^{q/2}/q!\le e^q$. The sum begins at $n\ge\lceil v/2\rceil$ and is bounded by a convergent geometric series in $C_1\epsilon$, giving \cref{eq:Haar-bound}. \src{F5}{V12}
\end{proof}

This closes the exact-Haar component, not the whole Wilson-density transfer. The exact action, normalization, assembled interpolation, and tree assignment must still have the correct fixed-tree grade.

\section{Forest projection versus physical source occurrences}
The later compact heat and L\'evy formulation replaces several coefficientwise arguments by an exact probabilistic interpolation. Optional support labels are summed before norms using convex contractions of the form $(1-\lambda)I+\lambda D^R(U)$. At the actual hierarchical interpolation matrices, the resulting weights are subprobabilities rather than Bell-number counts. Weighted Cayley disintegration, random-root orientation, and marginalization of unused completion branches give the strong vertex-accreting forest-projection estimate with total base $K^m$.

A physical source occurrence is not an unused completion branch. It owns a disjoint nonempty subforest and can carry a forced component excess. Marginalizing that ownership would remove a real derivative payment. Thus the complete forest-projection result and the source-occurrence closure are different achievements. \src{F6}{V1--V7}

One useful conditional reduction is as follows. On a carrier graph of degree $\Delta$, suppose source activities $z_A$ satisfy, for each component $\alpha$ and connected carrier $\Gamma$,
\begin{equation}\label{eq:carrier-profile}
\sum_{\substack{A\ni\alpha\\\Gamma(A)=\Gamma}}|A|z_A
\le C_0|\Gamma|^r a^{|\Gamma|}.
\end{equation}
Then the all-rank incidence row is bounded by
\begin{equation}\label{eq:carrier-row}
\kappa_{\rm hyp}\le C_0a\sum_{n\ge1}n^r(\Delta^2a)^{n-1},
\qquad \Delta^2a<1.
\end{equation}
The graph enumeration and the implication are closed. The physical carrier-activity profile \cref{eq:carrier-profile} is an input, not a consequence of the abstract forest projection. \src{F6}{V16}
\sourcebasis{The harmonic, collision, arity, and source-occurrence developments are distributed across \cite{F1,F2,F3,F4,F5,F6}. The factorial-grade restriction of \src{F5}{V12} governs every use of the all-order bound.}

\chapter{Conditional geometry, collars, and bad-field localization}
\label{ch:geometry}
\section{Local control must include the conditioned exterior}
An isolated tree gauge can simplify a finite block, but changing that gauge may also change its exterior data. Such a transformation relates conditional fibres; it is not necessarily a symmetry within one fixed conditional fibre. Weak-coupling exterior staples can create competing central wells, so an unconditional good-field probability cannot justify a uniform conditional convexity assertion.

The source analysis keeps the fixed exterior explicit. For six $\SU(3)$ staples $S_j$, set
\[
F=\min_V\sum_{j=1}^{6}\bigl(3-\Re\Tr(VS_j)\bigr).
\]
After the declared alignment, the staple sum differs from $6I$ by at most $\sqrt{12F}$ in Frobenius norm. A sufficiently small value of $F$ gives the single-well local quadratic control used by the good-field analysis. The direct star defect
\[
D_e=\sum_{j=1}^{6}(3-\Re\Tr S_j)
=\frac12\sum_{j=1}^{6}\norm{S_j-I}_{\rm F}^2
\]
satisfies $F_e\le D_e$. These are local, conditioned statements. \src{F6}{terminal staple analysis}; \src{F7}{opening geometric audit}

\section{The finite collar reconstruction}
The one-layer collar used in the source consists of $72$ plaquettes, $123$ edges, and $64$ vertices. Its scalar curvature map has rank $60$ and kernel equal to the $63$-dimensional gauge image. Fixing the full tree leaves $60$ chord variables. The reported integer certificate contains a $60\times60$ minor with determinant $-1$ and $102$ nonzero entries, giving the explicit bound
\begin{equation}\label{eq:collar-sigma}
\sigma_{\min}^2\ge102^{-59}.
\end{equation}
A $60$-generator, $72$-relator elimination certificate identifies the relevant complex as simply connected. Consequently the zero-curvature completion is unique after all prescribed tree frames are fixed, and a small-defect estimate extends to
\begin{equation}\label{eq:collar-distance}
\dist(U,\mathcal Z_{\rm flat})^2\le C_{\rm glob}D(U)
\end{equation}
on the declared finite collar with its boundary-frame convention. These are strong finite reconstruction results. The finite certificate is reported in the cited construction and is not independently recomputed here.

The geometric reconstruction does not close three further estimates: restricted smoothing for the full conditional action, stability of the physical response and ground law, and bounded-congestion coverage of the translated collars with the required physical form incidence. Keeping these separate prevents a topological statement from being used as a global Poincar\'e theorem. \src{F6}{final collar ledger}; \src{F7}{initial continuation}

\begin{proposition}[No finite Wilson-action-closed update block]\label{prop:no-finite-block}
On the infinite lattice there is no nonempty finite edge set closed under adding every edge of every plaquette touching an updated edge.
\end{proposition}
\begin{proof}[Geometric argument]
Take an updated edge at an extremal coordinate of the finite set. A plaquette containing it and extending in an available transverse direction introduces an edge beyond that extremum. Repeated plaquette closure therefore cannot terminate in a nonempty finite set. The conditional formulation must retain the boundary interactions or control their influence; it cannot assume a literal finite action-closed block. \src{F7}{V1}
\end{proof}

\section{Smoothing and rarity are different estimates}
Let $\E$ be the conditional projection, $D=I-\E$, and $P$ a conditional Markov smoothing with $\E P=P\E=\E$. Suppose its restricted bad-set multiplier has the bound encoded by $\eta_B$. The proved projected smoothing inequality is
\begin{equation}\label{eq:smoothbad}
\norm{M_BDF}^2
\le(1+\theta)\eta_B\norm{DF}^2
 +(1+\theta^{-1})\ip{F}{(I-P)F},\qquad\theta>0.
\end{equation}
Its useful datum is the restricted operator stock $\eta_B$, not merely $\mu(B)$. Ground transformation also changes the conditional projections and their multipliers by explicit Perron factors. These factors must be retained before comparing bad capacity between a reference law and the physical one.

\subsection{A quantitative core--bad-set compiler}
The restricted multiplier in \cref{eq:smoothbad} can be combined with a
conditional core gap without treating rarity as a functional inequality.
Let $\mathsf E_e$ be conditional expectation on a declared enlarged
block, $\mathsf D_e=I-\mathsf E_e$, and $\mathsf P_e$ a positive
self-adjoint conditional Markov transfer with
$\mathsf E_e\mathsf P_e=\mathsf P_e\mathsf E_e=\mathsf E_e$.
All operators act on the same physical law $\mu$. Let $B_e$ be the bad
set, $K_e=B_e^c$, and let $w_e\ge0$ be the incidence weights. Norms of
conditional operators below are uniform over the admitted exteriors.

\begin{lemma}[Gluing the conditional core mean]
\label{lem:core-mean-gluing}
For a probability law $\nu$, $p=\nu(B)<1$, $q=1-p$, and $\nu g=0$,
\begin{equation}\label{eq:core-mean-gluing}
\int_{B^c}|g|^2\dd\nu
\le q\Var_{\nu(\cdot\mid B^c)}(g)
  +\frac pq\int_B|g|^2\dd\nu.
\end{equation}
Moreover, if $P$ is Markov, then
$\nu(B)\le\norm{M_BP}^2$.
\end{lemma}
\begin{proof}
Writing the two conditional means as $m_K,m_B$, the centering identity
$q m_K+p m_B=0$ gives
$q|m_K|^2\le(p/q)\int_B|g|^2\dd\nu$. Add the conditional variance.
For the second assertion apply $M_BP$ to the unit constant, using
$P\one=\one$. \src{F6}{V97}
\end{proof}

\begin{criterion}[Core gap and restricted smoothing]
\label{crit:core-smoothing}
Let $\mathcal E_\mu$ be the declared physical energy form. Assume
approximate tensorization and core incidence in the form
\begin{align}
\Var_\mu(f)&\le C_*\mathcal S(f),
 &\mathcal S(f)&=\sum_e w_e\norm{\mathsf D_ef}^2,
 \label{eq:core-tensorization}\\
\sum_e w_e\E\!\left[q_e\Var_{\mu_e(\cdot\mid K_e)}
                    (\mathsf D_ef)\right]
 &\le C_{\rm core}\mathcal E_\mu(f),
 &q_e&=\mu_e(K_e). \label{eq:core-incidence}
\end{align}
Assume also
\begin{equation}\label{eq:core-transfer-incidence}
\sup_e\norm{M_{B_e}\mathsf P_e}^2\le\eta<1,
\qquad
\sum_e w_e\ip{f}{(I-\mathsf P_e)f}
\le m_P\delta\,\mathcal E_\mu(f).
\end{equation}
Here $\delta$ is the transfer duration in the declared incidence
normalization. For any $\theta>0$ with $(2+\theta)\eta<1$,
\begin{equation}\label{eq:core-smoothing-gap}
\Var_\mu(f)
\le
\frac{C_*\{(1-\eta)C_{\rm core}
 +(1+\theta^{-1})m_P\delta\}}
 {1-(2+\theta)\eta}\,\mathcal E_\mu(f).
\end{equation}
\end{criterion}
\begin{proof}
Set $\mathcal A(f)=\sum_e w_e\norm{M_{B_e}\mathsf D_ef}^2$.
On each conditional fibre, \cref{lem:core-mean-gluing} and
$p_e\le\eta$ give
\[
\mathcal S(f)\le C_{\rm core}\mathcal E_\mu(f)
 +\frac{\mathcal A(f)}{1-\eta}.
\]
The projected smoothing estimate and
\cref{eq:core-transfer-incidence} imply
\[
\mathcal A(f)\le(1+\theta)\eta\mathcal S(f)
 +(1+\theta^{-1})m_P\delta\,\mathcal E_\mu(f).
\]
Substitute, multiply by $1-\eta$, and absorb the coefficient
$(1+\theta)\eta$. The positive remaining coefficient is
$1-(2+\theta)\eta$. Apply \cref{eq:core-tensorization}.
\src{F6}{V97}
\end{proof}

If the normalized physical core differs from a bare core with
Poincar\'e floor $s c_a$ by a potential of oscillation at most $K$,
and the gradient incidence is $m_{\rm core}$, the source comparison
supplies
\begin{equation}\label{eq:core-oscillation}
C_{\rm core}\le\frac{m_{\rm core}e^K}{s c_a}.
\end{equation}
This is useful only when the oscillation bound includes the crossing
plaquettes, response terms, and Perron profile of the complete
conditional action. Neither the collar rank nor the bare Wilson
minimum proves such a bound. Likewise, the approximate tensorization
and both incidences in \cref{crit:core-smoothing} are inputs, not
consequences of conditional site gaps. The criterion gives the exact
constant obtained when these distinct physical estimates are available.

\section{A curvature-to-defect bridge}
The later exact-block analysis constructs tree--chord fillings that express local displacement in terms of plaquette curvature. In the stated saturation convention,
\begin{equation}\label{eq:filling}
\sum_{y\in Y}D_y(z;W)
\le20\sum_{p\in\mathcal P(Y)}\norm{U_p-I}_{\rm F}^2.
\end{equation}
For a coarse-flat background, including the declared commuting-toron family, the zero-defect set is the unique moving section. For a nonflat coarse background it is empty. The section is therefore background dependent and cannot be replaced by a fixed coordinate origin without paying the induced response.

If a connected set $Y$ is $r$-bad, the filling estimate gives the Wilson-action lower bound
\begin{equation}\label{eq:Peierls}
S_{\rm W}(Y)\ge\frac{r^2}{120}|Y|.
\end{equation}
With $r_\beta=\beta^{-2/5}$ and the source small-ball normalization, the primitive far-field stock is
\begin{equation}\label{eq:qfar}
q_{\rm far}\le C_-^{-45}(\beta/120)^{180}
\exp(-\beta^{1/5}/120),
\end{equation}
while the regular Wilson remainder satisfies
\begin{equation}\label{eq:etaW}
\eta_{\rm W}\le7680\pi^3\beta^{-1/5}.
\end{equation}
The power, volume cost, and exponential belong to the same normalized conditional integral. A bare action penalty alone is not the probability estimate. \src{F8}{curvature localization and final good--bad development}; \src{F9}{initial compact rebase}

\section{What is transported, and what is not}
The collar and curvature estimates provide a finite geometric scaffold for good-field localization, rare exits, and analytic normalization. They do not bound the complete retained slow field after conditional elimination. The identity
\begin{equation}\label{eq:totalcov}
\Cov(F,G)=\E\bigl[\Cov(F,G\mid W)\bigr]
+\Cov\bigl(\E[F\mid W],\E[G\mid W]\bigr)
\end{equation}
exhibits the missing slow contribution directly. Strong mixing inside every detail fibre leaves the second term untouched. The purpose of the renormalization construction is to keep, rather than silently erase, that retained law.
\sourcebasis{Conditional score and capacity distinctions: \cite{F6,F7}. Curvature localization, moving sections, and bad-field normalization: \cite{F8,F9}.}

\chapter{Exact compact blocking and the Gaussian principal part}
\label{ch:blocking}
\section{The action is the logarithm of the exact pushforward}
For a declared compact block map $P:X\to Y$, disintegrate the reference Haar measure and define the coarse action by
\begin{equation}\label{eq:exact-block}
e^{-S_+(y)}\dd\nu_+(y)
=P_*\bigl(e^{-S(x)}\dd\nu(x)\bigr),
\end{equation}
with its normalization fixed. This is the definition of an exact blocking step. It carries all generated interactions, the appropriate Jacobian, and any retained sectors. The quadratic principal part is extracted from this object; it is not used to redefine it.

For the nonlinear root-path block, \cref{thm:global-pivot,cor:compact-product}
construct such a pushforward explicitly on a smooth compact carrier.
The Haar product is retained as an inverse-pivot amplitude; its bounds
are local per factor, not a uniform bound on the total product. Agreement
with the original block near the identity is an equality of germs, with
the separate global-law comparison specified in
\cref{scope:completion-germ}.

When a tree--chord change of variables yields product Haar measure, its normalization is exact. Overlapping path blocks do not automatically have a product law: shared edges require their common conditional kernel. Under absolute convergence, regrouping the interaction by its declared coarse footprint is exact. For the selected-tree construction in the source, the resulting norm has the schematic but explicit form
\begin{equation}\label{eq:selected-tree}
\mathfrak O_\mu(\overline\Phi)
\le D_{\rm dec}\left(\rho+\frac{rs}{1-\epsilon}\right),
\end{equation}
with $\rho,r,s,\epsilon$ the inherited primitive, marked-response, and propagation stocks. This implication requires the full source-dependent conditional polymer hypothesis. A pair-covariance estimate alone is not sufficient for all higher cumulants. \src{F7}{V92--V94}

\subsection{The exact retained sector and conditional remainder}
Let $\mathcal R_P(S)=S_+$ denote the fixed-reference blocking map in
\cref{eq:exact-block}, and let $\nu_{S,y}$ be its normalized conditional
law. Define retained pullback and conditional expectation by
\[
(\mathsf J_P q)(x)=q(Px),\qquad
(\mathsf K_S h)(y)=\E_{\nu_{S,y}}h.
\]
The reference measures and block map are held fixed in the following
variation. Any later calibration or reference reset has its own
measured contribution.

\begin{theorem}[Exact differential and retained-pullback sector]
\label{thm:exact-block-differential}
For bounded perturbations $h_1,\ldots,h_n$,
\begin{equation}\label{eq:exact-block-cumulants}
D^n\mathcal R_{P,S}[h_1,\ldots,h_n](y)
=(-1)^{n-1}\kappa_{S,y}(h_1,\ldots,h_n).
\end{equation}
In particular, $D\mathcal R_{P,S}=\mathsf K_S$ and
\begin{equation}\label{eq:retained-pullback}
\mathcal R_P(S+\mathsf J_Pq)=\mathcal R_P(S)+q,
\qquad \mathsf K_S\mathsf J_P=I.
\end{equation}
Every bounded real perturbation has the exact decomposition
$h=\mathsf J_Pm+v$, where $m=\mathsf K_Sh$ and $\mathsf K_Sv=0$, and
\begin{equation}\label{eq:exact-centered-remainder}
\mathcal R_P(S+h)-\mathcal R_P(S)
=m-\log\E_{\nu_{S,y}}e^{-v}.
\end{equation}
If $\delta_v(y)$ is the essential range length of $v$ on the fibre,
then
\begin{equation}\label{eq:centered-range-bound}
-\frac{\delta_v(y)^2}{8}
\le\mathcal R_P(S+h)(y)-\mathcal R_P(S)(y)-m(y)\le0.
\end{equation}
\end{theorem}
\begin{proof}
The difference of blocked actions is the negative logarithm of
$\E_{\nu_{S,y}}\exp(-\sum_i t_i h_i)$. Its mixed derivatives give
\cref{eq:exact-block-cumulants}. A retained pullback is constant on each
fibre and factors out of this expectation, proving
\cref{eq:retained-pullback,eq:exact-centered-remainder}. Jensen's
inequality gives the upper bound in \cref{eq:centered-range-bound}.
For the lower bound let $\psi(t)=\log\E e^{-tv}$. Its tilted variance
obeys $\psi''(t)\le\delta_v(y)^2/4$, while
$\psi(0)=\psi'(0)=0$. Integrating twice up to $t=1$ gives
$\psi(1)\le\delta_v(y)^2/8$. \src{F7}{V99}
\end{proof}

If the chosen fine and coarse norms make $\mathsf J_P$ isometric,
\cref{eq:retained-pullback} forces
$\norm{D\mathcal R_{P,S}}\ge1$. Thus strict contraction on the complete
interaction space is excluded by an exact identity, not merely by a
weak estimate. Conditional centering removes the first-order detail
response, but the pointwise bound \cref{eq:centered-range-bound} is not
a volume-uniform polymer norm: its range may be extensive. Calibrated
marginal transport, local marked estimates, and stable-complement
contraction remain separate parts of the renormalization argument.

\section{Gaussian conditioning and exact harmonic extension}
Let $L$ be positive on the declared gauge quotient and let $Q$ be a full-rank retained linear map. Put $G=L^{-1}$ and
\begin{equation}\label{eq:gauss-conditional}
G_+=QGQ^*,\qquad J=GQ^*G_+^{-1},\qquad
\Gamma=G-GQ^*G_+^{-1}QG.
\end{equation}
Then
\begin{equation}\label{eq:harmonic}
QJ=I,\qquad J^*LJ=G_+^{-1},\qquad Q\Gamma=0,\qquad\Gamma\succeq0.
\end{equation}
The retained harmonic field and the detail innovation are energy orthogonal. In particular, the exactly Gaussian detail has no residual retained feedback after the harmonic recentering. Gauge gradients and any retained flat cohomology are removed or retained according to the declared quotient; they are not inverted as if massive.

\begin{proof}[Proof of the conditioning identities]
The first two identities follow by substitution. For the last, write
\[
\Gamma=G^{1/2}(I-\mathsf P)G^{1/2},\quad
\mathsf P=G^{1/2}Q^*(QGQ^*)^{-1}QG^{1/2}.
\]
The operator $\mathsf P$ is an orthogonal projection. This also proves positivity and the annihilation by $Q$.
\end{proof}

\section{Determinant ownership and complete derivatives}
A normalized stationary Gaussian reference carries its full background-dependent determinant. It is not a small polymer vertex. For its eliminated Hessian $L(a)$,
\begin{align}
D\log\det L&=\Tr(L^{-1}DL),\label{eq:detfirst}\\
D_2D_1\log\det L
&=\Tr\bigl(L^{-1}D_2D_1L-L^{-1}D_2L\,L^{-1}D_1L\bigr).
\label{eq:detsecond}
\end{align}
The second formula includes the full return term. Taking norms before the trace words and their Ward partners are assembled loses cancellations needed by the marginal response.

These determinant derivatives are taken in the declared background of the
exact pushed action. A later change of fluctuation representative must
transport the covariance and vertex contractions coherently; a change of
background is a separate operation. In particular, the fixed-background
banks of \cref{sec:background-prescription} are not substituted for the
harmonic background of the blocked action without the comparison in
\cref{crit:background-ADMC}.

Uniformly gapped finite-range eliminated Hessians admit exponential locality bounds. The corresponding absolute four-dimensional massless response can have logarithmic growth. A locality theorem with a gap hypothesis cannot be applied to a massless retained mode by omitting that hypothesis. For a normalized conditional current, three background derivatives can require a four-entry connected cumulant. The Hessian of its pushed Haar-Jacobian term requires the third current jet, as specified in \cref{thm:connected-score-hierarchy,prop:Jacobian-third-jet}. \src{F7}{V95--V100}; \src{F8}{V60--V61}

\section{Measured coordinate changes and the stable complement}
\label{sec:measured-redundancy}
An analytic change of variables acts on the measure, not only on the
displayed action. Infinitesimally its response combines action transport
with the divergence of the coordinate vector field. The exact
pushforward identifies how such a redundant tangent is represented on
the retained field, even when the fine vector field is not vertical for
the block map.

Let $\mathcal B:X\to Y$ be a smooth submersion between compact oriented
finite-dimensional manifolds without boundary, with smooth positive
reference measures $\nu_X,\nu_Y$. Set
\[
\rho\,\dd\nu_Y=\mathcal B_*(e^{-S}\dd\nu_X),
\qquad A=-\log\rho.
\]
For a smooth fine vector field $V$, define its vector-density current
$\mathcal J_V$ weakly by
\begin{equation}\label{eq:pushed-vector-current}
\int_Y\alpha(\mathcal J_V)\dd\nu_Y
=\int_Xe^{-S(x)}\alpha_{\mathcal B(x)}
       (D\mathcal B(x)V(x))\dd\nu_X(x)
\end{equation}
for every smooth one-form $\alpha$ on $Y$.

\begin{theorem}[Measured redundancy commutes with exact blocking]
\label{thm:measured-redundancy}
For the measured tangent
$\mathcal O_V=DS[V]-\operatorname{div}_{\nu_X}V$,
\begin{equation}\label{eq:redundant-pushforward}
\mathcal B_*(\mathcal O_Ve^{-S}\dd\nu_X)
=-\operatorname{div}_{\nu_Y}\mathcal J_V\,\dd\nu_Y.
\end{equation}
If $\rho>0$ and $\overline V=\mathcal J_V/\rho$, the first variation of
$A$ under $S+\varepsilon\mathcal O_V$ is
\begin{equation}\label{eq:coarse-redundant-tangent}
\dot A=DA[\overline V]-\operatorname{div}_{\nu_Y}\overline V,
\qquad
\overline V(y)=\E[D\mathcal B\,V\mid\mathcal B=y].
\end{equation}
Consequently the exact RG differential descends to the quotient by
measured field-redefinition tangents.
\end{theorem}
\begin{proof}
Apply the divergence theorem to $e^{-S}(f\circ\mathcal B)V$:
\[
\int_X(f\circ\mathcal B)\mathcal O_Ve^{-S}\dd\nu_X
=\int_Xe^{-S}Df_{\mathcal B(x)}(D\mathcal B V)\dd\nu_X
=-\int_Y f\operatorname{div}_{\nu_Y}\mathcal J_V\dd\nu_Y.
\]
This proves \cref{eq:redundant-pushforward}. Differentiating the pushed
density gives $\dot\rho=\operatorname{div}_{\nu_Y}(\rho\overline V)$.
Since $D\rho=-\rho\,DA$, division by $-\rho$ proves
\cref{eq:coarse-redundant-tangent}. The resulting tangent has the same
measured form on $Y$, which is the required quotient invariance.
\src{F8}{V58}
\end{proof}

Strict verticality $D\mathcal B\,V=0$ is sufficient to make the pushed
current vanish, but is not necessary for redundancy. In particular,
the natural equation-of-motion vector is generally not vertical for
straight dyadic blocking. The theorem replaces that false requirement
by the conditional current of the actual pushed measure. For compact
link groups it is formulated directly in product Haar measure and does
not require a global logarithmic chart.

\subsection{The Gaussian equation-of-motion tangent}
Use the Gaussian data $L,Q,J,G_+$ of
\cref{eq:gauss-conditional}, and put $L_+=G_+^{-1}$. For the linear
fine vector field $V(x)=\tfrac12Lx$, its coarse current is
\begin{equation}\label{eq:Gaussian-coarse-velocity}
\overline V(y)=R_Vy,
\qquad R_V=\frac12QLJ=\frac12QQ^*L_+.
\end{equation}
The integration-by-parts argument remains valid for this Gaussian
law and polynomial vector field by its integrable tails.

\begin{corollary}[No Gaussian marginal feed from the measured EOM tangent]
\label{cor:Gaussian-EOM-no-feed}
The pushed tangent is
\begin{equation}\label{eq:Gaussian-EOM-owner}
\mathcal O_{\overline V}(y)
=\ip{L_+y}{R_Vy}-\Tr R_V.
\end{equation}
If $L_+(k)=O(|k|^2)$ and $QQ^*(k)=O(1)$ near zero in the declared
transverse sector, its quadratic symbol is $O(|k|^4)$. It therefore
has no transverse $F^2$ moment; the displayed constant is a vacuum
contact.
\end{corollary}
\begin{proof}
The conditional mean is $Jy$, so the conditional fine velocity is
$\tfrac12LJy$. Applying $Q$ gives
\cref{eq:Gaussian-coarse-velocity}. Insert this linear field into
\cref{eq:coarse-redundant-tangent}; its divergence is $\Tr R_V$.
The symmetric quadratic kernel is
$\tfrac12(L_+R_V+R_V^*L_+)=O(|k|^4)$ under the stated symbol bounds.
\src{F8}{V58}
\end{proof}

This is an identity for the complete measured Gaussian tangent. It does
not say that the raw Jacobian term of an arbitrary nonlinear current
has zero marginal coefficient. Its subtraction is a coordinate choice
only after the complete pushed tangent and calibrated quotient have
been identified.

\subsection{The derivative order required by the Haar Jacobian}
\label{sec:Jacobian-third-jet}
In retained product-Haar exponential coordinates, write
$\dd\nu_Y=h(y)\dd y$ and $\gamma_i=\partial_i\log h$. At the identity,
$\gamma_i(0)=0$. For $U=U^i\partial_i$,
\[
\operatorname{div}_{\rm H}U=\partial_iU^i+\gamma_iU^i,
\]
where repeated component indices are summed.

\begin{proposition}[Third-current-jet requirement]
\label{prop:Jacobian-third-jet}
If $U(0)=0$, then
\begin{equation}\label{eq:Jacobian-third-jet}
\begin{aligned}
\partial_a\partial_b\operatorname{div}_{\rm H}U\big|_0
={}&\partial_a\partial_b\partial_iU^i\big|_0\\
 &+(\partial_a\gamma_i)(0)\,\partial_bU^i(0)
  +(\partial_b\gamma_i)(0)\,\partial_aU^i(0).
\end{aligned}
\end{equation}
Thus the Hessian of the pushed Haar-Jacobian term requires the third
background jet of $\overline V$.
\end{proposition}
\begin{proof}
Differentiate the divergence twice. In the product rule for
$\gamma_iU^i$, the terms with $U(0)$ and with $\gamma_i(0)$ vanish;
the two mixed first-derivative terms remain. The divergence of $U$
contributes its third derivative. \src{F8}{V60}
\end{proof}

Apply \cref{thm:connected-score-hierarchy} to the local current
$j=D\mathcal B\,V$. Local score derivatives through order three and
connected decay for up to four marked observables provide the input
for $D^3\overline V$. With the local Haar coefficients and summability
of the support incidence, \cref{eq:Jacobian-third-jet} then gives a
local Jacobian Hessian. The explicit sufficient activity criterion is
\cref{crit:nonlinear-current-gate}. A determinant locality estimate or
a two-point covariance bound alone does not supply these hypotheses.

In four dimensions, the $F^2$ and $F\widetilde F$ quadratic structures
are marginal. The properly identified dimension-six and higher local
jets contract at tree level by at most $b^{-2}$ under a block factor
$b>1$. The physical marginal and measured redundant directions must
be projected out before asserting contraction of the stable
complement. The retained unit sector is the exact one in
\cref{thm:exact-block-differential}; invariance of the redundant quotient
is supplied by \cref{thm:measured-redundancy}. Neither statement proves
the nonlinear stable estimate by itself. \src{F7}{V99--V100};
\src{F8}{V58--V61}; \src{F9}{V9--V11}

\begin{scope}[Linear versus compact blocking]
The explicit Gaussian block-average map in \cref{ch:tower} is a linear map on the gauge-field tangent problem. A complete interacting construction must supply a compatible compact nonabelian blocking, its measure disintegration, and the uniform estimates connecting its transported vertices to the Gaussian tower. These are not consequences of the linear formulas alone.
\end{scope}
\sourcebasis{Exact coarse action, normalized Gaussian extraction, retained pullback, and conditional expansions: \src{F7}{V92--V100}. Measured redundant pushforward and the higher current/Jacobian hierarchy: \src{F8}{V58--V61}. Measured calibration and sector transport: \src{F9}{V9--V11}.}

\chapter{Compact analytic forests and nonduplicated ownership}
\label{ch:forests}
\section{Soft reference factors and compact heat kernels}
A nontrivial compactly supported analytic cutoff does not exist. The source therefore replaces differentiated hard selectors by analytic soft factors and exact add--subtract identities. The regular soft-good factor is absorbed into the stronger normalized reference. It is not reintroduced as a forest vertex. The rebased bad factor is nonnegative and is estimated in that same law.

A global logarithmic chart with a subtracted Euclidean Gaussian is invalid across compact-group sheets. The compact heat-kernel reference retains all sheets, and its replica interpolation factorizes at the endpoints required by the forest formula. Right-invariant derivatives are arranged to commute with the left-weakened generators. This avoids the derivative commutator loss of a mismatched interpolation. \src{F9}{V1--V4}

\section{The correct all-tree estimate}
Let $z$ be the primitive analytic activity on radius $R$, let $r<R$ be the output radius, let $h_*$ be the fixed local incidence stock, and let $G_\mu$ be the weighted covariance stock. Define
\begin{equation}\label{eq:KHK}
K_{\rm HK}=\frac{8h_*G_\mu}{(R-r)^2}.
\end{equation}
In a local real coordinate chart, the Taylor--Wiener norm is
\[
\norm{F}_R=\sup_{\varphi\in\R^q}\sum_{\alpha\in\N^q}
\frac{R^{|\alpha|}}{\alpha!}|\partial^\alpha F(\varphi)|.
\]
Its compact heat-kernel version uses the declared invariant derivative bank and source radii. For a multi-index $\nu$,
\[
\norm{\partial^\nu F}_r\le
\frac{\nu!}{(R-r)^{|\nu|}}\norm{F}_R.
\]
The Taylor--Wiener derivative estimate spends the available radius once for all derivatives incident at a vertex. The fixed-tree bound retains the product of vertex-degree factorials. Removing that product is neither necessary nor justified.

\begin{lemma}[Degree-faithful labeled-tree census]\label{lem:degree-census}
For $n\ge2$,
\begin{equation}\label{eq:degree-census}
\sum_{T\in\mathcal T_n}\prod_{i=1}^n d_i(T)!
=(n-2)!\binom{3n-3}{n-2}.
\end{equation}
\end{lemma}
\begin{proof}
A degree sequence $d_i\ge1$ with $\sum_i d_i=2n-2$ occurs $(n-2)!/\prod_i(d_i-1)!$ times by the Pr\"ufer code. Multiplying by $\prod_i d_i!$ gives $(n-2)!\prod_i d_i$. With $e_i=d_i-1$, the remaining sum is the coefficient of $x^{n-2}$ in
\[
\left(\sum_{e\ge0}(e+1)x^e\right)^n=(1-x)^{-2n},
\]
which is the binomial coefficient in \cref{eq:degree-census}.
\end{proof}

\begin{theorem}[Connected compact forest bound]\label{thm:forest}
Under the compact analytic reference, covariance, incidence, and endpoint-factorization hypotheses above, if $8zK_{\rm HK}<1$, then
\begin{equation}\label{eq:forest}
u_{\rm good}\le C_{\rm own}
\left(z+\frac{4z^2K_{\rm HK}}{1-8zK_{\rm HK}}\right).
\end{equation}
The degree factorials in \cref{eq:degree-census} are included in this estimate.
\end{theorem}
\begin{proof}[Proof structure]
Apply the compact interpolation forest formula with the one-radius derivative bound. Retain $\prod_i d_i!$ at fixed tree, sum it by \cref{lem:degree-census}, and include the connected-series normalization. The resulting geometric majorant has ratio $8zK_{\rm HK}$, yielding \cref{eq:forest}. The owner embedding contributes $C_{\rm own}$. \src{F9}{V1--V4}
\end{proof}

A degree-free strengthening is not an additional requirement for this argument. The degree-faithful theorem already closes the one-scale combinatorics; what remains is compatibility of its constants along the chosen regulator family. \src{F10}{V100}

\section{An explicit regulator-growth criterion}
Assume
\[
z_\beta\le C_{\rm an}\bigl(e^{c_*\beta^{-1/5}}-1\bigr),
\qquad c_*=7680\pi^3,
\qquad K_{\rm HK}\le C_K\beta^p,
\]
with constants uniform in lattice spacing and volume. Put $A=eC_{\rm an}c_*$. For $p<1/5$ and
\begin{equation}\label{eq:beta-threshold}
\beta\ge\max\left\{1,c_*^5,(16AC_K)^{1/(1/5-p)}\right\},
\end{equation}
we have $8z_\beta K_{\rm HK}\le1/2$ and
\begin{equation}\label{eq:forest-beta}
u_{\rm good}\le C_{\rm own}(A+8A^2C_K)\beta^{-1/5}.
\end{equation}
At $p=1/5$, the sufficient critical condition is $Q=8AC_K<1$, with coefficient $C_{\rm own}[A+4A^2C_K/(1-Q)]$. These follow by $e^x-1\le ex$ for $0\le x\le1$ and substitution into \cref{eq:forest}.

Thus the surviving forest compatibility question is a concrete growth estimate for $8h_*G_\mu/(R-r)^2$, not a new enumeration problem. Fixed positive radius loss and fixed covariance/incidence stocks give $p=0$. The theorem does not establish that those stocks remain fixed in a full interacting continuum construction. \src{F10}{V100}

\section{Morse gluing and all-colour scheduling}
The parameter-dependent Morse construction exactly Gaussianizes the local action around its moving stationary center. The remaining regular interaction is the transformed Jacobian and the explicitly extracted higher-order density. Exact add--subtract gluing keeps the Gaussian principal part whole and avoids derivatives of a cutoff indicator. Constant, linear, and quadratic jets are absorbed into the calibrated reference; interior third-order terms and exterior exits have different owners.

The all-colour extension requires a triangular elimination schedule.
Every primitive has one historical birth support, one saturation, and
one owner. A later colour update may use an already exposed contribution
but cannot charge it again as a new primitive. Under the common
analytic-domain and source hypotheses, this gives the complete one-scale
owner product. The marked criterion below states how its activity must
be controlled to obtain the higher normalized current derivatives; its
activity is not the physical inverse-coupling coordinate.
\src{F8}{V61--V65}; \src{F9}{V5--V8}

\section{Four-mark clustering and the nonlinear current interface}
\label{sec:four-mark-interface}
The scalar owner bound becomes a statement about differentiated
conditional currents only after local complex sources have been
included. Work in a finite volume of the actual contour graph
$\mathcal G$, with maximum degree $\Delta$, and let polymers be its nonempty finite connected
vertex sets. Two polymers are incompatible when they intersect. Assume
an exact source-dependent representation
\begin{equation}\label{eq:marked-polymer-partition}
\mathcal Z(t)=\sum_{\Gamma\ {\rm compatible}}
                 \prod_{Y\in\Gamma}z_t(Y),
\qquad |z_t(Y)|\le u^{|Y|}
\end{equation}
uniformly on the declared complex source polydisc. Source-independent
normalizing factors are kept outside $\mathcal Z$; every
source-dependent local factor must be retained in this representation
or in its explicitly bounded local contribution.

\begin{lemma}[Explicit contour-entropy criterion]
\label{lem:polymer-entropy-gate}
For $a,\mu>0$ put
\begin{equation}\label{eq:polymer-Q}
Q_{a,\mu}(u)=\frac{u e^{a+\mu}}
                 {1-\Delta^2u e^{a+\mu}}.
\end{equation}
If $\Delta^2u e^{a+\mu}<1$ and $Q_{a,\mu}(u)\le a$, then
\begin{equation}\label{eq:polymer-KP}
\sum_{Y':\,Y'\cap Y\ne\varnothing}
 |z_t(Y')|e^{(a+\mu)|Y'|}\le a|Y|.
\end{equation}
For the filling contours of \cref{eq:filling}, the graph is
$y\sim_* y'$ when $0<\norm{y-y'}_\infty\le1$, hence
$\Delta=3^4-1=80$. A sufficient choice with $a=1$ is
\begin{equation}\label{eq:four-mark-6401}
6401e^{1+\mu_*}u\le1.
\end{equation}
\end{lemma}
\begin{proof}
A deterministic depth-first traversal encodes a connected set of size
$n$ containing a fixed vertex by $2(n-1)$ steps. Its number is therefore
at most $\Delta^{2(n-1)}$. Summing these sets with the weight in
\cref{eq:marked-polymer-partition} gives $Q_{a,\mu}(u)$ at each vertex.
Sum over the vertices of $Y$ to obtain \cref{eq:polymer-KP}. For $a=1$
the condition is $(\Delta^2+1)u e^{1+\mu}\le1$, which gives
\cref{eq:four-mark-6401}. The smaller constant $65$ applies only to a
degree-eight primitive graph, not to these filling contours.
\src{F8}{V61, V63}
\end{proof}

For local supports $X_0,\ldots,X_r$, $r\le3$, assume that $z_t(Y)$
depends on $t_i$ only if $Y\cap X_i\ne\varnothing$. Let $\rho_i>0$ be
the source radii. Assume the activities are holomorphic on the source
polydisc and obey the displayed bound on every smaller closed
polydisc. Let $\tau_\mathcal G(X_0,\ldots,X_r)$ be the length
of a shortest graph tree meeting all the supports.

\begin{theorem}[Volume-uniform four-mark clustering]
\label{thm:four-mark-clustering}
Under \cref{eq:marked-polymer-partition,eq:polymer-KP}, the partition
function has a nonvanishing branch connected to $\mathcal Z=1$ at zero
activities, its connected expansion is normally convergent, and
\begin{equation}\label{eq:four-mark-clustering}
\left|\partial_{t_0}\cdots\partial_{t_r}\log\mathcal Z(0)\right|
\le 2^{r+1}|X_0|Q_{a,\mu}(u)
 \left(\prod_{i=0}^r\rho_i^{-1}\right)
 e^{-\mu\tau_\mathcal G(X_0,\ldots,X_r)}.
\end{equation}
The constant is independent of total volume.
\end{theorem}
\begin{proof}[Connected expansion and marked estimate]
Expand the hard-core logarithm in its connected Ursell coefficients.
The connected-graph deletion bound majorizes each coefficient by the
sum of incompatibility spanning trees. Rooting at a polymer $Y_0$,
\cref{eq:polymer-KP} sums the terminal branches with a total rooted
majorant at most $e^{a|Y_0|}$. This proves absolute normal convergence
and the nonvanishing logarithmic branch.

Inclusion--exclusion with $t_i=0$ selects clusters depending on every
marked source and costs at most $2^{r+1}$. Their connected union meets
every $X_i$, so their total polymer size is at least
$\tau_\mathcal G(X_0,\ldots,X_r)$. Extract the resulting exponential
weight, root at a polymer meeting $X_0$, and sum the root activities
using \cref{eq:polymer-Q}. The anchor sum contributes $|X_0|$.
Multivariate Cauchy differentiation contributes
$\prod_i\rho_i^{-1}$, proving \cref{eq:four-mark-clustering}.
\src{F8}{V61}
\end{proof}

\begin{criterion}[Complete nonlinear-current locality]
\label{crit:nonlinear-current-gate}
Suppose the exact good--bad expansion of the normalized reduced Wilson
fibre, including the current source and up to three local score
sources, satisfies \cref{eq:marked-polymer-partition} on the filling
contour graph, uniformly in the retained background chart. Assume its
nonduplicated activity obeys
\begin{equation}\label{eq:current-activity-sum}
u(\beta)\le q_{\rm far}(\beta)+u_{\rm good}(\beta),
\qquad
6401e^{1+\mu_*}\{q_{\rm far}(\beta)+u_{\rm good}(\beta)\}\le1,
\end{equation}
where the two stocks are in the same rebased law, owner convention,
and source polydisc. If the required local current and score jets
have bounded incidence and uniform local norms, then the conditional
current has three exponentially summable background jets in every
smaller decay weight. Its pushed Haar-Jacobian Hessian is therefore a
well-defined local kernel.
\end{criterion}
\begin{proof}
Apply \cref{thm:four-mark-clustering} to the mixed derivatives defining
the connected cumulants of the current and score insertions. Every
term of \cref{eq:connected-score-hierarchy} at order at most three has
at most four entries. The finite partition sum and weighted
support-tree sums give the stated locality, with the source radii and
local jet constants retained. Then use
\cref{prop:Jacobian-third-jet}. \src{F8}{V59--V63}
\end{proof}

The forest bound \cref{thm:forest} and far-field stock \cref{eq:qfar}
supply activities only through a common owner embedding in the exact
marked representation. With shared source domains and covariance
stocks, $p<1/5$ in \cref{eq:forest-beta} sends both activities to zero
and yields \cref{eq:current-activity-sum} for sufficiently large $\beta$;
the critical case retains its separate strict condition. This implication
neither permits a hard or merely smooth selector in place of analytic
rebasing nor restores the degree-free tree hypothesis.


\begin{remark}[Two distinct factorial issues]
The factorial-grade obstruction of \cref{prop:factorial} concerns the strong all-order source coefficient assigned to each labeled tree. The degree factorials of \cref{lem:degree-census} belong to a different compact analytic forest expansion and are correctly summed by its own generating coefficients. Closing the latter does not retroactively prove the former, and the former does not invalidate the latter.
\end{remark}
\sourcebasis{Analytic forests and ownership: \src{F9}{V1--V8}; corrected four-mark and current estimates: \src{F8}{V60--V65}; degree-faithful regulator growth: \src{F10}{V100}.}

\chapter{Calibrated renormalization and the physical scale}
\label{ch:rg}
\section{Coordinates and the direction of the iteration}
Write $u=\beta^{-1}$ for the declared inverse-coupling coordinate. A calibrated renormalization state consists of $u$, the retained marginal anisotropies and sector labels, and stable coordinates $z_i$ carrying specified powers $p_i$. Measured redundant directions have already been removed. The flow is not a contraction on the entire effective action: it is a controlled marginal recursion coupled to a contracting stable complement.

We use increasing $j$ for the ultraviolet orientation, so $a_j\downarrow0$ and an asymptotically free branch has
\[
\beta_{j+1}-\beta_j\longrightarrow b>0,
\qquad u_{j+1}=u_j-bu_j^2+o(u_j^2).
\]
A coarse blocking computation naturally gives the opposite orientation. Its shell coefficient must be converted before it is compared with $b$. The current response programme keeps this sign and its field normalization explicit.

\section{Invariant power cones}
Suppose the stable recursion has the nonnegative majorant
\begin{equation}\label{eq:cone-rec}
z_i^+\le\sum_jQ_{ij}z_j+d_i u^{p_i},
\qquad \frac{u^+}{u}\ge\sigma>0,
\qquad 0<u\le1.
\end{equation}
Assume $Q_{ij}\ne0$ only when $p_j\ge p_i$. Put
\[
(Q_p)_{ij}=\sigma^{-p_i}Q_{ij},\qquad
(d_p)_i=\sigma^{-p_i}d_i.
\]

\begin{criterion}[Power-preserving stable cone]\label{thm:cone}
If the spectral radius of $Q_p$ is less than one, then
\[
H=(I-Q_p)^{-1}d_p\ge0
\]
defines an invariant cone $0\le z_i\le H_i u^{p_i}$ for \cref{eq:cone-rec}. For a triangular $Q_p$, it is enough that every diagonal entry be less than one.
\end{criterion}
\begin{proof}
If $z_j\le H_ju^{p_j}$, then $u^{p_j}\le u^{p_i}$ wherever $Q_{ij}\ne0$. Consequently
\[
z_i^+\le\left(\sum_jQ_{ij}H_j+d_i\right)u^{p_i}
\le H_i(\sigma u)^{p_i}\le H_i(u^+)^{p_i}.
\]
The middle inequality is the resolvent equation for $H$. Positivity follows from the convergent Neumann series for the nonnegative matrix $Q_p$. \src{F9}{V12}
\end{proof}

\section{Reciprocal coupling shields stable errors}
The exact physical coupling reset has the form
\begin{equation}\label{eq:reciprocal}
F_r(u,x)=\frac1{u^{-1}+c_r(u^{-1},x)}.
\end{equation}
If $|c_r|\le M$, $\norm{D_xc_r}\le L$, and $u\le(2M)^{-1}$, then
\begin{equation}\label{eq:shield}
\norm{D_xF_r}\le4Lu^2.
\end{equation}
Therefore a stable displacement $x=O(u^{6/5})$ affects the physical coupling only at order $u^{16/5}$. The factor $u^2$ follows from differentiating the actual reciprocal map; it is not an independently chosen smallness assignment.

If $c_r(\beta,0)=b+O(\beta^{-\delta})+\varepsilon_r$, then expansion of \cref{eq:reciprocal} gives
\[
F_r(u,0)=u-bu^2+O(u^{2+\delta})-\varepsilon_r u^2+O(u^3)
\]
under the displayed boundedness assumptions. The exact sign of the $\varepsilon_r$ term follows the convention in $c_r$; its absolute size is the relevant error stock. \src{F9}{V12--V14}

\begin{lemma}[Ces\`aro marginal recursion]\label{lem:cesaro-beta}
Suppose a positive sequence satisfies
\[
\beta_{j+1}=\beta_j+b+d_j,\qquad b>0,
\qquad \frac1N\sum_{j<N}d_j\longrightarrow0.
\]
Then
\begin{equation}\label{eq:beta-linear}
\frac{\beta_N}{N}\longrightarrow b,\qquad
u_N:=\beta_N^{-1}\sim\frac1{bN}.
\end{equation}
\end{lemma}
\begin{proof}
Sum the recursion and divide by $N$. The fixed initial value contributes $\beta_0/N\to0$, and the assumed average error vanishes. Invert the resulting positive asymptotics.
\end{proof}

\section{Logarithmic ultraviolet control is not a power of the cutoff}
For geometric spacing $a_j=a_0B^{-j}$ and \cref{eq:beta-linear},
\begin{equation}\label{eq:clock}
\frac{\beta_j}{\log(a_0/a_j)}\longrightarrow\frac b{\log B},
\qquad
u_j\log(a_0/a_j)\longrightarrow\frac{\log B}{b}.
\end{equation}
Thus $\beta_j^{-p}$ corresponds to $(\log(a_0/a_j))^{-p}$, not to $a_j^q$ for a positive $q$. The ratio $j^{-p}B^{qj}$ diverges. Any regulator estimate that requires a power of $a_j$ needs an additional argument; the marginal recursion alone does not provide it.

\begin{proposition}[Stable memory preserves logarithmic powers]\label{prop:stable-memory}
If $x_{j+1}\le\kappa x_j+C\beta_j^{-p}$ with $0\le\kappa<1$, $p>0$, and $\beta_j/j\to b>0$, then
\begin{equation}\label{eq:stable-memory}
\limsup_{j\to\infty}j^p x_j\le\frac{Cb^{-p}}{1-\kappa}.
\end{equation}
Its Ces\`aro average is $O(N^{-p})$ for $p<1$, $O(\log N/N)$ for $p=1$, and $O(N^{-1})$ for $p>1$, apart from the exponentially decaying initial contribution.
\end{proposition}
\begin{proof}[Proof outline]
Iterate the recursion into a geometric convolution. For each fixed backward offset, $(j/\beta_{j-\ell})^p\to b^{-p}$. The recent part is bounded by a geometric series and the remote history has an exponentially small weight. This proves \cref{eq:stable-memory}; summing $j^{-p}$ gives the stated averages. \src{F11}{V1}
\end{proof}

The contraction constant in this proposition must be smaller than one. The finite Wilson covariance-derivative ceiling smaller than $4/3$ in \cref{ch:oracle} is a different object and cannot be used as $\kappa$.
\sourcebasis{Calibrated power cones, reciprocal shielding, and the marginal recursion: \src{F9}{V12--V14}. Cofinal clock and stable-tail estimates: \src{F11}{V1}.}

\chapter{Anomaly matching and the remaining scalar class}
\label{ch:anomaly}
\section{Determinant accumulation without duplicate resets}
Let $\mathfrak b$ extract the normalized transverse $F^2$ coefficient from the complete quadratic background germ. Before a reference reset, exact nested Gaussian elimination gives
\begin{equation}\label{eq:det-telescope}
\Gamma_G^{[0,N)}=\sum_{j<N}\Gamma_{G,j}+C_N,
\qquad \mathfrak b(C_N)=0,
\end{equation}
where $C_N$ is background independent in the declared extraction. A reset may create an additional endpoint or stable-history term; it must be included once, in its own ledger.

\begin{definition}[Averaged determinant matching card]\label{def:ADMC}
The averaged determinant matching card, abbreviated ADMC, requires
\begin{equation}\label{eq:ADMC}
\mathfrak b\bigl(\Gamma_G^{[0,N)}\bigr)=Nb_{\rm YM}+o(N),
\qquad R_N^{\rm reset}+R_N^{\rm boundary}=o(N),
\end{equation}
with one common compact determinant, Haar, quotient, sector, and normalization convention throughout the block family.
\end{definition}

Together with the calibrated analytic domain and the marked-response bounds, ADMC supplies the Ces\`aro error required by \cref{lem:cesaro-beta}. It is stronger than a correct coefficient at one scale and weaker than demanding a pointwise error estimate at every scale. The nonlinear inverse-coupling remainder is controlled separately; it is not inserted into the principal determinant a second time. \src{F9}{V13--V15}

\section{Compact--continuum intertwiners}
\begin{definition}[Sublinear compact--continuum intertwiner card]\label{def:SCIC}
The sublinear compact--continuum intertwiner card, abbreviated SCIC, is a decomposition
\begin{equation}\label{eq:SCIC}
\Gamma^{\rm cmp}_j-\Gamma^{\rm cov}_j
=\mathcal A_{j+1}-\mathcal A_j+\mathcal E_j+C_j,
\end{equation}
with
\[
\mathfrak b(C_j)=0,\qquad
\mathfrak b(\mathcal A_N)=o(N),\qquad
\frac1N\sum_{j<N}\mathfrak b(\mathcal E_j)\longrightarrow0,
\]
uniformly over the retained sectors needed by the application.
\end{definition}

If the covariant comparison has averaged coefficient $b_{\rm YM}$, telescoping \cref{eq:SCIC} supplies the Gaussian part of ADMC. An exactly measured change of variables can contribute a coboundary, and an exactly compensated pure-gauge quotient can contribute no anomaly. These facts do not show that every compact--continuum difference is one of those terms.

The Ward-recombined homotopy formulation makes the residual explicit. For the complete scalar response $W_{j,t}(k)$,
\begin{equation}\label{eq:Wardhomotopy}
\partial_tW_{j,t}
=\operatorname{div}_kV_{j,t}
+D_{j+1,t}-D_{j,t}+R_{j,t}+C_{j,t},
\qquad \avg{C_{j,t}}=0.
\end{equation}
Periodic divergences integrate to zero. Artificial low/edge partition derivatives cancel only after both pieces have been recombined. With sublinear endpoint and Ces\`aro-null residual integrals, \cref{eq:Wardhomotopy} implies SCIC. \src{F9}{V15--V16}

\section{Analytic alias cohomology}
Use normalized Haar average on $\T^4$ and define the dyadic inverse-branch operator
\begin{equation}\label{eq:alias-P}
(\mathcal Pf)(q)=\frac1{16}\sum_{\sigma\in\{0,1\}^4}
f(q/2+\pi\sigma).
\end{equation}
For the analytic Wiener norm
\[
\norm{f}_\rho=\sum_{m\in\Z^4}e^{\rho|m|_1}|\widehat f(m)|,
\]
one has $\avg{\mathcal Pf}=\avg f$ and
\begin{equation}\label{eq:Fourier-decimation}
\widehat{\mathcal P^nf}(m)=\widehat f(2^nm).
\end{equation}

\begin{theorem}[The constant mode is the only analytic obstruction]\label{thm:alias}
If $\avg f=0$ and $0<\rho'\le\rho$, then
\begin{equation}\label{eq:alias-decay}
\norm{\mathcal P^nf}_{\rho'}
\le e^{-(\rho2^n-\rho')}\norm f_\rho,\qquad n\ge1.
\end{equation}
Consequently every analytic real scalar response has the unique decomposition
\begin{equation}\label{eq:alias-decomp}
f=\avg f+(I-\mathcal P)a,\qquad \avg a=0,
\end{equation}
where $a=\sum_{n\ge0}\mathcal P^n(f-\avg f)$ converges in each stated smaller analytic norm. The quotient by analytic alias coboundaries is one dimensional.
\end{theorem}
\begin{proof}
For $m\ne0$, the ratio between the output and input analytic weights is bounded by $e^{-(\rho2^n-\rho')}$ because $|m|_1\ge1$. Apply \cref{eq:Fourier-decimation} and sum to prove \cref{eq:alias-decay}. The Poisson series converges, and its partial sums telescope under $I-\mathcal P$. If $\mathcal Pa=a$, then $\widehat a(m)=\widehat a(2m)$; absolute summability forces all nonzero Fourier coefficients to vanish. Haar invariance forces the constant in \cref{eq:alias-decomp}. \src{F9}{V17}
\end{proof}

For a uniformly analytic, fully Ward-recombined edge family $f_{j,t}$ in the common rescaled variable, put $m_{j,t}=\avg{f_{j,t}}$. Under the source's compatible edge-transport identification, the nonconstant part gives a bounded endpoint cochain. The exact remaining edge criterion is
\begin{equation}\label{eq:edge-scalar}
\frac1N\sum_{j<N}\int_0^1m_{j,t}\,\dd t\longrightarrow0.
\end{equation}
If $m_{j,t}\to m_{\infty,t}$ in $L^1(\dd t)$, it reduces to $\int_0^1m_{\infty,t}\dd t=0$. The analytic cohomology theorem does not evaluate that integral. Subtracting an unknown mean in order to solve the Poisson equation is not a proof that the mean vanishes.

\section{What a finite quadrature can certify}
For the $M^4$ equally spaced grid on $\T^4$, let $\langle f\rangle_M$ denote its normalized mean. The source analytic stock gives the explicit certificate
\begin{equation}\label{eq:quadrature}
\abs{\langle f\rangle_M-\avg f}
\le\norm f_\rho
\left[\left(\frac{1+e^{-\rho M}}{1-e^{-\rho M}}\right)^4-1\right].
\end{equation}
The proof sums precisely the nonzero Fourier modes divisible by $M$ in each coordinate. Thus a finite grid can certify a mean when paired with a uniform analytic norm. A grid value without that stock is not a Haar integral, and a pointwise sign at finitely many momenta is not an all-torus sign theorem.

\section{The obstruction to a fixed-reference comparison}
The finite response calculations produce exact local algebra. They do not license a global Neumann approximation outside its certified domain. The finite test in \src{F11}{V31} finds that the relevant frozen-chart residual test holds at only one of the $256$ quarter-turn points and fails at the other $255$. A constant Laurent coefficient of a truncated polynomial remains exact as polynomial algebra, but is not an approximation to the rational Haar mean without a global remainder bound.

There is also a more structural obstruction. If every rescaled step uses the same frozen endpoint pair and its scalar mean difference is a constant $m_G$, then the average in \cref{eq:edge-scalar} is $m_G$. A certified nonzero value would refute SCIC for that reset comparison, not prove anomaly matching. The correct nested relation has the form
\begin{equation}\label{eq:nested-shell-cob}
b_{G,j}=b_{\rm YM}+c_j-c_{j+1}+e_j,
\qquad c_N=o(N),\qquad \frac1N\sum_{j<N}e_j\to0.
\end{equation}
The nested Gaussian tower of \cref{ch:tower} avoids this repeatedly reset comparison; it is not inferred from a frozen-chart sign calculation. \src{F11}{V31--V32}
\sourcebasis{ADMC, SCIC, Ward recombination, and alias cohomology: \src{F9}{V13--V17}. Correct next marginal row: \src{F10}{V100}. Frozen-chart limitations and nested replacement: \src{F11}{V31--V34}.}

\chapter{The finite response oracle and its actual scope}
\label{ch:oracle}
\section{The reduced Hessian and its full edge origin}
The parity-tree calculation starts from a $64$-edge fibre and removes the declared gauge directions to obtain a $45$-coordinate reduced Hessian. The response oracle retains the noncommuting colour algebra, harmonic routers, cubic and quartic vertices, and the retree transformations. A coordinate restriction is valid only for the vertical space for which it was constructed. A different block map can require full-edge vertex banks before projection.

The source gives all-torus floors
\begin{equation}\label{eq:finite-H-floor}
H_{\rm W}(k)\succeq\frac{49}{1250}I_{45},\qquad
H_*(k)\succeq\frac{343}{1250}I_{45}
\end{equation}
for its declared endpoint Hessians. The proof uses the exact alias-energy and reciprocal estimates, not interpolation between a finite set of samples. These are finite eliminated-fibre statements, not floors for a physical continuum Hamiltonian. \src{F9}{V63 and endpoint development}; \src{F10}{inherited oracle}

\section{Moving-frame covariance}
Let $V$ be the normalized alias isometry. In the source convention, the frame connection is
\[
\Gamma_\mu=-\frac{\ii}{2}\diag(s_\mu),\qquad
\boldsymbol\nabla_\mu X=\partial_\mu X-[X,\Gamma_\mu].
\]
The full compressed symbol satisfies the exact moving-frame derivative identity, while commutators telescope in complete typed traces. A derivative of an individual matrix entry is not the corresponding covariant response. Likewise, an isolated fine $4\times4$ curl block has a moving gauge kernel; dropping the $45$-dimensional constrained range can destroy the positivity needed by a proposed estimate.

The tree-gauge map used for the tower is not merely a unitary change of
frame inside this compressed symbol. It selects a representative of the
fine gauge quotient before contraction with fixed fluctuation-leg banks.
The finite discrepancy in \cref{prop:representative-dependence} therefore
cannot be removed by invoking moving-frame covariance alone. The
covariance must be pushed through the specified representative map.

\section{A global relative covariance-derivative margin}
Write the exact finite Laurent representation
\[
H_{\rm W}(k)=\sum_{\alpha\in\mathcal S}G_\alpha e^{\ii\alpha\cdot k},
\qquad \overline M_{\rm W}=\sum_{\alpha\in\mathcal S}\norm{G_\alpha}_{\rm F}.
\]
The completed finite atlas reports the strictly positive rational margin
\begin{equation}\label{eq:theta97}
\theta_{97}=
\frac{23784808996629854813771747874389}
{241558005344776833024000000000000000000}
\end{equation}
for
\begin{equation}\label{eq:relative-floor}
\frac43H_{\rm W}\pm\boldsymbol\nabla_1H_{\rm W}\succeq\theta_{97}I_{45}.
\end{equation}

\begin{theorem}[Normalized derivative consequence]\label{thm:covariance-margin}
For this fixed-cutoff Wilson fibre,
\begin{align}
2H_{\rm W}\pm\boldsymbol\nabla_1H_{\rm W}
&\succeq\left(\frac{49}{1875}+\theta_{97}\right)I_{45},\label{eq:H2margin}\\
\norm{H_{\rm W}^{-1/2}(\boldsymbol\nabla_1H_{\rm W})H_{\rm W}^{-1/2}}_{\op}
&\le\kappa_{\rm W}:=\frac43-\frac{\theta_{97}}{\overline M_{\rm W}}<\frac43.
\label{eq:kappaW}
\end{align}
The same bound holds for
$H_{\rm W}^{1/2}(\boldsymbol\nabla_1H_{\rm W}^{-1})H_{\rm W}^{1/2}$.
\end{theorem}
\begin{proof}
Add $(2/3)H_{\rm W}$ to \cref{eq:relative-floor} and use \cref{eq:finite-H-floor}. Since $H_{\rm W}\preceq\overline M_{\rm W}I$, congruence of \cref{eq:relative-floor} by $H_{\rm W}^{-1/2}$ gives eigenvalue bounds $\pm(4/3-\theta_{97}/\overline M_{\rm W})$ for the normalized derivative. Finally,
\[
\boldsymbol\nabla_1H^{-1}=-H^{-1}(\boldsymbol\nabla_1H)H^{-1}
\]
proves the inverse-covariance assertion. \src{F10}{V97, V100}
\end{proof}

The rational atlas is a reported finite certificate. The consequence just proved is an elementary deduction from that certificate and the stated Hessian floor. Neither inequality evaluates the complete Ward scalar. In particular, $\kappa_{\rm W}<4/3$ is not a contraction smaller than one, not a regulator-uniform Schwinger-function derivative bound, and not a mass gap.

\section{The retained role of the oracle}
The oracle supplies exact vertices, normalized derivative control, and
finite tests of complete trace-word cancellation. For the block-average
tower, its extension to the full fine-edge space is necessary but not
sufficient: the shell covariance must also be put in the tree-gauge
representative used by the banks. The resulting corrected finite route
agrees with the slice route at the reported nonzero corners. The fixed
background and the block-rule harmonic background then give different
second-order corner means, recorded separately in
\cref{sec:fixed-background-cards}. This preserves the finite oracle as a
usable component without confusing its cards with a physical anomaly.
\sourcebasis{The finite oracle and its global relative margin are
developed in \cite{F9,F10}; full-edge reconstruction and the representative
correction are in \src{F11}{V39, V41}.}

\chapter{The explicit nested block-average Gaussian tower}
\label{ch:tower}
\section{Why this is a new principal object}
The earlier base-point path rules solve the axis problem and, after symmetrization, yield an off-axis fixed point with contraction rate $17/36$. Their global nondegeneracy requires an additional inequality. The block-average rule changes the linear map upstream: its alias blocks are diagonal, so the covariance tower and the limiting nondegeneracy can be written as positive lattice sums. The unresolved inequality for the older rule is not an extra hypothesis of the new one. \src{F11}{V32--V37}

Let
\begin{equation}\label{eq:qPi}
q_\mu(k)=e^{\ii k_\mu}-1,\qquad
\lambda(k)=\sum_{\mu=1}^{4}|q_\mu(k)|^2,
\qquad \Pi(k)=I-\frac{q(k)q(k)^*}{\lambda(k)}
\end{equation}
for $k\in\T^4\setminus\{0\}$. The free transverse covariance is
\[
G^{(0)}(k)=\frac{\Pi(k)}{\lambda(k)}.
\]
All inverses below are transverse inverses, or equivalently the indicated Moore--Penrose inverses on the gauge quotient.

\section{The block rule and its Ward identity}
For a cell $\{0,1\}^4$ and direction $\mu$, define
\begin{equation}\label{eq:block-average}
c_\mu=\frac1{16}\sum_{x\in\{0,1\}^4}
\bigl(a_{(x,\mu)}+a_{(x+e_\mu,\mu)}\bigr).
\end{equation}
Links leaving the cell carry their neighbor-cell phases. For the aliases $p_\epsilon=k/2+\pi\epsilon$, $\epsilon\in\{0,1\}^4$, put
\begin{equation}\label{eq:alias-B}
s_\epsilon=\frac1{16}\prod_{\nu=1}^{4}(1+e^{\ii p_{\epsilon,\nu}}),
\qquad
\Delta_\epsilon=\diag_\mu(1+e^{\ii p_{\epsilon,\mu}}),
\qquad B_\epsilon=s_\epsilon\Delta_\epsilon.
\end{equation}
The basic Ward identity is
\begin{equation}\label{eq:block-Ward}
B_\epsilon(k)q(p_\epsilon)=s_\epsilon q(k).
\end{equation}
It follows from $(1+w_\mu)(w_\mu-1)=w_\mu^2-1$. Thus the average of a fine pure gauge is a coarse pure gauge, and fine gauge changes with zero averaged potential are annihilated.

In the slice normalization of the source, the covariance recursion is
\begin{equation}\label{eq:tower-rec}
G^{(j+1)}(k)=\frac1{16}\Pi(k)
\sum_{\epsilon}B_\epsilon(k)G^{(j)}(p_\epsilon)B_\epsilon(k)^*\Pi(k).
\end{equation}
The segment family in \cref{eq:block-average} is covariant under coordinate permutations and reflections up to the stated cell-translation phase. This gives the correct hypercubic symmetry convention; reflection does not mean merely changing a momentum sign while leaving the edge frame fixed.

\section{A closed formula at every level}
For $x=k+2\pi M$ set
\begin{equation}\label{eq:Fj}
F_\nu^{(j)}(x)=\prod_{\ell=1}^{j}|1+e^{\ii x_\nu/2^\ell}|^2
=\frac{\sin^2(k_\nu/2)}{\sin^2(x_\nu/2^{j+1})}.
\end{equation}
At $k_\nu=0$, the convention is $F_\nu^{(j)}=4^j$ when $M_\nu\equiv0\pmod{2^j}$ and zero otherwise. Empty products at $j=0$ equal one.

\begin{theorem}[Explicit finite-level propagator]\label{thm:tower-explicit}
For every $j\ge0$ and $k\ne0$,
\begin{equation}\label{eq:Gj}
G^{(j)}(k)=\Pi(k)\diag_\mu\bigl(\gamma_\mu^{(j)}(k)\bigr)\Pi(k),
\end{equation}
where
\begin{equation}\label{eq:gammaj}
\gamma_\mu^{(j)}(k)=16^{-3j}
\sum_{M\bmod2^j}
\frac{\bigl(\prod_\nu F_\nu^{(j)}(k+2\pi M)\bigr)
F_\mu^{(j)}(k+2\pi M)}
{\lambda((k+2\pi M)/2^j)}.
\end{equation}
\end{theorem}
\begin{proof}
Iterate \cref{eq:tower-rec}. The chain block $W_M$ is the product of the diagonal $B$ matrices along the alias path. The inner fine transverse projection can be removed after the outer compression because the iterated Ward identity sends its gauge direction to $q(k)$. The squared scalar part of $W_M$ is $16^{-2j}\prod_\nu F_\nu^{(j)}$, its directional squared part is $F_\mu^{(j)}$, and the slice covariance contributes $16^{-j}$. This gives \cref{eq:gammaj}. \src{F11}{V37}
\end{proof}

\section{The positive lattice-sum fixed point}
Define
\[
\sigma_\nu(k,M)=\frac{\sin^2(k_\nu/2)}{(k_\nu+2\pi M_\nu)^2},
\]
with $\sigma_\nu=(1/4)\delta_{M_\nu,0}$ at $k_\nu=0$.

\begin{theorem}[Explicit fixed point and uniform difference control]\label{thm:tower-fixed}
The limiting propagator is
\begin{align}
G^*(k)&=\Pi(k)\diag_\mu\bigl(\Gamma_\mu^*(k)\bigr)\Pi(k),\label{eq:Gstar}\\
\Gamma_\mu^*(k)&=1024\sum_{M\in\Z^4}
\frac{\bigl(\prod_\nu\sigma_\nu(k,M)\bigr)\sigma_\mu(k,M)}
{|k+2\pi M|^2}.
\label{eq:positive-lattice-sum}
\end{align}
The series is absolutely convergent with nonnegative terms for $k\ne0$. If
\[
d_1=\sup_{k\ne0}\norm{G^{(1)}(k)-G^{(0)}(k)},
\]
then $d_1<\infty$ and
\begin{equation}\label{eq:tower-rate}
\sup_{k\ne0}\norm{G^{(j)}(k)-G^*(k)}
\le\frac43d_1\,4^{-j}.
\end{equation}
The uniform statement concerns the bounded difference of covariances sharing an infrared singularity; it does not assert that $G^*$ itself is bounded near zero.
\end{theorem}
\begin{proof}[Proof structure]
Choose symmetric representatives of the aliases in \cref{eq:gammaj}. The sine inequalities on $[-\pi/2,\pi/2]$ compare each summand to a fixed multiple of the positive summand in \cref{eq:positive-lattice-sum}; the coordinatewise product gives the summable lattice majorant. Taking the limit yields the displayed series.

For the linear covariance map $\mathcal L$ in \cref{eq:tower-rec}, direct alias summation gives
\begin{equation}\label{eq:Lcontraction}
\mathcal L[\Pi](k)=\frac14\Pi(k)
\diag_\mu\left(1-\frac{|q_\mu(k)|^2}{8}\right)\Pi(k)
\preceq\frac14\Pi(k).
\end{equation}
By positivity, this controls any bounded transverse difference by a factor $1/4$. The common infrared leading term makes the first difference bounded. Therefore consecutive differences are bounded by $d_1 4^{-j}$; summing their geometric tail gives \cref{eq:tower-rate}. \src{F11}{V35, V37}
\end{proof}

\begin{theorem}[Global transverse nondegeneracy]\label{thm:tower-floor}
Let $c_0=256/\pi^{12}$. Then
\begin{equation}\label{eq:Gfloor}
G^*(k)\succeq c_0\Pi(k),\qquad
\norm{S^*(k)}\le\frac{\pi^{12}}{256},\qquad S^*=(G^*)^+.
\end{equation}
For
\[
c_j=c_0-\frac43d_1 4^{-j}>0,
\]
the kernels $S^{(j)}=(G^{(j)})^+$ satisfy
\begin{equation}\label{eq:Srate}
\sup_{k\ne0}\norm{S^{(j)}(k)-S^*(k)}
\le\frac{d_1 4^{-j}}{(3/4)c_j^2}.
\end{equation}
\end{theorem}
\begin{proof}
Take the $M=0$ summand in \cref{eq:positive-lattice-sum} with $k\in[-\pi,\pi]^4$. Each of the five sine ratios is at least $\pi^{-2}$, and $|k|^2\le4\pi^2$, giving $1024\pi^{-10}/(4\pi^2)=c_0$. Compression by $\Pi$ proves \cref{eq:Gfloor}. The difference estimate gives the finite-level floor $c_j\Pi$. The inverse identity and \cref{eq:tower-rate} bound the inverse difference by $(4/3)d_1 4^{-j}/c_j^2$. \src{F11}{V37}
\end{proof}

\begin{scope}[Nondegeneracy is not a mass]
The bound is a lower bound on the covariance and an upper bound on the kernel. On a coordinate axis, for a transverse component,
\begin{equation}\label{eq:axis}
G^*_{22}(t,0,0,0)=\frac{2+\cos t}{12\sin^2(t/2)},
\qquad
S^*_{22}(t,0,0,0)=\frac{12\sin^2(t/2)}{2+\cos t}
=t^2+O(t^4).
\end{equation}
The Gaussian fixed point is infrared massless. Its uniform invertibility away from the gauge direction, in the precise sense of \cref{eq:Gfloor}, is useful analytic control and not the conclusion of the Yang--Mills mass-gap problem.
\end{scope}

\section{Harmonic extensions and positive shells}
Let $p_M=(k+2\pi M)/2^j$ and
\[
W_M=\prod_{\ell=1}^{j}B(e^{\ii(k+2\pi M)/2^\ell}).
\]
The minimizing fine field generated by a level-$j$ transverse coarse field $c$ is
\begin{equation}\label{eq:tower-minimizer}
a(p_M)=\mathcal M_M^{(j)}(k)c,
\qquad
\mathcal M_M^{(j)}=16^{-j}G^{(0)}(p_M)W_M^*S^{(j)}(k).
\end{equation}
The exact constraint and energy identities are
\begin{equation}\label{eq:tower-min-identities}
\Pi\sum_MW_M\mathcal M_M^{(j)}=\Pi,
\qquad
16^j\sum_M(\mathcal M_M^{(j)})^*S^{(0)}(p_M)\mathcal M_M^{(j)}=S^{(j)}.
\end{equation}
Successive minimizer maps compose. Intermediate transverse projectors cancel in the appropriate places by the Ward identity, not by treating them as identity matrices on the full edge space.

On a fixed finest alias class, let $\mathcal C_\ell$ denote the part of the fine covariance explained by the level-$\ell$ retained field. At the top level $j$ its blocks are
\begin{equation}\label{eq:explained-cov}
\mathcal C_j(M,M')=16^{-2j}G^{(0)}(p_M)W_M^*S^{(j)}W_{M'}G^{(0)}(p_{M'}).
\end{equation}
The unrestricted fine covariance in these units is $16^{-j}G^{(0)}(p_M)\delta_{MM'}$.

\begin{theorem}[Positive Gaussian shell decomposition]\label{thm:shells}
The differences
\[
\Gamma_\ell=\mathcal C_\ell-\mathcal C_{\ell+1}
\]
are positive semidefinite and
\begin{equation}\label{eq:shell-telescope}
\mathcal C_0=\sum_{\ell=0}^{j-1}\Gamma_\ell+\mathcal C_j.
\end{equation}
They are the one-step conditional vertical covariances lifted to the fine field. In the fixed-fine-momentum sense stated in the source, the explained tail tends to zero, giving the corresponding infinite shell expansion of the free covariance.
\end{theorem}
\begin{proof}
For nested Gaussian retained sigma fields, the explained covariances are obtained by orthogonal projection in the Gaussian Hilbert space. Projection to a coarser retained field can only reduce the explained covariance. Its difference is the covariance of the conditional innovation, hence positive. Telescope the differences and apply the source's explained-tail limit. The explicit formulas follow from \cref{eq:tower-minimizer}. \src{F11}{V38}
\end{proof}

The one-step vertical block formula makes its normalization explicit:
\begin{equation}\label{eq:vertical}
\begin{aligned}
\operatorname{Cov}^{\rm vert}_{\epsilon\epsilon'}
={}&\frac{\delta_{\epsilon\epsilon'}}{16}G^{(j)}(p_\epsilon)\\
&-\frac{\overline s_\epsilon s_{\epsilon'}}{256}
G^{(j)}(p_\epsilon)\Delta_\epsilon^*S^{(j+1)}(k)
\Delta_{\epsilon'}G^{(j)}(p_{\epsilon'}).
\end{aligned}
\end{equation}
It is positive and annihilated by the retained block map. Shell positivity gives an exact Gaussian multiscale decomposition; it does not give a nonlinear measure or a physical transfer contraction.

\section{Full-edge vertices and the fluctuation representative}
\label{sec:shell-functional}
The positive shells of \cref{thm:shells} are defined on transverse alias
spaces. A background-response calculation requires, in addition, a
representative of each fluctuation and a background at which the vertices
are differentiated. These are distinct choices. Positivity of a covariance
on the gauge quotient does not make a partially assembled vertex
contraction independent of its fluctuation representative.

Let $Q_0(b)$ and $R_0(b)$ denote the full fine-edge quadratic and linear
background vertices, with the colour trace and the paired-vertex signs
fixed as in the finite response construction. The full-edge construction
extends the cubic and quartic banks from the $45$ retained coordinates to
all $64$ edges. Their restrictions reproduce the original banks. The
additional columns comprise fifteen tree links and four second links; the
cubic extension has $378$ nonzero entries across the three
external-momentum orders. These are reported exact finite identities of
\src{F11}{V39}; they make the vertex forms available on the $49$-coordinate
tree slice before the coarse constraint is imposed.

\begin{proposition}[Representative dependence of the finite response]
\label{prop:representative-dependence}
The Wilson response formed by contracting these vertex banks with a
transverse alias covariance need not equal the response in the tree-gauge
slice, even when the two fluctuation fields represent the same gauge
class. In the stated normalization, the reported order-zero, order-one,
and order-two cards for the straight rule at $k=(\pi,0,0,0)$ are
\begin{align}
\mathfrak c_{\perp}(k)
 &=\left(-\frac{118777}{51450},\ 0,
          -\frac{3401839261}{1512630000}\right),
 \label{eq:transverse-corner}\\
\mathfrak c_{\rm tree}(k)
 &=\left(\frac{636}{245},\ 0,
          -\frac{1110886837}{43218000}\right).
 \label{eq:tree-corner}
\end{align}
Consequently, direct substitution of the transverse covariance into the
fixed tree-gauge vertex banks is not a valid representative-independent
prescription.
\end{proposition}
\begin{proof}[Finite-certificate implication]
The unequal zeroth-order entries already disprove representative
independence. The exact cards are reported in \src{F11}{V41}; their
full computational derivation is not repeated here. The distinction
concerns the fluctuation legs of the specified response banks. It is not a
failure of gauge invariance of the complete Wilson action or of a fully
matched change of variables in the measure.
\end{proof}

\section{The tree-gauge map and covariance transport}
\label{sec:tree-gauge}
Fix a nonzero coarse Bloch momentum $k$, a rooted spanning tree of the
sixteen sites of a cell, and the coarse-link map $L$ of the chosen rule.
Let $a$ be a cell edge field. Integrate $a$ along the fifteen tree links to
obtain a site potential $\phi_a$ with $\phi_a(0)=0$, and extend that
potential to neighbouring cells with the Bloch phases. Then
$a'=a-d\phi_a$ vanishes on the tree. Define the residual gauge field
\[
v_{(s,\mu)}(k)=q_\mu(k)\one_{\{s_\mu=1\}},
\]
which is generated by a constant potential on each cell, and set
\begin{equation}\label{eq:tree-gauge-map}
c_a(k)=\frac{q(k)^*L(a')}{\lambda(k)},\qquad
T_{\rm g}(k)a=a-d\phi_a-c_a(k)v(k).
\end{equation}
This formula specifies the map on $k\ne0$; it does not prescribe an inverse
of $\lambda$ on the toron sector.

\begin{proposition}[Exact representative constraints]
\label{prop:tree-gauge-constraints}
The map $T_{\rm g}$ is linear. Its image has zero tree links and zero
longitudinal coarse link, and
\begin{equation}\label{eq:tree-gauge-constraints}
q(k)^*L(T_{\rm g}a)=0,\qquad
\Pi(k)L(T_{\rm g}a)=\Pi(k)L(a).
\end{equation}
In particular, a transversely vertical field is sent to the full
coarse-link constraint $L(T_{\rm g}a)=0$. On the specified nonzero momentum
sector this is the unique representative in its fine gauge class with
these constraints, and $T_{\rm g}^2=T_{\rm g}$.
\end{proposition}
\begin{proof}
Tree integration is linear and fixes the potential relative to the root.
The residual field $v$ vanishes on the tree and obeys $L(v)=q(k)$. The
Ward identity of the block map sends every fine gradient to a coarse
gradient. Thus subtracting $d\phi_a$ and $c_av$ preserves the transverse
coarse link, while the definition of $c_a$ removes its longitudinal
component. If two representatives in the same fine gauge class have zero
tree links, their potential difference is constant within each cell. Their
remaining difference is therefore a multiple of $v$. The longitudinal
condition forces that multiple to vanish because $\lambda(k)>0$.
Applying the map to a field already satisfying the constraints makes both
subtractions zero. These are the representative identities underlying the
construction of \src{F11}{V41}.
\end{proof}

The normalization of the frame change is part of the covariance transport.
At level zero write
\[
D_0(k)=\diag_{\epsilon}G^{(0)}(p_\epsilon),\qquad
\mathsf B(k)=\bigl(B_\epsilon(k)\bigr)_\epsilon,
\]
and let
$E_{(s,\mu),(\epsilon,\nu)}=e^{\ii p_\epsilon\cdot s}\delta_{\mu\nu}$.
In the colour and Fourier convention of the response banks,
\begin{align}
\widehat C_0(k)
 &=3\left[D_0(k)-\frac1{16}D_0(k)\mathsf B(k)^*
 S^{(1)}(k)\mathsf B(k)D_0(k)\right],
 \label{eq:response-alias-cov}\\
C^{\rm edge}_0(k)&=\frac1{16}E(k)\widehat C_0(k)E(k)^*,
 \label{eq:response-edge-cov}\\
\Xi_0(k)&=T_{\rm g}(k)C^{\rm edge}_0(k)T_{\rm g}(k)^*.
 \label{eq:response-tree-cov}
\end{align}
The factor $3$ is the inverse of the one-colour Hessian normalization; the
factor $1/16$ in the edge conversion is a Fourier-frame factor. They are
not interchangeable. In particular, the bracket in
\cref{eq:response-alias-cov} is sixteen times the level-zero covariance in
\cref{eq:vertical}; this reconciles the two displayed covariance
conventions without inserting a further colour factor into the vertices.

\begin{proposition}[Covariance transport on a fixed quotient]
\label{prop:tree-covariance}
For a centered edge field with covariance $C^{\rm edge}\succeq0$, its
representative $T_{\rm g}a$ has covariance
\[
\Xi=T_{\rm g}C^{\rm edge}T_{\rm g}^*\succeq0.
\]
If the original covariance is supported on the transversely vertical
space, then $L\Xi=0$. When an alias calculation and a slice calculation
realize the same quotient Gaussian law in this unique representative,
their covariances and their contractions with the same vertex banks agree.
\end{proposition}
\begin{proof}
Covariance is transported by linear congruence. For every test vector $h$,
$\ip{h}{\Xi h}=\ip{T_{\rm g}^*h}{C^{\rm edge}T_{\rm g}^*h}\ge0$.
The constraint follows from \cref{prop:tree-gauge-constraints}. Equality
for two realizations of the same quotient law follows from uniqueness of
the representative; equality of their matrix contractions is then
algebraic. The same-quotient-law premise is essential.
\end{proof}

For the finite level-zero response, the exact-arithmetic comparison in
\src{F11}{V41} reports agreement of the corrected alias route with the
straight-rule reference cards and with the block-average slice route at
all fifteen nonzero corners and all three external-momentum orders.
The zero-momentum toron corner is treated by the regular slice route:
the transverse alias representation there has a singular constant-alias
kernel, so the ordinary Taylor inversion used away from zero is not
applicable. The finite agreement does not certify a full-torus quadrature.
The precise punctured-torus derivative bounds, rather than a uniform
pointwise bound at the toron, are given in
\cref{prop:toron-regularity}.

For a higher shell, let $C_j^{\rm edge}$ denote the lifted covariance
$\Gamma_j$ after the corresponding edge-frame and colour-normalization
conversion. The required covariance in the fixed fine-edge banks is
\begin{equation}\label{eq:nested-tree-cov}
\Xi_j=T_{{\rm g},j}C_j^{\rm edge}T_{{\rm g},j}^*.
\end{equation}
The nested representative is fixed scale by scale, first on the retained
lattice and then on each finer cell. This gives an algebraic prescription;
regulator-uniform estimates for that prescription and its momentum jets
remain separate analytic requirements. Positivity of each congruence does
not furnish those derivative bounds.

\subsection{Toron regularity of the representative}
\label{sec:toron-regularity}
The projector in \cref{eq:tree-gauge-map} cannot be treated as a finite
Laurent polynomial. Write $T_0$ for tree zeroing before the final
longitudinal subtraction, and let $R_{\partial}:\C^4\to\C^{64}$ place
one value on the eight boundary edges of its family. Thus
$T_0R_{\partial}=R_{\partial}$ and $LR_{\partial}=I$. At nonzero
momentum,
\[
 T_{\rm g}=(I-R_{\partial}P_{\rm long}L)T_0,\qquad
 P_{\rm long}=qq^*/\lambda,\qquad
 T_{\rm g}R_{\partial}=R_{\partial}(I-P_{\rm long}).
\]

\begin{proposition}[Punctured-torus regularity and integrable jets]
\label{prop:toron-regularity}
The implemented representative is bounded and smooth away from zero,
but has no continuous extension at zero on the full edge space. For
$j=0,1,2$ it satisfies
\[
 \norm{D^jT_{\rm g}(k)}\le C_j|k|^{-j}
\]
near zero. With any value assigned at zero it belongs to
$W^{1,p}(\T^4)$ for $1\le p<4$ and to $W^{2,p}(\T^4)$ for
$1\le p<2$. In particular $DT_{\rm g}\in L^2$ and
$D^2T_{\rm g}\in L^1$.
\end{proposition}
\begin{proof}
Along $k=te_1$ the map sends $R_{\partial}e_1$ to zero; along
$k=te_2$ it fixes that nonzero vector. The limits disagree. Boundedness
follows from $\norm{P_{\rm long}}=1$ and the finite Laurent matrices
$R_{\partial},L,T_0$. Since $\lambda\asymp|k|^2$, differentiating
$qq^*/\lambda$ gives the displayed bounds. Radial integration in four
dimensions gives $jp<4$. Integration by parts outside a radius-$r$
ball has boundary errors $O(r^3)$ for the first derivative and
$O(r^2)$ for the second; hence these are also the weak derivatives.
\src{F11}{V42}
\end{proof}
Finite products of translated such representatives and compatible
jointly $C^2$ bounded matrix factors are twice differentiable into
$L^1$. A second derivative contains either one $L^1$ second derivative
or two $L^2$ first derivatives; all other factors are bounded.
Consequently their second derivatives may be integrated. If a product
trace has bound $|\partial_t^2F(k,0)|\le B_2|k|^{-2}$ near zero, its
omitted normalized integral over $|k|<\delta$ is at most
$\pi^2B_2\delta^2/(2\pi)^4$. These finite-product estimates do not by
themselves bound a growing all-level alias trace. They refine, rather
than eliminate, the toron requirement in \cref{crit:background-ADMC}.

\section{Fixed and harmonic background prescriptions}
\label{sec:background-prescription}
The natural background of an exactly blocked tree-level action is the
harmonic extension
\begin{equation}\label{eq:harmonic-background}
b_j^{\rm harm}=\mathcal M^{(j+1)}\beta
\end{equation}
of the level-$(j+1)$ coarse polarization. This remains the background
associated with the minimizing construction. It cannot be inserted into
zero-background-momentum vertex banks as though it were a single fine
plane wave. Its nonconstant aliases have fine momenta
$2^{-j-1}(q+2\pi M)$ with $M\ne0$ and amplitudes of order $q$; they can
therefore contribute to the second external-momentum derivative.

A second, explicitly fixed prescription uses the fine tree-gauge router
extension of the reference coarse polarization plane wave in direction
$2$, with the same external-momentum jets at zero at every level. Denote
this background by $b^{\rm fix}$. Its vertex banks are fixed, while the
scale dependence enters through \cref{eq:nested-tree-cov}. The two
prescriptions must be kept distinct. Their limiting $F^2$ coefficients
are not identified by the gauge-representative map, which acts on
fluctuation legs rather than replacing a background variation.
\src{F11}{V41}

With this fixed prescription the tree-level Wilson shell density is
\begin{equation}\label{eq:shell-functional}
\begin{aligned}
\Omega_j^{\rm fix}(k,q)
={}&\frac12\Tr\bigl[\Xi_j(k)Q_0(b^{\rm fix};q)\bigr]\\
&-\frac13\Tr\bigl[\Xi_j(k+qe_1)R_0(b^{\rm fix};q)
\Xi_j(k)R_0(b^{\rm fix};q)^*\bigr].
\end{aligned}
\end{equation}
The trace runs over the full corresponding fine alias class, in the same
normalization as the banks. The downward coefficient is
\begin{equation}\label{eq:b-shell}
b_{G,j}^{\mathrm{fix},\downarrow}
=\frac32\int_{\T^4}
\left.\partial_q^2\Omega_j^{\rm fix}(k,q)\right|_{q=0}\dd k.
\end{equation}
Both the tadpole and bubble are differentiated completely, including the
momentum dependence of the covariance, its representative map, and the
vertices. They are assembled before absolute estimates. A coefficient for
the harmonic prescription instead requires the full background jets of
\cref{eq:harmonic-background}, including its nonzero aliases. No equality
between these two responses is imposed by notation.

\begin{scope}[What the fixed prescription supplies]
The fixed background makes the existing fine-edge banks usable in an
explicit shell-response prescription at every level. It does not itself
prove that this prescription equals the derivative of the exact nested
blocked action in its harmonic background. Such an identification requires
a background-compatibility estimate, or a direct calculation in the
harmonic prescription. It is also distinct from the complete Haar,
quotient, and determinant normalization required by ADMC.
\end{scope}

\section{Exact finite cards in the declared convention}
\label{sec:fixed-background-cards}
The three entries below are the corner means of the order-zero,
order-one, and order-two true external-momentum derivatives of the
level-zero Wilson response. For the block-average rule with the fixed
background, the reported exact values are
\begin{equation}\label{eq:corners-fixed}
\overline{\mathfrak c}^{\,\rm fix}_{\rm block}
=\left(\frac{5587}{1568},\ 0,
-\frac{407872449589}{47416320000}\right).
\end{equation}
For the same block rule with its own harmonic background, they are
\begin{equation}\label{eq:corners}
\overline{\mathfrak c}^{\,\rm harm}_{\rm block}
=\left(\frac{5587}{1568},\ 0,
-\frac{254367096889}{47416320000}\right),
\end{equation}
whereas the straight-rule reference gives
\begin{equation}\label{eq:corners-straight}
\overline{\mathfrak c}_{\rm straight}
=\left(\frac{176943}{627200},\ 0,
-\frac{2822949104411}{265531392000}\right).
\end{equation}
The first two block-rule entries agree, while their order-two difference
is exactly
\begin{equation}\label{eq:background-corner-difference}
\overline{\mathfrak c}^{\,\rm fix}_{{\rm block},2}
-\overline{\mathfrak c}^{\,\rm harm}_{{\rm block},2}
=-\frac{3480847}{1075200}.
\end{equation}
These identities are reported finite certificates from
\src{F11}{V39, V41}. Their difference shows why the background convention
must accompany any quoted second-order card. It does not prove that the
two full-torus integrals or their limiting tower coefficients are unequal.
Conversely, calling the difference a lattice-scale effect is not a proof
that its accumulated contribution vanishes. The cards are level-zero
quantities, not higher-shell coefficients, not exact torus integrals, and
not an evaluation of $b^*$.

\section{The fixed-point response and background-compatibility criteria}
\label{sec:background-matching}
\begin{criterion}[Gaussian matching from a nested fixed point]
\label{thm:fixedpoint-anomaly}
Suppose the nested kernels converge at rate $C\rho^j$, $\rho<1$, and the
complete shell response in a declared prescription has a common Lipschitz
control along the tower. That control includes the vertex banks,
representative maps, external derivatives, infrared treatment, trace
normalization, and all Gaussian-principal contributions needed by the
chosen determinant. If the consistently oriented coefficients satisfy
\[
b_{G,j}=b^*+O(\rho^j),
\]
then
\begin{equation}\label{eq:fixedpoint-ADMC}
\sum_{j<N}b_{G,j}=Nb^*+O(1).
\end{equation}
If these are the coefficients of the actual nested determinant and
$b^*=b_{\rm YM}$ in the same convention, its Gaussian ADMC follows.
Exact nesting creates no reference-reset term.
\end{criterion}
\begin{proof}
Sum the geometric coefficient error. For coefficients of the actual
nested determinant, the normalization telescopes as in
\cref{eq:det-telescope}. The physical coefficient identification is an
independent input. This is the nested fixed-point implication of
\src{F11}{V34, V38}, with the representative and background requirements
made explicit.
\end{proof}

The $4^{-j}$ covariance and inverse-kernel estimates establish the kernel
part of this criterion. Multilinearity of \cref{eq:shell-functional}
gives Lipschitz bounds on any fixed finite alias class. It does not remove
the need for a common bound as the class grows, or justify differentiating
through the singular toron representation without an appropriate
alternative chart or integrable control.

For a fixed-background evaluation, the missing connection to the actual
nested determinant can be stated at precisely the averaged strength
required by ADMC. Let $b_{G,j}^{\rm fix}$ and $b_{G,j}^{\rm harm}$ use the
same upward orientation, physical normalization, and complete principal
contributions, and define
\begin{equation}\label{eq:background-defect}
\delta_j^{\rm bg}=b_{G,j}^{\rm harm}-b_{G,j}^{\rm fix}.
\end{equation}

\begin{criterion}[Background compatibility at the ADMC interface]
\label{crit:background-ADMC}
Assume the fixed-background coefficients obey
$b_{G,j}^{\rm fix}=b_{\rm fix}^*+O(\rho^j)$ and
\begin{equation}\label{eq:background-Cesaro}
\frac1N\sum_{j<N}\delta_j^{\rm bg}\longrightarrow0.
\end{equation}
Then
\[
\sum_{j<N}b_{G,j}^{\rm harm}=Nb_{\rm fix}^*+o(N).
\]
In particular, $b_{\rm fix}^*=b_{\rm YM}$ supplies the Gaussian
coefficient part of ADMC for the harmonic determinant, provided its other
normalization and endpoint requirements hold. A sufficient form of
\cref{eq:background-Cesaro} is
\begin{equation}\label{eq:background-coboundary}
\delta_j^{\rm bg}=d_{j+1}-d_j+r_j,\qquad
d_N=o(N),\qquad \frac1N\sum_{j<N}r_j\longrightarrow0.
\end{equation}
\end{criterion}
\begin{proof}
Sum \cref{eq:background-defect} and apply
\cref{eq:fixedpoint-ADMC,eq:background-Cesaro}. For the sufficient
condition, the $d_j$ terms telescope to $d_N-d_0$, which is sublinear.
This is the SCIC telescoping argument of \cref{def:SCIC} applied to the
specific background discrepancy; it does not establish the discrepancy
estimate itself.
\end{proof}

The criterion expresses, rather than resolves, the background-independence
assumption. It permits either direct computation with the full harmonic
background or a fixed-background computation accompanied by a tested
comparison. The level-zero difference in
\cref{eq:background-corner-difference} alone decides neither route.

\section{Physical normalization and the unresolved scalar}
The colour and action convention is
\[
-\tr_{\rm ad}(\operatorname{ad}X\operatorname{ad}Y)=C_A\ip{X}{Y},
\qquad S_{\rm bg}=\frac1{4u}\int|\overline F|^2.
\]
Its upward dyadic comparison coefficient is
\begin{equation}\label{eq:bYM}
b_{\rm YM}=\frac{22C_A}{3(4\pi)^2}\log2.
\end{equation}
The one-colour incidence Hessian carries a factor $1/3$ relative to the
colour-traced convention, hence the propagator factor $3$ in
\cref{eq:response-alias-cov}. The reference flat Schur form has
$s_0(q)=q^2/3+O(q^4)$, so a response
$\gamma(q)=\kappa_{\rm tree}q^2+O(q^4)$ shifts the downward inverse
coupling by $3\kappa_{\rm tree}$. This accounts for the factor $3/2$
multiplying the second true derivative in \cref{eq:b-shell}. The direction
of the iteration must then be converted to the upward convention of
\cref{eq:bYM}; a downward card is not compared to that coefficient without
this conversion.

The one-step nonlinear response of \cref{sec:actual-theta} supplies an
additional actual-fibre integral, with its minimizing lift, fourth jets,
and Haar terms. Its $\SU(2)$ coefficient $\theta$ is not by definition
the complete all-level colour-normalized coefficient in
\cref{eq:bYM}. The flat cancellation of \cref{thm:flat-bloch} validates
one normalization interface, not that physical identification.

The finite representative correction and the frame factors are now
explicit. The outstanding calculation is nevertheless a full-torus,
complete-principal response along the nested tower, with control of the
toron region, the growing alias trace, the background comparison where
needed, and the continuum normalization. No corner mean is identified
with \cref{eq:bYM}, and no principal-margin protection theorem supplies
that identification. These distinctions leave the Gaussian anomaly,
the interacting continuation, and the physical infrared contraction as
separate requirements.
\sourcebasis{The explicit propagator, inverse bounds, extensions, and
positive shells are established in \src{F11}{V37--V38}; the full-edge
banks and finite response cards in \src{F11}{V39}; the representative map
and fixed-background prescription in \src{F11}{V41}. The averaged
background criterion is the elementary specialization of
\cref{def:SCIC} displayed above.}

\chapter{Nonlinear compact fibres and the measured local response}
\label{ch:compact-response}
The Gaussian tower of \cref{ch:tower} fixes a principal quadratic
object, but not a nonlinear nonabelian fibre or its measured response.
This chapter supplies the corresponding one-step construction. The
root-path block, its solved pivots, the moving constrained minimum, and
the Haar density are kept together. The resulting local coefficient has
a thermodynamic limit and a volume-uniform one-loop remainder on its
specified box. Neither statement identifies that box with the complete
compact fibre or supplies an iterated physical trajectory.

\section{A nonlinear block with the prescribed linear germ}
\label{sec:nonlinear-root-block}
Let $G$ be a compact matrix gauge group. Divide the fine lattice into
sixteen-site parity cells and use the rooted tree which clears the
highest active parity bit. Tree retraction leaves fifteen tree links
fixed to identity in each cell. Of the remaining forty-nine group
coordinates, seventeen are internal non-tree links and thirty-two are
boundary links. In each of the four coarse directions choose one root
boundary link as a pivot. The forty-five other groups are independent
fibre coordinates.

For a coarse edge $e$, the sixteen chronological two-link paths have two
identical root products $B_e$ on this tree slice. Write
$P_{e,1},\ldots,P_{e,14}$ for the other products. None contains a pivot;
each depends only on the nonpivot boundary groups of this edge and the
internal groups of its two endpoint cells. On the common logarithm
domain the block is
\begin{equation}\label{eq:root-block-local}
K_e(Y,B)=B_e\exp\left\{\frac1{16}\sum_{j=1}^{14}
                  \log(B_e^{-1}P_{e,j}(Y))\right\}.
\end{equation}
Its differential at the identity is the sixteen-path linear average
$L$. A right inverse $R$ places the same coarse Lie-algebra value in all
eight boundary links of a family and zero in its internal links; then
$LR=I$. A free-coordinate map $E$ parametrizes $\Ker L$. In particular,
the four pivot rows of $E$ are determined by $LE=0$, not selected as
additional independent directions. \src{F11}{V57--V58, V64--V66, V76}

\begin{proposition}[Why the linear average is not a nonlinear gauge map]
\label{prop:nonabelian-average}
For a nonabelian Lie algebra $\mathfrak g$, a scalar-weight map
$A:\mathfrak g^{\oplus m}\to\mathfrak g$ of the form
$A(X_1,\ldots,X_m)=\sum_i w_iX_i$ can be a Lie algebra homomorphism
only if $w_iw_j=0$ for $i\ne j$ and $w_i^2=w_i$ for every $i$.
Consequently a genuine multi-entry average is not the differential of
a Lie-group homomorphism with these domain and range groups.
\end{proposition}
\begin{proof}
Elements supported in different summands commute in the direct-sum
algebra. Their images have bracket $w_iw_j[X,Y]$, which must vanish
for every $X,Y$. Within one summand, bracket preservation instead
requires $w_i^2[X,Y]=w_i[X,Y]$. Nonabelianity gives both assertions.
The derivative of a Lie-group homomorphism preserves the bracket.
\end{proof}
This obstruction concerns a proposed homomorphism, not gauge covariance
of \cref{eq:root-block-local}. Endpoint path dressing makes that block
covariant under the residual cell-constant gauge transformations. The
field dependence of this dressing is precisely what a frozen linear
average omits. \src{F11}{V56--V57}

\subsection{The actual constrained Hessian and density}
Near the identity write the local fibre as
\[
x=X(z,c)=Ez+Rv(z,c),\qquad F(X(z,c))=c,\qquad F=\log K.
\]
The implicit construction has a common local maximum-norm domain, and
its fixed-order local derivatives are uniform in volume. Put
$f(z,c)=S_0(\exp X(z,c))$. At a constrained stationary point, with
$\dd S_0=\lambda\,\dd F$, the pullback Hessian is
\begin{equation}\label{eq:actual-constrained-Hessian}
D_z^2 f=(D_zX)^*\bigl(D^2S_0-\lambda\cdot D^2F\bigr)D_zX.
\end{equation}
Indeed the ordinary chain rule contains the additional term
$\dd S_0[D_z^2X]$, while twice differentiating $F(X)=c$ gives
$DF[D_z^2X]=-D^2F[D_zX,D_zX]$. Thus neither the moving tangent nor its
acceleration can be dropped from an off-shell response.

With a fixed ordering of the free coordinates, the local measure
relative to $\dd z\,\dd c$ is, up to a constant independent of $z,c$,
\begin{equation}\label{eq:root-local-density}
 e^{-\beta f(z,c)+\ell(z,c)}\dd z\,\dd c,\qquad
 \ell=\sum_{\text{non-tree }l}\log j_G(X_l)
             -\log\det\bigl(DF(X)R\bigr).
\end{equation}
The determinant is positive on the chosen local branch. The pivot
structure makes its relevant coarse blocks finite and local. Its
second derivative involves the third jet of $F$; the quadratic action
alone does not determine the measured order-one response. This is the
concrete fibre realization of the Jacobian-jet requirement in
\cref{prop:Jacobian-third-jet}. \src{F11}{V55, V58--V60}

\section{A smooth global compact completion}
\label{sec:global-pivots}
Fix a smooth nonincreasing cutoff $\chi:[0,\infty)\to[0,1]$, equal to
one on $[0,1/2]$ and zero on $[1,\infty)$. For a sufficiently small
fixed $\rho>0$ inside an injective logarithm neighbourhood, define
\[
 g_\rho(U)=\chi(\norm{\log U}/\rho)\log U
\]
there and extend it by zero. This is a smooth conjugation-equivariant
function on $G$. Define the completed block $\widetilde K$ by replacing
every logarithm in \cref{eq:root-block-local} with $g_\rho$.

\begin{lemma}[An upper derivative bound for the truncated logarithm]
\label{lem:pivot-cutoff}
For a fixed cutoff and sufficiently small $\rho$, constants independent
of the group input satisfy
\[
 \norm{g_\rho}\le\rho,\qquad
 \norm{D_Lg_\rho}\le C_\chi,\qquad
 \operatorname{Sym}D_Lg_\rho\preceq(1+C\rho)I.
\]
Here $D_Lg(U)v=\left.\frac{\dd}{\dd t}g(e^{tv}U)\right|_{t=0}$.
\end{lemma}
\begin{proof}
In logarithm coordinates $h(z)=\chi(\norm z/\rho)z$ has derivative
\[
 Dh(z)=\chi(s)I+s\chi'(s)\frac{z\otimes z}{\norm z^2},
 \qquad s=\norm z/\rho.
\]
Its eigenvalues are at most one, and their absolute values are bounded
by a constant depending only on $\chi$. The left logarithmic derivative
is $I+O(\norm z)$ uniformly on the support. Composing these two
derivatives proves the estimates; the zero extension is smooth.
No positive lower bound for the cutoff derivative is asserted.
\end{proof}

\begin{theorem}[Globally invertible compact pivots]
\label{thm:global-pivot}
For sufficiently small fixed $\rho$, arbitrary compact inputs
$P_1,\ldots,P_{14}$, and $0\le t\le1$, the maps
\[
 F_t(B)=B\exp(tX(B)),\qquad
 X(B)=\frac1{16}\sum_{j=1}^{14}g_\rho(B^{-1}P_j)
\]
are smooth diffeomorphisms of $G$. All differential singular values are
bounded above and bounded below by $1/16$, after decreasing $\rho$.
Their inverses have uniform derivatives of every fixed order in the
compact source parameters.
\end{theorem}
\begin{proof}
For the input variation $Be^{sv}$, put
$A_B=16^{-1}\sum_jD_Lg_\rho(B^{-1}P_j)$. Up to an orthogonal adjoint
factor the derivative of $F_t$ is
\[
 M_t=I-t\mathcal J(tX)A_B,\qquad
 \mathcal J(Y)=\int_0^1e^{s\operatorname{ad}Y}\dd s.
\]
The preceding lemma gives $\norm X\le14\rho/16$,
$\norm{A_B}\le14C_\chi/16$, and
$\operatorname{Sym}A_B\preceq(14/16)(1+C\rho)I$.
Since $\mathcal J(tX)=I+O(\rho)$,
\[
 \operatorname{Sym}M_t\succeq(1-14t/16-C'\rho)I
                         \succeq\tfrac1{16}I
\]
for sufficiently small $\rho$. This bounds the inverse differential.
The two duplicate root paths are essential to the strict margin.

To pass from local to global invertibility, solve the time-dependent
vector field $V_t=-(DF_t)^{-1}\partial_tF_t$ on the compact group.
Its complete flow $\Psi_t$ satisfies
$\frac{\dd}{\dd t}F_t(\Psi_t(B))=0$, hence
$F_t\circ\Psi_t=I$. The flow is a diffeomorphism, so $F_t=\Psi_t^{-1}$.
Smooth parameter dependence and repeated implicit differentiation on
the fixed compact parameter carrier give all stated derivative bounds.
They may depend on $\rho$, which is not sent to zero in this theorem.
\end{proof}

\begin{corollary}[Exact global product carrier and smooth density]
\label{cor:compact-product}
For $n$ coarse cells, the map
\[
 (Y,B)\longmapsto(Y,\widetilde K(Y,B)),\qquad
 Y\in G^{45n},\quad B\in G^{4n},
\]
is a global diffeomorphism. Its inverse $B=B(Y,C)$ has
\begin{equation}\label{eq:compact-haar-density}
 \dd Y\,\dd B=j(Y,C)\dd Y\,\dd C,\qquad
 j=\prod_e j_e,\qquad 0<c\le j_e\le C_1<\infty.
\end{equation}
Each factor and each solved pivot has a fixed finite stencil and
uniform fixed-order derivative bounds. For every finite volume and
finite $\beta$, the completed blocked law has a positive smooth density
relative to coarse Haar,
\begin{equation}\label{eq:completed-Z}
 \widetilde Z_\beta(C)=\int_{G^{45n}}
       e^{-\beta S_0(Y,B(Y,C))}j(Y,C)\dd Y.
\end{equation}
It is gauge invariant. If
$V_\beta=\beta S_0(Y,B(Y,C))-\log j(Y,C)$ and
$\widetilde W=-\log\widetilde Z_\beta$, then in any fixed coarse chart
\begin{equation}\label{eq:compact-score}
 \partial_a\widetilde W=\E V_{\beta,a},\qquad
 \partial_a\partial_b\widetilde W
   =\E V_{\beta,ab}-\Cov(V_{\beta,a},V_{\beta,b}).
\end{equation}
\end{corollary}
\begin{proof}
At fixed $Y$, each output depends only on its own pivot. Apply
\cref{thm:global-pivot} to each edge. The total Jacobian is triangular,
so its inverse Haar factor is the product in
\cref{eq:compact-haar-density}. Compactness permits differentiation of
\cref{eq:completed-Z}. Endpoint gauge transformations conjugate each
relative source, preserve Haar measure and the Wilson action, and
transport the inverse pivots; this proves invariance. Differentiating
the normalized integral gives \cref{eq:compact-score}.
\end{proof}
The per-factor bounds in \cref{eq:compact-haar-density} are not a
volume-independent bound on the complete product $j$. The local score
coefficients have size $O(\beta+1)$, but this alone gives no uniform
spatial decay of their covariance. Also, \cref{eq:completed-Z} is
already a density relative to coarse Haar: no second coarse-Haar
factor is inserted in \cref{eq:compact-score}. \src{F12}{V18}

\begin{scope}[Agreement of germs and identification of laws]
\label{scope:completion-germ}
The completed block agrees with the original logarithmic block wherever
all relative logarithms lie inside the inner cutoff region. Their local
action and measured jets therefore agree on a possibly smaller common
box. A piecewise compact fallback and this smooth completion need not
have the same global pushforward. Equality of their germs does not
identify their normalized compact conditionals, their reflected laws,
or their iterates. A change of completion requires its own comparison
in \cref{hyp:compact}; it is not absorbed into Gaussian notation.
\end{scope}

\subsection{The metric and completion equation in the explicit SU(2) fibre}
\label{sec:su2-pivot-equation}
The explicit compact spectral application uses $G=\SU(2)$ with its
invariant Hilbert--Schmidt Riemannian metric. Write $d$ for its geodesic
distance and $d_{\rm ch}(U,V)=\norm{U-V}_{\rm HS}$ for chordal distance.
On the logarithm support $d(I,U)=\norm{\log U}_{\rm HS}$, and
$d_{\rm ch}\le d$. The working cutoff is $\rho=1/64$ in the geodesic
convention. Chordal lower margins used below are sufficient inactivity
conditions; they do not redefine the cutoff.

There is also a direct characterization at this specified radius. Set
\[
 \Phi_\rho(r)=\int_0^r t\chi(t/\rho)\dd t,\qquad
 R_\rho=7\rho/8,\qquad
 \kappa_\rho=\tfrac18-\tfrac{49\rho^2}{256}.
\]
For arbitrary fourteen source products define, on the geodesic ball
$\overline B_{R_\rho}(C)$,
\[
 \mathcal V_C(B)=\tfrac12d(B,C)^2
             -\tfrac1{16}\sum_j\Phi_\rho(d(B,P_j)).
\]

\begin{proposition}[Strongly convex characterization of the actual pivot]
\label{prop:pivot-potential}
For $0<\rho\le1/4$,
$\operatorname{Hess}\mathcal V_C\succeq\kappa_\rho I>0$ on this
ball. Its unique critical point $B_*$ is the unique solution of the
completed coarse equation. In right-trivialized coordinates its
residual and a certified error estimate are
\begin{align}
 \mathcal R_C(B)&=\log(C^{-1}B)
             +\tfrac1{16}\sum_j g_\rho(B^{-1}P_j),
             \label{eq:pivot-residual}\\
 d(B,B_*)&\le\kappa_\rho^{-1}\norm{\mathcal R_C(B)}.
             \label{eq:pivot-residual-bound}
\end{align}
In particular, $B_*=C$ if and only if
$\sum_jg_\rho(C^{-1}P_j)=0$, and
$\kappa_{1/64}=131023/1048576$.
\end{proposition}
\begin{proof}
The group is the round sphere of radius $\sqrt2$. A radial potential
$f(r)$ has radial Hessian $f''(r)$ and transverse Hessian
$f'(r)\cot(r/\sqrt2)/\sqrt2$. For $r^2/2$ the latter is at least
$1-r^2/4$ on this ball. For $\Phi_\rho$, the radial eigenvalue is
$\chi(s)+s\chi'(s)\le1$ and the transverse eigenvalue is at most one;
outside its support both vanish. Subtracting fourteen such upper
bounds gives $1-R_\rho^2/4-14/16=\kappa_\rho$.
The radial derivative at the boundary is strictly positive because
$\norm{g_\rho}<\rho$. Hence the minimum is interior, and strict
geodesic convexity gives uniqueness. Its gradient is
\cref{eq:pivot-residual}. Every solution of the completed equation
lies within the same ball, so the characterization is global for the
pivot equation. Integrating the Hessian lower bound along the geodesic
from $B_*$ to $B$ and applying Cauchy--Schwarz proves
\cref{eq:pivot-residual-bound}.
\end{proof}
This convexity belongs to an auxiliary pivot-solving potential, not to
the Wilson action. Small independent pivot caps are not replacements
for \cref{eq:pivot-residual}: actual source balance, including possible
cancellation of active sources, is required. \src{F12}{V98}

\section{The moving minimum and its measured coefficient}
\label{sec:actual-one-loop}
Return to the local germ of \cref{sec:nonlinear-root-block} and to
orthonormal Lie coordinates. The free fibre Hessian obeys the
volume-uniform eliminated-space bound
\[
 f_{zz}(0,0)=E^*d^*dE\succeq hI,\qquad h=24/2327.
\]
The finite incidence construction and its coercivity estimate concern
this fibre, not the retained gauge field. On a common smaller
maximum-norm domain there is a unique nearby stationary section
$z_*(c)$, with
\begin{equation}\label{eq:moving-minimum}
 f_z(z_*,c)=0,\qquad H_c=f_{zz}(z_*,c)\succeq h_1I,\qquad
 D z_*=-H_c^{-1}f_{zc},
\end{equation}
where $h_1>0$ is uniform in volume. The inverse Hessian and the
fixed-order local derivatives of this section have exponential spatial
locality. Defining $G(c)=f(z_*(c),c)$ gives
\begin{equation}\label{eq:minimized-Schur}
 D^2G=f_{cc}-f_{cz}H_c^{-1}f_{zc}.
\end{equation}
This is the nonlinear constrained counterpart of the Schur identity
in \cref{prop:schur}; its retained gauge null directions at zero have
not been removed by the fibre bound. \src{F11}{V66, V80--V81}

Let $\mathcal D$ be the fixed local integration box and
$N=45(\dim G)n$. With the actual density \cref{eq:root-local-density},
put
\begin{align}
 Z_\beta^{\rm loc}(c)&=\int_{\mathcal D}
                e^{-\beta f(z,c)+\ell(z,c)}\dd z,
     &W_\beta^{\rm loc}&=-\log Z_\beta^{\rm loc},\notag\\
 \widehat W_\beta^{\rm loc}&=W_\beta^{\rm loc}+\log j_c(c),
     &\widehat Q(c)&=\tfrac12\log\det H_c-\ell(z_*(c),c)
                                     +\log j_c(c).
 \label{eq:normalized-local-Q}
\end{align}
Here $j_c(c)=\prod_ej_G(c_e)$ converts the coordinate density to a
density relative to coarse Haar. This conversion is needed in
\cref{eq:normalized-local-Q}, unlike \cref{eq:completed-Z}.

\begin{theorem}[Volume-uniform measured one-loop approximation on the box]
\label{thm:uniform-local-loop}
On the declared common local domain, let
\[
 \mathcal R_\beta=\widehat W_\beta^{\rm loc}-\beta G
                  -\tfrac N2\log(\beta/2\pi)-\widehat Q.
\]
For sufficiently large $\beta$, constants independent of finite volume
and of the admitted small coarse field satisfy
\begin{align}
 N^{-1}|\mathcal R_\beta|+|\partial_a\mathcal R_\beta|
      &\le C\beta^{-1/2},\label{eq:loop-error-zero}\\
 |\partial_a\partial_b\mathcal R_\beta|
      &\le C\beta^{-1/2}e^{-\gamma d(a,b)},\label{eq:loop-error-local}\\
 \norm{D_c^2\mathcal R_\beta}_{\rm row}
       +\norm{D_c^2\mathcal R_\beta}_{\op}
      &\le C\beta^{-1/2}.\label{eq:loop-error-row}
\end{align}
The indices $a,b$ are fixed local coarse directions, and the row norm
sums the norms of their spatial Hessian blocks.
\end{theorem}
\begin{proof}
Use the centered convex extension on the common local domain. Its
localized expansion, with the actual minimum, Hessian, and density,
has an $O(\beta^{-1/2})$ error per degree of freedom and the localized
first and second derivative bounds of \cref{eq:loop-error-local}.
The comparison with the original potential on the same box is
$O(e^{-\kappa\beta})$ in those norms. These are the local extension
and same-box estimates of \src{F12}{V10--V11}.

For completeness, the nontrivial boundary comparison is a normalized
wall ratio, not the removal of a rare global union. Introduce a
nonnegative coarse-independent wall $\Psi$ and the potential
$E_\beta+\lambda\Psi$, $0\le\lambda<\infty$. For its free energy
$F_\lambda$, normalized differentiation gives
\begin{align*}
 \partial_\lambda F_\lambda&=\E_\lambda\Psi,\\
 \partial_\lambda\partial_aF_\lambda
              &=-\Cov_\lambda(E_{\beta,a},\Psi),\\
 \partial_\lambda\partial_a\partial_bF_\lambda
              &=-\Cov_\lambda(E_{\beta,ab},\Psi)
                +\kappa_{3,\lambda}(E_{\beta,a},E_{\beta,b},\Psi).
\end{align*}
The localized wall estimates give, after summing the local wall pieces,
\[
 |\partial_\lambda\partial_a\partial_bF_\lambda|
 \le C\beta^2e^{-a_0\beta}(1+\lambda)^{-9/8}
                                      e^{-\gamma d(a,b)},
\]
with the analogous first-derivative and per-site zeroth-order bounds.
Since the strength integral is $8$, the limiting hard-wall ratio and
its first two derivatives are exponentially small. At finite volume,
dominated convergence identifies the limit with the actual box integral;
the wall has no omitted coarse derivative. Adding the extension error,
same-box comparison, and wall ratio proves
\cref{eq:loop-error-zero,eq:loop-error-local}. Summability of the
spatial exponential gives \cref{eq:loop-error-row}.
\src{F12}{V12}
\end{proof}
This theorem strengthens a fixed-volume Laplace expansion. It is not
an all-orders volume-uniform expansion, and it gives no derivative
estimate for the logarithm of the full compact complement.

\subsection{A complete fourth-jet formula}
At the identity set $P=f_{zz}(0)^{-1}$, $A=f_{zc}(0)$ and
$Ju=(-PAu,u)$. Let $f_3,f_4$ be the third and fourth derivatives of
the actual pullback action and define fibre matrices
\[
 B(\xi)_{ij}=f_3[e_i^z,e_j^z,\xi],\qquad
 C(\xi,\eta)_{ij}=f_4[e_i^z,e_j^z,\xi,\eta].
\]
For the semisimple gauge groups and the normalized convention under
consideration, invariance makes the relevant first derivative at zero
vanish. The complete measured Hessian is
\begin{equation}\label{eq:complete-fourth-jet}
\begin{split}
 D^2\widehat Q(0)[u,v]={}&\tfrac12\Tr\{P C(Ju,Jv)
                       -P B(Ju)P B(Jv)\}\\
 &-\ell_2[Ju,Jv]
       +\tfrac1{12}\sum_e\Tr_{\mathfrak g}
                       (\operatorname{ad}u_e\operatorname{ad}v_e).
\end{split}
\end{equation}
Indeed differentiating $\tfrac12\log\det H_c$ gives the two trace
terms; differentiating the density gives the last line. Invariance
removes the term in the second derivative of the minimizing section,
not its first derivative $J$. The formula retains the block constraint,
pivot acceleration, and Haar normalization through the actual jets.
\src{F11}{V91--V93}

Exponential locality of these kernels gives an exponentially local
infinite-volume coefficient kernel and a row-summable finite-volume
periodization error. Its Fourier symbol $\mathcal K(p)$ is analytic
and obeys the Ward identity
\[
 \mathcal K(p)g(p)=0,\qquad g_\mu(p)=1-e^{\ii p_\mu},\qquad
 \mathcal K(0)=0.
\]
For the last assertion differentiate the Ward identity once at zero;
the derivatives of $g$ span the four coarse directions. This is a
thermodynamic limit of a coefficient, not of the entire quantum law.

As an exact finite-regulator normalization test, the reported
$\SU(2)$ checkerboard calculation uses
$c_{r,1}=(-1)^{r_0}T$, $\norm T_{\rm HS}^2=2$, on sixteen coarse cells.
The contributions to $n^{-1}D^2\widehat Q(0)[c,c]$ are
\[
 \frac{19}{3}-\frac{689}{384}
          +\frac{79227}{4480}-\frac{166901}{5760}
                      =-\frac{272249}{40320}.
\]
These are respectively the normalized measure contribution, neutral
term, charged second-vertex term, and subtracted first-vertex square.
The equality is a reported finite certificate in \src{F11}{V94--V99}.
The external momentum is $\pi$; its negative value is not a mass or an
infrared coefficient.

\section{The actual small-momentum coefficient}
\label{sec:actual-theta}
The remaining one-step transverse coefficient can be formulated without
a small real-momentum difference quotient. For the $45$-component
scalar-colour fibre write
\[
 H(k)=E^\#(k)D^\#(k)D(k)E(k),\qquad P(k)=H(k)^{-1}.
\]
Here $A^\#(k)=A(-k)^T$ is the Laurent continuation of the real-torus
adjoint. It is not complex conjugation of the complex momentum.
The finite coefficient enumeration reports $33$ nonzero Laurent
coefficients $H_r$, with
\[
 R_H=\max|r|_1=2,\qquad
 W_H=\sum_r\max(\norm{H_r}_1,\norm{H_r}_\infty)=1041/2.
\]
Together with $H(k)\succeq hI$, $h=24/2327$, these data give
\begin{equation}\label{eq:actual-holomorphic-strip}
 s=\frac{h}{12W_HR_H}=\frac2{2422407},\qquad
 \norm{P(k)}\le2/h\quad(\norm{\operatorname{Im}k}_\infty\le3s).
\end{equation}
To see this, bound the perturbation off the real torus by
$W_H(e^{3sR_H}-1)\le h/2$ and use a Neumann series.
The integer coefficient data are reported in \src{F12}{V5}; the
analytic implication is independent of momentum sampling.

Fix a unit real external polarization $v$ and external momentum $pe_0$.
Let $V(k,p)$ be the actual cubic vertex along $J(p)v$, carrying fibre
momentum $k$ to $k+p$, and let $C(k,p)$ be the zero-transfer fourth
vertex for the pair of opposite external modes. Define
$\widetilde V(k,p)=V(k+p,-p)$. These are the vertices of
\cref{eq:complete-fourth-jet}, not frozen fine-action vertices.

\begin{theorem}[Differentiated determinant representation]
\label{thm:theta-integral}
In the stated unit-polarization normalization,
\begin{align}
 g(k,p)&=\tfrac12\Tr\{P(k)C(k,p)
       -P(k+p)V(k,p)P(k)\widetilde V(k,p)\},\notag\\
 d_2(k)&=[p^2]g(k,p),\qquad
 \theta=\frac{19}{48}
           +\int_{\T^4}d_2(k)\frac{\dd k}{(2\pi)^4}.
 \label{eq:actual-theta}
\end{align}
Coefficient extraction commutes with integration. If $M$ is an upper
bound for $|g|$ on the complex domain
$\norm{\operatorname{Im}k}_\infty\le s$, $|p|\le s$, then the
four-dimensional tensor-grid mean $\theta_N$ satisfies
\begin{equation}\label{eq:theta-quadrature}
 |\theta_N-\theta|\le\frac{M}{s^2}
 \left\{\left(\frac{1+e^{-sN}}{1-e^{-sN}}\right)^4-1\right\}.
\end{equation}
\end{theorem}
\begin{proof}
Apply \cref{eq:complete-fourth-jet} to opposite commensurate plane waves
on finite tori. The fourth vertex has zero transfer, while the cubic
pair shifts momentum by $p$ and back. The trace per cell is therefore
the discrete mean of $g$. Locality and the inverse bound give the
thermodynamic integral; the normalized Haar term has quadratic
coefficient $19/48$ in this convention. The finite-range jets and
\cref{eq:actual-holomorphic-strip} make the integrand jointly
holomorphic on a neighbourhood of the indicated compact complex set.
Cauchy's coefficient formula justifies interchange of extraction and
integration and bounds $|d_2|$ there by $M/s^2$.

Contour shifting bounds the Fourier coefficient at $r$ by
$(M/s^2)e^{-s|r|_1}$. The tensor grid retains exactly the Fourier
indices $r=Nl$. Summing the nonzero aliases gives
\cref{eq:theta-quadrature}. \src{F12}{V5}
\end{proof}
For direct evaluation, the coefficients of $P(k+p)$ are
\[
 P_0=P,\qquad P_1=-PH'P,\qquad
 P_2=PH'PH'P-\tfrac12PH''P,
\]
and
\[
 d_2=\tfrac12\Tr\left\{PC_2-
             \sum_{a+b+c=2}P_aV_bP\widetilde V_c\right\}.
\]
All quantities are Taylor coefficients, including their factorials.
The opposite vertex must be differentiated in both its arguments, as
must the momentum-dependent minimizing lift. A uniform absolute error
$\varepsilon$ in each $d_2$ contributes at most $\varepsilon$ to the
average, in addition to separately bounded summation errors.

The strip in \cref{eq:actual-holomorphic-strip} is very conservative.
The supplied calculation does not populate a practically useful bound
$M$ and a certified numerical enclosure of $\theta$. Thus no sign of
$\theta$, and no identification with the all-level physical coefficient
$b_{\rm YM}$ of \cref{eq:bYM}, follows from this representation.

\section{Exact flat-background cancellation and finite-volume aliases}
\label{sec:flat-bloch-cancellation}
For the $\SU(2)$ constant flat coarse field in direction $1$, with
amplitude $tT$ and $\norm T_{\rm HS}^2=2$, put
$q(t)=e^{\ii t}\sin(t)/t$. The actual constrained charged fibre is
related to the free fibre at $k'=k+2te_1$ by a $45$-dimensional
coordinate map $J_t(k)$ satisfying
\begin{equation}\label{eq:flat-fibre-congruence}
 \det_{\C}J_t=q(t)^7,\qquad
 H_t^{\rm ch}(k)=J_t(k)^*H^{\rm free}(k')J_t(k).
\end{equation}
There are eight fine boundary coordinate factors and one coarse factor;
the exponent seven is their quotient. The neutral fibre is unchanged.
Taking the determinant in \cref{eq:flat-fibre-congruence} contributes
$14\log(\sin t/t)$. The fine Haar density contributes
$16\log(\sin t/t)$ and the coarse Haar factor contributes
$2\log(\sin t/t)$, so the normalized measure cancels that term exactly.

\begin{theorem}[Flat one-loop identity at every real momentum]
\label{thm:flat-bloch}
Let $f(k)=\log\det H^{\rm free}(k)$, in the charged-coordinate
normalization of \cref{eq:flat-fibre-congruence}. The normalized
order-one integrand changes by
\[
 \delta\widehat q_t(k)=f(k+2te_1)-f(k),\qquad
 \mathcal I_0(k)=2\partial_{k_1}^2f(k).
\]
Consequently its infinite-volume flat variation vanishes exactly. On a
side-$L$ coarse torus, if $f_r$ are the Fourier coefficients of $f$,
\begin{align}
 \frac{\widehat Q_L(c(t))-\widehat Q_L(0)}{L^4}
   &=\sum_{l\in\Z^4}f_{Ll}(e^{2\ii tLl_1}-1),\label{eq:flat-alias}\\
 L^{-4}\sum_{k\in\mathcal G_L}\mathcal I_0(k)
   &=-2L^2\sum_{l\ne0}l_1^2 f_{Ll}.\label{eq:flat-alias-second}
\end{align}
\end{theorem}
\begin{proof}
The charged real determinant contributes one complex log determinant
to $\tfrac12\log\det_{\R}H$. Combining the congruence and the exact
Haar cancellation gives the first identity. Half its second amplitude
derivative is $2\partial_{k_1}^2f$, the indicated unit-polarization
normalization. Translation invariance of Haar measure on momentum
torus makes the integral of the difference zero.

The gapped fibre gives a periodic holomorphic logarithm of the positive
real-torus determinant on a smaller strip. Its Fourier coefficients
are exponentially summable, including two derivatives. A side-$L$
grid averages $e^{\ii k\cdot r}$ to zero unless $r=Ll$, proving
\cref{eq:flat-alias,eq:flat-alias-second}. \src{F12}{V9}
\end{proof}
The reported side-two constant mean $7499/15680$ is such a Fourier
alias. It is neither a nonzero bulk constant nor a correction which
may be subtracted from a nonflat quadratic card to produce $\theta$.
The cancellation is integrated, not pointwise. In particular, the
absence of a constant term does not determine the transverse $p^2$
coefficient of \cref{thm:theta-integral}.

\section{What remains between the local box and the full compact law}
\label{sec:compact-local-interface}
The global completed integral and the controlled local integral solve
different parts of the construction. The former has a fixed compact
carrier and exact score identities; the latter has the stronger
localized derivative estimates in \cref{thm:uniform-local-loop}.
Neither positivity of the former nor concentration in the latter
bounds the logarithmic derivatives of their ratio without additional
normalized estimates.

For a retained sigma-field $\mathscr S$, the exact identity
\begin{equation}\label{eq:conditional-means-retained}
 \Cov(f,g)=\E\Cov(f,g\mid\mathscr S)
          +\Cov\bigl(\E[f\mid\mathscr S],\E[g\mid\mathscr S]\bigr)
\end{equation}
keeps the conditional-mean term visible. A decay bound for the first
term alone cannot be transferred to the full covariance. Similarly,
a small exceptional-set probability is not a bound on the Dirichlet
energy of every normalized function concentrated on that set.
These distinctions are already present in the normalized source
calculus and the core--bad-set compiler; they remain essential after
compact completion. \src{F12}{V19, V23--V26, V33}

The nonlinear construction therefore strengthens the input to the
owner and forest analysis of \cref{ch:remainder}, without discharging
its cofinal uniformity hypotheses. In the auxiliary spectral problem,
\cref{ch:compact-repair} makes the obstruction concrete: a well with
exponentially small mass can still prevent a temperature-uniform
hidden diffusion gap. The valid replacement uses same-law visible
repair and states separately the coverage needed to turn local
estimates into an unrestricted interface inequality.

\chapter{Nonlinear remainder protection and sector discipline}
\label{ch:remainder}
The local compact construction of \cref{ch:compact-response} supplies
an actual action and density to which this bookkeeping can be applied.
Its $O(\beta^{-1/2})$ one-loop Hessian estimate is a restricted-box
estimate in the displayed local norm. It is not substituted for the
stronger source-dependent owner stock or its cofinal propagation
assumptions below. Localized well removal in \cref{prop:well-removal}
is likewise an explicitly compared finite contribution, not a deletion
of the entire bad-field class.

\section{An exhaustive owner ledger}
The remainder analysis begins with the one-scale compact analytic and ownership estimates. At a finite regulator, fixed retained zero-theta sector, and exact triangular elimination, primitives carry birth tags for regular local, bad, reset, or genuine boundary contributions. The Gaussian principal shell and extracted finite jets are outside this remainder alphabet.

A connected word is classified by the first applicable alternative
\[
\text{boundary}\succ\text{bad}\succ\text{reset}\succ
\text{regular multi-vertex}\succ\text{regular single-vertex}.
\]
This is an exhaustive disjoint partition, not an additional payment rule. Historical support and ownership are unchanged. Under normal convergence on the common smaller source polydisc, the normalized marginal response therefore has the exact identity
\begin{equation}\label{eq:owner-ledger}
\mathcal R_{\rm exact}=\mathcal R_{\rm main}
+E_{\rm vert}+E_{\rm polymer}+E_{\rm bad}
+E_{\rm reset}+E_{\rm multiscale}.
\end{equation}
The last term includes the genuine boundary class and the endpoint of propagated stable history. Derivatives pass through the series only after the normal-convergence hypothesis is supplied. Identity \cref{eq:owner-ledger} is not by itself a smallness estimate. \src{PAR}{V1}

\section{Power-preserving nonlinear control}
The common primitive stock is built from the two-jet-subtracted regular density, the normalized kernel defect, saturated exits, and compensation terms. For $u=\beta^{-1}$, the owner expansion gives, on its declared common domain,
\begin{equation}\label{eq:owner-power}
w(u)\le C_wu^{6/5}.
\end{equation}
The stretched exponentials are converted to this common power by the elementary maximum of $\beta^Ae^{-c\beta^d}$. Radius restriction gives, for the high-grade local density,
\[
\norm{D^kV_{\rm hi}}_r\le\frac{k!C_{\rm loc}}{(R-r)^k}u^{6/5},
\qquad k=0,1,2,
\]
with the zeroth-order formula interpreted on its original radius. Source derivatives have their corresponding Cauchy-radius losses. The marked connected-polymer and bad-field bounds use the same $w$, and the stable-history propagation preserves the $6/5$ power under the imported cone hypothesis. \src{PAR}{V2--V3}

\begin{criterion}[Physical marginal remainder and sign protection]\label{thm:remainder}
Assume the common-regulator owner, analytic-domain, fixed-sector, and stable-history hypotheses of the owner expansion. Suppose
\[
\beta_{j+1}=\beta_j+c_j,\qquad c_j=b_j+e_j,
\qquad |e_j|\le C_{\rm NP}u_j^{6/5},\quad |c_j|\le M_c,
\]
and $M_cu_j\le1/2$. Then
\begin{equation}\label{eq:physical-remainder}
u_{j+1}=u_j-b_ju_j^2+r_j,\qquad
|r_j|\le C_{\rm NP}u_j^{16/5}+2M_c^2u_j^3.
\end{equation}
If also $|b_j|\ge m_{\rm main}>0$, then at sufficiently small common $u_j$,
\begin{equation}\label{eq:sign-protection}
\frac{|r_j|}{|b_j|u_j^2}
\le\frac{C_{\rm NP}}{m_{\rm main}}u_j^{6/5}
+\frac{2M_c^2}{m_{\rm main}}u_j\le\frac12.
\end{equation}
Thus the exact real increment has the sign of $-b_ju_j^2$.
\end{criterion}
\begin{proof}
Use the exact identity
\[
\frac{u}{1+uc}=u-cu^2+\frac{c^2u^3}{1+uc}.
\]
Subtract the main increment $-bu^2$, bound $e$ and the denominator, and divide by $|b|u^2$. The common smallness threshold also retains every inherited owner and analytic-domain condition. \src{PAR}{V4}
\end{proof}

This is a protection theorem for a supplied principal margin. It does
not compute $b_j$, prove its sign, or identify its limiting value with
$b_{\rm YM}$. The principal coefficient must be that of the same calibrated
exact response. A margin measured with fixed-background Wilson banks
cannot be inserted as the physical principal margin before the background
and complete-determinant comparisons of \cref{sec:background-matching}
are established. Those discrepancies are not silently absorbed into the
$O(u^{6/5})$ nonlinear remainder. The principal-response calculation and the
remainder theorem are complementary, not substitutes.

\section{Zero theta, compact labels, and the correct thermodynamic target}
The exact compact construction retains its sector labels and the declared zero-theta section through the chosen transport. Reflection positivity is not used as a general argument that every possible theta parameter must vanish. Gauge quotient labels, commuting torons, and global charge labels require their own treatment.

A proposed volume-uniform tail for total sector complexity encounters a precise tensorization obstruction. A local nonzero charge costs finite action and finite one-component free energy, while its number of possible placements grows with volume. A global total-charge bound cannot be inferred from a local defect penalty.

\begin{theorem}[Tensorization obstruction to global sector tightness]\label{thm:sector-no-go}
Let $X$ be the bounded integer charge on a fixed finite component, and suppose $X$ is not almost surely zero. On $N$ disjoint copies, let $X_i$ be independent copies and set
\[
\Psi_N=\sum_{i=1}^N|X_i|,\qquad Q_N=\sum_{i=1}^NX_i.
\]
For every fixed $R<\infty$,
\begin{equation}\label{eq:sector-escape}
\Pr(\Psi_N\ge R)\longrightarrow1,\qquad
\Pr(|Q_N|\ge R)\longrightarrow1.
\end{equation}
Consequently no bound $Ce^{-cR}$ with positive $c$ and finite $C$ independent of $N$ controls either tail for all $N,R$.
\end{theorem}
\begin{proof}
If $p=\Pr(|X|\ge1)>0$, then $\Psi_N$ dominates a sum of independent Bernoulli variables with mean $Np$. If $\E X\ne0$, the law of large numbers gives escape of $Q_N$. If $\E X=0$, nontriviality gives positive variance and the central limit theorem implies that its probability in any fixed bounded interval tends to zero. A nonzero constant charge is covered by the nonzero-mean case. Choosing $R$ with $Ce^{-cR}<1/2$ contradicts the limiting probabilities. \src{PAR}{V5}
\end{proof}

The theorem applies exactly to the locality-compatible class allowing disjoint unions. It does not by itself prove anti-concentration on a sequence restricted to connected periodic tori; that would require a positive susceptibility or a suitable insertion/mixing estimate. Fixed physical volume can also change the target. The retained alternatives are local-window or charge-density bounds, explicit phase or sector decompositions, and exact compact transport. The global extensive tail is not an unpaid theorem that one should continue trying to obtain from the same insufficient energy--entropy mechanism.

\section{Refinement coherence is retained independently}
Let $P_j$ be the field block map, $q_j$ the label map, and $r_j$ the label refinement. Suppose $r_jq_{j+1}=q_jP_j$ outside an event of probability $\epsilon_j$, and
\[
d_j=\norm{(P_j)_*\mu_{j+1}-\mu_j}_{\rm TV}.
\]
Then data processing gives
\begin{equation}\label{eq:label-TV}
\norm{(r_j)_*(q_{j+1})_*\mu_{j+1}-(q_j)_*\mu_j}_{\rm TV}
\le\epsilon_j+d_j.
\end{equation}
Summable defects yield a projectively Cauchy family at every fixed level. This coherence estimate survives the tightness obstruction. Coherence and tightness are different: a coherent family can still place its mass on escaping labels. \src{PAR}{V5}
\sourcebasis{The exact owner partition, power bounds, stable-history control, and physical remainder are \src{PAR}{V1--V4}; the topology and refinement results are \src{PAR}{V5}.}

\chapter{Compact spectral geometry and local interface repair}
\label{ch:compact-repair}
The compact completion permits an auxiliary spectral analysis on the
actual fibre, without replacing it by a Gaussian or an independent
product reference. Two features must be treated together. A normally
coercive symmetry orbit supports a quantitative local Poincar\'e
inequality, but a second, higher-energy orbit creates metastability
for the full hidden diffusion. Local visible repair can nevertheless
be proved before integrating the hidden coordinates and transferred
to the original interface heat bath. This chapter develops that route
and its remaining local-to-global requirement. Every heat-bath rate
below is an auxiliary coordinate-update rate, not the physical clock
of \cref{ch:spectral}.

\section{The actual finite parent and its stationary orbits}
\label{sec:compact-orbit-data}
Use the $\SU(2)$ completion of
\cref{sec:su2-pivot-equation}, with $\rho=1/64$ and action convention
$s_p=2-\Tr Q_p$. On the fine torus of side six, fix the coarse field
$C=I$ and the parent of coarse cells $P=\{0,1\}^4$. Its independent
hidden carrier is
\[
 M=\SU(2)^{944},\qquad 944=272+672,\qquad \dim M=2832.
\]
The $272$ internal non-tree groups and $672$ nonpivot boundary groups
are independent; the $96$ pivots are smooth functions of them. Fix
the central exterior sheet specified by
\begin{equation}\label{eq:central-sheet}
 \epsilon_{(x,\mu)}=
 \begin{cases}
 -1,&\mu=2,\ x_2=1,\text{ and some }x_a\text{ is odd},\ a\in\{0,1,3\},\\
 +1,&\text{otherwise}.
 \end{cases}
\end{equation}
The parent law, including its inherited additive action normalization,
is
\begin{equation}\label{eq:actual-parent-law}
 \dd\mu_\beta=Z_\beta^{-1}a(Y)e^{-\beta S(Y)}\dd\operatorname{vol}_M(Y).
\end{equation}
Here $a$ is the full positive inverse-pivot amplitude, independent of
$\beta$, and the metric is product Hilbert--Schmidt.

\begin{certificate}[Two stationary orbits]
\label{def:two-orbit-input}
For this fixed parent and exterior, the finite certificates locate two
full stationary simultaneous-conjugation orbits
$\mathcal O_1,\mathcal O_2\simeq S^2$. Their Hessian kernels are
exactly the two orbit directions. On the $2830$-dimensional normal
spaces the respective lower bounds are $1/5000$ and $31/125000$ in
the Hilbert--Schmidt metric. Their energies satisfy
\[
 \frac{113936882369}{200000000}\le E_1
          \le\frac{569684411883}{10^9},\qquad
 569.68444636<E_2<569.68444638,
\]
and hence $E_2-E_1>34477/10^9>0$.
\end{certificate}
The first orbit and its full normal Hessian are reported in
\src{F12}{V89--V90}; the second, including the transition-source pivot
acceleration, is reported in \src{F13}{V8}. The associated interval
and integer factor certificates are not independently rerun here.
The analytic results below use these explicitly identified finite
inputs. No classification of the global ground set is assumed.
In particular, the second orbit is not globally minimizing, whereas
global minimality of the first is not established by these data.

\section{An equivariant tube inequality with the actual amplitude}
\label{sec:actual-orbit-tube}
Choose $p\in\mathcal O_1$ and let $K\simeq\mathrm U(1)$ be its
stabilizer. The normal bundle is $G\times_KN_p\mathcal O_1$, not an
assumed globally trivial product. For a sufficiently small fixed
normal ball $B_\varepsilon\subset N_p\mathcal O_1$, the map
\[
 F(g,z)=g\cdot\exp_pz
\]
parametrizes the orbit tube $\mathcal T_\varepsilon$ through that
associated bundle. Invariance of the action, amplitude, and metric
gives a smooth positive $K$-invariant function $b$ such that
\begin{equation}\label{eq:tube-law-pushforward}
 F_*(\dd g\otimes\dd\nu_\beta)
      =\mu_\beta(\,\cdot\mid\mathcal T_\varepsilon),\qquad
 \dd\nu_\beta(z)=
 \frac{b(z)e^{-\beta s(z)}\dd z}
      {\int_{B_\varepsilon}b(w)e^{-\beta s(w)}\dd w},
 \quad s(z)=S(\exp_pz).
\end{equation}
Indeed $b$ is the original amplitude times the normal-exponential
volume Jacobian. Haar measure pushes forward to normalized orbit
volume, while equivariance transports the same normal density around
the orbit. This retains the Jacobian on the whole tube; it does not
replace it by its value at $p$.

\begin{theorem}[Local hidden diffusion inequality]
\label{thm:actual-tube-gap}
For a fixed sufficiently small $\varepsilon>0$ and a finite threshold
$\beta_0$, every $\beta\ge\beta_0$ and every weighted $H^1$ function
on the tube satisfy
\begin{equation}\label{eq:actual-tube-gap}
 \Var_{\mu_\beta(\cdot\mid\mathcal T_\varepsilon)}f
 \le\left(3776+\frac{80000}{\beta}\right)
     \int_{\mathcal T_\varepsilon}|\nabla_Mf|^2
                              \dd\mu_\beta(\cdot\mid\mathcal T_\varepsilon).
\end{equation}
For conjugation-invariant $f$, the constant $3776$ can be omitted.
The local Neumann diffusion gap is therefore at least $1/83776$ for
$\beta\ge\max\{1,\beta_0\}$.
\end{theorem}
\begin{proof}
Normal positivity in \cref{def:two-orbit-input} and continuity permit
$D^2s\succeq I/10000$ and $\norm{D\exp_p|_N}\le2$ on the fixed
ball. With $K_b=\sup\norm{D^2\log b}$, take
$\beta_0\ge20000K_b$. Then
$D^2(\beta s-\log b)\succeq(\beta/20000)I$.
The weighted Bochner identity on a convex Euclidean ball, including
the nonnegative second-fundamental-form boundary term, gives the
Neumann Poincar\'e constant $20000/\beta$ for $\nu_\beta$.

Haar measure on $\SU(2)$ in this metric has a sufficient Poincar\'e
constant one: the group is the round three-sphere of radius $\sqrt2$
and its Ricci tensor equals the metric, so the integrated Bochner
identity yields this bound. For $h=f\circ F$,
\[
 |D_zh|^2\le4|\nabla_Mf|^2\circ F,\qquad
 |\nabla_Gh|^2\le4\cdot944|\nabla_Mf|^2\circ F.
\]
The latter uses $\norm{[Z,U]}_{\rm HS}\le2\norm Z_{\rm HS}$
in each of the $944$ coordinates; it is an operator-norm bound and
has no additional factor $\dim G$. Variance decomposition in the
product law of \cref{eq:tube-law-pushforward}, the two Poincar\'e
bounds, and Jensen's inequality for the differentiated Haar average
give \cref{eq:actual-tube-gap}. For invariant $f$ the group derivative
vanishes. Density of smooth functions in the weighted form domain
completes the proof. \src{F13}{V1}
\end{proof}
The tube radius, amplitude derivative bound, and onset $\beta_0$ are
existential constants, not numerically enclosed parameters. The
inequality is a local Neumann statement, with no boundary condition
imposed on its test functions. It is distinct from a killed escape
inequality and from the physical orbit criterion in \cref{thm:orbit}.

\section{A secondary well excludes a uniform hidden diffusion gap}
\label{sec:hidden-metastability}
Let $\lambda_\beta$ be the gap of the full hidden diffusion with
Dirichlet form
$\mathcal E_\beta(f)=\int_M|\nabla f|^2\dd\mu_\beta$.

\begin{theorem}[Metastability of the actual fixed-parent diffusion]
\label{thm:hidden-metastability}
For the law in \cref{eq:actual-parent-law} and the finite inputs in
\cref{def:two-orbit-input}, constants $C<\infty$, $\delta>0$, and
$\beta_1<\infty$ satisfy
\begin{equation}\label{eq:hidden-metastability}
 0<\lambda_\beta\le C\beta^{1415}e^{-\beta\delta}
                   \qquad(\beta\ge\beta_1).
\end{equation}
The same upper bound holds for the lowest Dirichlet eigenvalue on the
complement of any sufficiently small closed tube about $\mathcal O_1$
which is disjoint from a neighbourhood of $\mathcal O_2$.
\end{theorem}
\begin{proof}
Choose a small normal tube around $\mathcal O_2$, disjoint from
$\mathcal O_1$. Normal positivity gives
$E_2+c_-|z|^2\le S\le E_2+c_+|z|^2$ in finitely many tubular
charts. Let $f$ equal one on its radius-$r$ inner tube and vanish
outside the radius-$2r$ tube. Its gradient is supported where
$S\ge E_2+\delta$, with $\delta=c_-r^2>0$. Smooth positive amplitude
and volume bounds give
\[
 \mathcal E_\beta(f)\le Z_\beta^{-1}C e^{-\beta(E_2+\delta)},\qquad
 \int f^2\dd\mu_\beta\ge Z_\beta^{-1}c\beta^{-1415}e^{-\beta E_2}.
\]
The exponent is half the normal dimension $2830$. The lower-energy
orbit independently gives
$Z_\beta\ge c'\beta^{-1415}e^{-\beta E_1}$. Thus
$\mu_\beta(\supp f)\le C'\beta^{1415}e^{-\beta(E_2-E_1)}\to0$.
Cauchy--Schwarz implies
\[
 \Var_{\mu_\beta}f\ge(1-\mu_\beta(\supp f))\int f^2\dd\mu_\beta.
\]
Its Rayleigh quotient proves the upper bound. The density is positive
and smooth on a compact connected manifold at every finite $\beta$,
so its elliptic diffusion gap is positive. The same supported test
function belongs to the Dirichlet domain outside the first tube;
its uncentered Rayleigh quotient proves the killed bound.
\src{F13}{V8}
\end{proof}
The barrier $\delta$ is not the energy difference $E_2-E_1$.
Exponential smallness of the higher well's equilibrium mass does not
remove its small Rayleigh quotient. Covering both wells by separate
local tubes therefore cannot restore a temperature-uniform gap for
this same diffusion. This obstruction does not assert a corresponding
upper bound for a nonlocal heat bath and does not concern the physical
Yang--Mills Hamiltonian. It removes one sufficient auxiliary route,
not the continuum target.

\subsection{Localized marginal removal does not require hidden mixing}
\label{sec:well-removal}
There is a useful simultaneous positive conclusion. Let a fixed smooth
cutoff $0\le\chi\le1$ be supported in the second orbit tube and vanish
near the first. For coarse parameters $q$ near the declared background,
write $N_\beta(q)$ for the integral with factor $\chi$,
$p_\beta=N_\beta/Z_\beta$, and
$\widetilde Z_\beta=Z_\beta-N_\beta=Z_\beta(1-p_\beta)$.
This is a specified comparison integral, not a redefinition of the
original law.

\begin{proposition}[A local normalized well-comparison estimate]
\label{prop:well-removal}
On a sufficiently small fixed coarse neighbourhood $U$, and for every
fixed integer $k\ge0$,
\begin{align}
 0\le p_\beta(q)&\le C\beta^{1415}e^{-\sigma\beta},
                  &\sigma&=1/50000,\label{eq:well-mass}\\
 \norm{\log Z_\beta-\log\widetilde Z_\beta}_{C^k(\overline U)}
        &\le C_k\beta^{1415+k}e^{-\sigma\beta}.\label{eq:well-derivatives}
\end{align}
No hidden Poincar\'e inequality is required.
\end{proposition}
\begin{proof}
Shrink $U$ so $|S(q,x)-S(q_0,x)|\le\gamma=1/200000$ uniformly
in the compact hidden carrier. The numerator is at most
$Ce^{-\beta(E_2-\gamma)}$, while the first orbit supplies a denominator
at least $c\beta^{-1415}e^{-\beta(E_1+\gamma)}$.
The energy intervals give $E_2-E_1-2\gamma>\sigma$, proving
\cref{eq:well-mass}. Derivatives of the actual integrand divided by
that integrand are bounded by $C_\alpha(\beta+1)^{|\alpha|}$.
The fixed nonnegative cutoff contributes no parameter derivatives.
Induction in the differentiated identity $Z_\beta p_\beta=N_\beta$
gives
$|\partial^\alpha p_\beta|\le C_\alpha(\beta+1)^{|\alpha|}p_\beta$.
Differentiate $\log(1-p_\beta)$ and use $p_\beta\le1/2$ for large
$\beta$. This proves \cref{eq:well-derivatives}, including the
normalization derivatives. \src{F13}{V9}
\end{proof}
For normalized retained laws with the same direct retained weight and
partition factors $Z_\beta$ and $\widetilde Z_\beta$ on a compact
domain in $U$, the density-ratio oscillation is at most
$(1-\varepsilon_\beta)^{-1}$, where
$\varepsilon_\beta=\sup_U p_\beta$. Their same-metric diffusion gaps,
same-block heat-bath gaps, and killed floors on the same event therefore
obey the two-sided comparison factors
$1-\varepsilon_\beta$ and $(1-\varepsilon_\beta)^{-1}$.
For heat baths this follows directly from
\[
 \E_\nu\Var(f\mid q_{B^c})
    =\inf_{g\ \mathscr F_{B^c}\text{-measurable}}\int(f-g)^2\dd\nu,
\]
so the compared update blocks and rates are not changed. This
localized removal estimate does not classify the remaining wells or
supply an unrestricted retained gap.

\section{A same-law diffusion-to-heat-bath comparison}
\label{sec:global-resolvent}
Let $\mu$ be a positive smooth density proportional to $e^{-W}$
relative to product Haar in a finite volume. Write
$L=\sum_iL_i\succeq0$ for its reversible product-metric diffusion,
$\Gamma_i$ for the conditional expectation given all coordinates
except group $i$, and
\[
 \mathsf H_\mu=\sum_i(I-\Gamma_i),\qquad
 \mathcal E_L(u)=\ip u{Lu},\qquad
 \mathcal E_{\mathsf H_\mu}(f)=\sum_i\norm{(I-\Gamma_i)f}_2^2.
\]
The symbol $\mathsf H_\mu$ denotes the rate-one coordinate heat bath;
it is not the reconstructed physical Hamiltonian $H$.
Assume the intrinsic mixed Hessian row bound
\[
 \sup_i\sum_{j\ne i}\sup\norm{W_{ij}^{\rm Hess}}_{\op}\le K.
\]

\begin{theorem}[Global resolvent domination]
\label{thm:global-resolvent}
On the same $L^2(\mu)$ space,
\begin{equation}\label{eq:global-resolvent}
 \mathsf H_\mu\succeq L(L+K)^{-1}.
\end{equation}
The right side is defined by spectral calculus, with value zero on
$\Ker L$; if $K=0$ it is the projection onto $(\Ker L)^\perp$.
In particular a diffusion gap $g>0$ implies
$\gap\mathsf H_\mu\ge g/(g+K)$.
\end{theorem}
\begin{proof}
The product-manifold and one-coordinate weighted Bochner identities
give
\[
 \sum_i\norm{L_iu}_2^2\le\norm{Lu}_2^2+K\mathcal E_L(u).
\]
All diagonal curvature and Ricci terms cancel on subtraction; the
remaining mixed Hessian squares are nonnegative, and the mixed
potential terms are bounded by the row sum and
$ab\le(a^2+b^2)/2$. No commutation of the $L_i$ is assumed.
Since $\Gamma_iL_iu=0$,
\[
 \ip f{Lu}=\sum_i\ip{(I-\Gamma_i)f}{L_i u}.
\]
Completing a square in this direct sum yields, for every smooth $u$,
\[
 2\operatorname{Re}\ip f{Lu}-\norm{Lu}_2^2-K\mathcal E_L(u)
             \le\mathcal E_{\mathsf H_\mu}(f).
\]
Taking the supremum over $u$ by the spectral theorem gives
\cref{eq:global-resolvent}. Closure of the forms extends the estimate
to its full domain. The scalar gap consequence follows by restricting
to the centered subspace. \src{F12}{V60, V77--V80}
\end{proof}
For the unintegrated completed fibre,
$W_\beta=\beta S-\log j$ is a sum of fixed smooth finite-stencil
terms. Bounded incidence and compact input derivatives give
\begin{equation}\label{eq:preintegrated-row}
 K_\beta\le C_{\rm f}(\beta+1)
\end{equation}
uniformly in the ambient volume and compact coarse data. After hidden
integration, the score covariance in \cref{eq:compact-score} can
instead produce a quadratic-in-$\beta$ Hessian bound. Keeping
\cref{eq:global-resolvent} before integration avoids this unnecessary
loss; it does not posit a uniform hidden diffusion gap.

\begin{lemma}[Exact domination of marginal heat-bath forms]
\label{lem:marginal-form-domination}
Let $\mu(\dd x,\dd v)$ be the joint fine law and $\nu(\dd v)$ its
exact retained interface marginal, using the same independent group
partition. For $Jf(x,v)=f(v)$,
\[
 \norm{Jf}_{L^2(\mu)}=\norm f_{L^2(\nu)},\qquad
 \mathcal E_{\mathsf H_\nu}(f)\ge
                         \mathcal E_{\mathsf H_\mu}(Jf).
\]
\end{lemma}
\begin{proof}
Hidden-group updates leave $Jf$ unchanged. For a retained group,
conditioning additionally on the hidden variables reduces the expected
conditional variance, by variance decomposition. Summing the retained
coordinate inequalities proves the claim.
\end{proof}
To transfer a killed estimate by this lemma, the event must constrain
visible coordinates only. Its lift must contain every hidden field;
a hidden force-region restriction would introduce a further posterior
coverage problem.

\section{From visible negative curvature to an original-clock repair floor}
\label{sec:repair-transfer}
For an operator $A\succeq0$ and a visible event $D$, write
\[
 \kappa_D^A=\inf_{0\ne f=\one_Df}
                 \frac{\ip f{Af}}{\norm f_2^2}
\]
for its killed floor, with the infimum taken in the appropriate form
domain. It is not a reflected or Neumann gap on $D$.

\begin{lemma}[Spectral localization in a visible cylinder]
\label{lem:repair-localization}
Suppose the same global diffusion satisfies, for every form vector $u$,
\begin{equation}\label{eq:repair-localization}
 \norm{\one_Du}_2^2\le\frac2\lambda\mathcal E_L(u)
                           +\frac{2M_\chi}{\lambda}\norm u_2^2,
 \qquad \lambda\ge8M_\chi.
\end{equation}
Together with \cref{eq:global-resolvent}, this gives
\begin{equation}\label{eq:repair-floor-general}
 \kappa_D^{\mathsf H_\mu}\ge
                 \frac{\lambda}{2(\lambda+8K)}.
\end{equation}
The identical lower bound transfers to the exact marginal heat bath
when $D$ is a cylinder restricting only retained coordinates.
\end{lemma}
\begin{proof}
For the global diffusion spectral projection
$\Pi_t=\one_{[0,t]}(L)$, \cref{eq:repair-localization} gives
$\norm{\one_D\Pi_t}^2\le2(t+M_\chi)/\lambda$.
Taking adjoints gives the same estimate for $\Pi_t\one_D$; no
commutation is asserted. At $t=\lambda/8$ the bound is at most $1/2$,
so a vector supported in $D$ has at least half its squared norm above
$t$. The global resolvent bound gives
$\ip f{\mathsf H_\mu f}\ge\lambda\norm f^2/[2(\lambda+8K)]$.
Use \cref{lem:marginal-form-domination} for the last assertion.
\end{proof}
The resolvent in this proof is always that of the global diffusion for
the same law. A Dirichlet resolvent is not substituted for it.

Here is the local geometric input to \cref{eq:repair-localization}.
In product normal coordinates on a finite set of visible groups,
leave all other groups uncharted. If a constant unit chart direction
$a$ obeys
$\partial_a^2\Phi_\beta\le-\eta\beta$ for the actual measured
potential, uniformly in all uncharted inputs, and the inverse metric
is at least $c_gI$, then ground-state conjugation gives
\[
 \mathcal E_L(h)\ge\frac{c_g\eta\beta}{2}\norm h^2
\]
for functions supported in that chart cylinder. Indeed, with
$u=h e^{-\Phi_\beta/2}$, the one-direction potential term is
$|\partial_a\Phi_\beta|^2/4-\partial_a^2\Phi_\beta/2$.
A cutoff $\chi$ equal to one on a smaller cylinder yields
\cref{eq:repair-localization} with
$\lambda=c_g\eta\beta/2$ and $M_\chi=\sup|\nabla\chi|^2$.
For a fixed finite cover, normalized cutoffs with
$\sum_j\chi_j^2\le1$, equal to one on the event, give the same estimate
with $M_\chi=\sup\sum_j|\nabla\chi_j|^2$.

Thus any such fixed visible certificate with $\lambda=c_0\beta$ and
\cref{eq:preintegrated-row} gives, for sufficiently large $\beta$,
\begin{equation}\label{eq:uniform-repair-template}
 \kappa_D^{\mathsf H_\nu}\ge
 \frac{c_0\beta}{2\{c_0\beta+8C_{\rm f}(\beta+1)\}}
 \ge\frac{c_0}{2(c_0+16C_{\rm f})}>0.
\end{equation}
The constant is in the original rate-one single-group clock; no
collective rotation has been inserted into the generator.

\section{A non-wrapping all-hidden-field shell}
\label{sec:nonwrapping-repair}
The following geometry is defined on the integer lattice before
embedding into a torus. It therefore retains the outer boundary cost
which a small periodic shell would miss. Keep the parent $P$ and the
central signs \cref{eq:central-sheet}. For $q\ge3$ set
\[
 Q_q=\{-1,0,\ldots,q-2\}^3\setminus\{0,1\}^3
\]
and select the seven direction-$2$ free links
\[
 ((2a+u,2b+v,1,2c+w),2),\quad
 (a,b,c)\in Q_q,\quad
 (u,v,w)\in\{0,1\}^3\setminus\{0\}.
\]
Their set $T_q$ has size $7(q^3-8)$. A selected link occurs in two
nonroot paths over its own coarse edge and in no other pivot source
family. Call a plaquette uncontrolled when it contains a parent-hidden
physical link or a parent-dependent pivot; the other incident
plaquettes are determined by retained inputs.

The finite incidence enumeration gives the following data. The last
column is the upper second derivative under a common selected-link
rotation, with arbitrary uncontrolled fields.
\begin{center}\small
\begin{tabular}{rrrrrrr}
\toprule
$q$ & $|T_q|$ & Two selected & Fixed $-$ & Fixed $+$ & Uncontrolled & $S''$ upper\\
\midrule
3 &133&261&102&126&48&144\\
4 &392&828&324&288&84&96\\
5 &819&1845&690&450&84&$-312$\\
\bottomrule
\end{tabular}
\end{center}
Use $q=5$, write $T=T_5$, and let $F$ be the following sufficient
visible collar. Include $T$, all free inputs of its $117$ source
families, and all free inputs of the fixed incident plaquettes. For
this closure, discard tree links and replace a pivot by the free
links in its fourteen nonroot paths; those paths contain no pivots,
so no further iteration is required. The reported finite data are
\begin{equation}\label{eq:shell-collar-data}
 |T|=819,\quad |F|=10034,\quad F\cap I_P=\varnothing.
\end{equation}
The closure of all incident plaquettes, including the uncontrolled
ones, uses $10431$ free groups. Its required vertices lie within
$[-3,9]\times[-3,9]\times[0,3]\times[-3,9]$. Translation by
$(8,8,8,8)$ embeds the whole closure without wraparound into every
even fine torus of side $L\ge32$. These integer-lattice counts and
source exclusions are the geometric certificates of
\src{F13}{V14}, not an extrapolation of the periodic side-six shell.

\begin{theorem}[Non-wrapping negative curvature and visible repair]
\label{thm:nonwrapping-repair}
At $C=I$, prescribe $Y_l=\epsilon_lI$ on $F$ and leave every other
free group arbitrary. Under $Y_l(t)=-e^{tE}$ for all $l\in T$, with
$E^2=-I$ and $\norm E_{\rm HS}^2=2$, the completed action obeys
\[
 S''(0)\le-312,\qquad
 \operatorname{Hess}S(a,a)\le-4/21
\]
for the corresponding unit direction $a$. If $F$ is retained in the
interface, a fixed smaller visible collar cylinder $D_L$ has the
original-clock floor \cref{eq:uniform-repair-template} with $c_0=1/84$.
All constants can be chosen independently of $L\ge32$ and of the
compatible hidden index sets disjoint from $F$.
\end{theorem}
\begin{proof}
The selected source products are near $-I$, outside the logarithmic
cutoff. The associated $117$ completed pivots are exactly $I$ on a
fixed neighbourhood and have inverse Jacobian one there. Other pivots
are independent of the selected motion. Two-selected plaquettes have
opposite orientations and cancel the common rotation exactly. A fixed
negative one-selected plaquette contributes $-2$ to $S''$, a positive
one contributes $2$, and each uncontrolled plaquette contributes at
most $2$ by unitarity of its other quaternion factors. Hence
\[
 S''(0)\le2(-690+450+84)=-312.
\]
The incidence check is $2\cdot1845+690+450+84=6\cdot819$.
The independent-coordinate speed squared is $2\cdot819=1638$, so
normalization gives $-312/1638=-4/21$.

Only a fixed finite compact input closure enters these derivatives.
Uniform continuity gives a visible chart on which
$\partial_a^2S\le-2/21$ for all uncharted fields. The inverse-pivot
amplitude is independent of the selected groups there; the
coordinate-Haar second derivative is bounded. For large $\beta$,
$\partial_a^2\Phi_\beta\le-\beta/21$. Choose inverse-metric bound
$c_g=1/2$. The local diffusion coefficient is then
$c_0=c_g/(2\cdot21)=1/84$. All chart and cutoff constants are
volume independent. Apply \cref{lem:repair-localization} and
\cref{lem:marginal-form-domination}; the visibility condition on $F$
ensures that the lifted event permits all hidden variables.
\end{proof}
The positive upper bounds for $q=3,4$ only mean that this particular
worst-case sufficient test fails there. The $q=5$ theorem is a local
repair statement, not an assertion that these collar events cover all
defects or have large probability.

\subsection{Transport to locally curved coarse backgrounds}
\label{sec:coarse-gauge-repair}
The needed coarse input closure consists of $550$ edges inside a
contractible box with eight vertices in each coordinate direction.
Assume only that the elementary coarse plaquettes in this box satisfy
$d(I,Q_p(C))\le\varepsilon$. In a monotone path gauge, move the last
step of an edge past at most $3\cdot7=21$ later-direction steps.
Each swap changes the transport by a conjugate plaquette holonomy.
Bi-invariance and the triangle inequality give a gauge $h$ with
\begin{equation}\label{eq:local-box-gauge}
 d(I,C_e^h)\le21\varepsilon
\end{equation}
for every edge in the box. This uses no commuting-field approximation.

The cell-constant fine gauge
$Y_{x,\mu}^h=h_rY_{x,\mu}h_s^{-1}$, where $r,s$ are endpoint coarse
cells, preserves tree links and is a product Haar-preserving isometry
of the independent groups. The completed pivot maps are equivariant,
so their inverse Haar Jacobians and the full fixed-coarse law transform
exactly. The same transformation conjugates the single-group heat
baths and every marginal with the same index partition. Since $h$
depends only on the fixed coarse input, no additional field-dependent
Jacobian is introduced.

Finite-input continuity of the shell certificate therefore supplies
an $\varepsilon_*>0$ such that the repair floor with $c_0=1/84$
persists, uniformly in ambient volume, whenever
$\varepsilon\le\varepsilon_*$. The visible cylinder is the pullback
of the central collar chart by this gauge. It does not require that
all coarse links globally be near identity or that wrapping holonomies
vanish. The constant $\varepsilon_*$ and the chart radius are not
numerically enclosed. \src{F13}{V15}

\section{A noncentral visible matrix certificate}
\label{sec:matrix-repair}
The shell need not remain near its central pattern. Keep the fixed
$T,F$ and coarse input set $K$ above. Let $\mathcal P_1$ be its $1140$
visible one-selected plaquettes and $\mathcal P_2$ its $1845$ visible
two-selected plaquettes. Impose the strict source margin
\begin{equation}\label{eq:inactive-margin}
 d_{\rm ch}(I,C_e^{-1}P_{e,j}(Y_F))\ge\rho+\sigma,
 \qquad \sigma>0,
\end{equation}
for all fourteen paths of the $117$ selected edges. Since
$d_{\rm ch}\le d$, the corresponding cutoffs are inactive in an
open neighbourhood; their actual pivots equal $C_e$ and their inverse
Jacobians are one. All other coarse and fine inputs can be arbitrary.

Write a quaternion as $(q_0,\mathbf q)$, with $\Tr q=2q_0$.
Cyclically start a two-selected plaquette at its positive selected
link, writing its word as $UAV^{-1}B=XB$, where $X=(x_0,x)$ and
$B=(b_0,b)$. Define the symmetric matrix
\begin{equation}\label{eq:visible-matrix}
 M_{\rm vis}=2\sum_{p\in\mathcal P_1}q_{p,0}I_3
  +8\sum_{p\in\mathcal P_2}
       \left\{\frac{x_pb_p^T+b_px_p^T}{2}-(x_p\cdot b_p)I_3\right\}.
\end{equation}
It is a smooth function of the actual completed visible plaquettes
and their fixed compact coarse inputs. It is a criterion in the
declared tree coordinates; a common rotation axis is not asserted to
be invariant under independent cell gauges.

\begin{theorem}[Matrix-certified repair with arbitrary hidden fields]
\label{thm:matrix-repair}
For fixed $\sigma,\eta>0$, let $D_C^{\sigma,\eta}$ be the visible
event defined by \cref{eq:inactive-margin} and
\begin{equation}\label{eq:matrix-negative-test}
 \lambda_{\min}(M_{\rm vis})\le-168-\eta.
\end{equation}
If $F$ is retained, the exact interface heat bath has the floor
\cref{eq:uniform-repair-template} on this event with
$c_\eta=\eta/26208$. Its constants are uniform in even $L\ge32$,
in all compact coarse inputs, and in compatible hidden index sets.
They may depend on the fixed shell, $\sigma$, and $\eta$.
\end{theorem}
\begin{proof}
Left multiply every selected free link by
$R_n(t)=\exp(tE_n)$, where $n\in\R^3$ is unit. A one-selected
plaquette contributes $2q_{p,0}$ to the second derivative. A
two-selected plaquette has moving trace $\Tr(R_nXR_n^{-1}B)$.
The identity
\[
 [E_n,[E_n,X]]=(0,4\{n(n\cdot x)-x\})
\]
gives its action second derivative
$8\{(n\cdot x)(n\cdot b)-x\cdot b\}$. Each of the $84$
uncontrolled plaquettes contributes at most two. Consequently
\[
 S''(0)\le n^TM_{\rm vis}n+168.
\]
Choose a least-eigenvalue axis. The speed squared is $1638$, so the
unit-direction Hessian is at most $-\eta/1638$ for all uncontrolled
fields.

The closed inequalities defining the event give a compact subset of
the fixed finite carrier $G^K\times G^F$. At each point choose one
axis and freeze it in orthonormal normal coordinates. Continuity and
the strict source margin give a chart with action second derivative
at most $-\eta/(2\cdot1638)$. The inverse-pivot amplitude is constant
in the selected variables; bounded coordinate-Haar terms give measured
curvature at most $-\beta\eta/(4\cdot1638)$ at large $\beta$.
With inverse metric at least $1/2$, the local diffusion coefficient
is $\eta/(16\cdot1638)=\eta/26208$.

A finite cover of this fixed compact carrier admits normalized smooth
cutoffs whose squared sum is one on the certified set and at most one
elsewhere, with a bounded sum of squared fine-coordinate derivatives.
No derivative of an eigenvector through a crossing is needed. This
gives \cref{eq:repair-localization} uniformly in $C$; apply the global
resolvent and exact marginal domination. If the event is empty, the
form inequality is vacuous. \src{F13}{V16}
\end{proof}
At the central collar, $M_{\rm vis}=-480I_3$, so adding the worst-case
$168I_3$ recovers $-312I_3$. The event is nonempty for $0<\eta<312$
and sufficiently small $\sigma$, and includes noncentral perturbations.
A weaker scalar sufficient test is
\[
 \tfrac13\Tr M_{\rm vis}
 =2\sum_{\mathcal P_1}q_{p,0}
       -\tfrac{16}{3}\sum_{\mathcal P_2}x_p\cdot b_p
                         \le-168-\eta.
\]
This is not equivalent to the least-eigenvalue test. The finite
covering here is over configurations on one fixed patch, not over a
number of spatial patches growing with volume.

\section{Repair with genuinely active pivot sources}
\label{sec:active-repair}
The inactive-source margin excludes an important complementary family.
At $C=I$, reverse the central signs on a nonempty subset $R\subset T$,
leaving the other groups in $F$ at the central pattern. Let
$k_e\in\{0,\ldots,7\}$ count the reversed links over edge $e$.
Under the same common selected-link rotation, the fourteen source
products on that edge are $2k_e$ copies of $R_n(t)$ and
$14-2k_e$ copies of $-R_n(t)$. The first are in the inner logarithm
region and contribute to the pivot derivative.

\begin{lemma}[Exact active pivot motion]
\label{lem:active-pivot-motion}
On a common sufficiently small interval in $t$, the actual pivot is
\[
 B_e(t)=\exp(-a_etE_n),\qquad a_e=\frac{k_e}{8-k_e},
 \qquad B'_e(0)=-a_eE_n,\quad B''_e(0)=-a_e^2I.
\]
Every other pivot is independent of this selected motion.
\end{lemma}
\begin{proof}
At the candidate pivot the positive relative sources have logarithm
$(1+a_e)tE_n$ in the inner cutoff region, while the negative ones
remain outside the support. The completed equation reduces to
$-a_e+(2k_e/16)(1+a_e)=0$. Its solution is the displayed $a_e$.
Global uniqueness of the pivot identifies this candidate with the
actual solution. There are only eight occupancies, so the interval
can be chosen common. Source incidence excludes all other pivots.
\end{proof}
A zero logarithm at the centre does not make an inner active source
an inactive source. Its first derivative and the pivot acceleration
both enter the completed Hessian.

Define the rational curvature cost
\begin{equation}\label{eq:active-debit}
 d(a)=\max\{2a^2+4a,4a^2\},\quad
 \ell(k)=24k+6d\!\left(\frac{k}{8-k}\right),\quad
 \mathcal L(R)=\sum_e\ell(k_e).
\end{equation}
The values for $k=0,\ldots,7$ are
\[
 0,\quad1356/49,\quad172/3,\quad2268/25,\quad132,
                   \quad580/3,\quad360,\quad1344.
\]

\begin{theorem}[Active-source curvature and original-clock repair]
\label{thm:active-repair}
At every prescribed $R$-pattern, uniformly in all unprescribed fields,
\[
 S''(0)\le-312+\mathcal L(R).
\]
For fixed $\eta>0$, consider all nonempty subsets with
$\mathcal L(R)\le312-\eta$. A union $D_\eta^{\rm act}$ of sufficiently
small closed inner collar charts around those patterns has the floor
\cref{eq:uniform-repair-template} with $c_\eta=\eta/26208$, uniformly
in even $L\ge32$ and compatible hidden sets, provided $F$ is retained.
The charts can be chosen outside every all-sources-inactive event of
\cref{thm:matrix-repair}.
\end{theorem}
\begin{proof}
First hold pivots artificially fixed only to compare derivatives.
Two-selected central plaquettes still cancel the common rotation.
Reversing a selected sign changes a fixed one-selected second
derivative by at most four, at six incident plaquettes, so this
auxiliary derivative is at most $-312+24|R|$.

Restore every actual pivot term using \cref{lem:active-pivot-motion}.
A pivot of speed $a$ and acceleration magnitude $a^2$ contributes at
most $2a^2+4a$ when the other moving parallel link is a selected free
link, and at most $2a^2$ if it is fixed. For two moving pivots with
speeds $a,b$, all new terms are at most
\[
 2a^2+2b^2+4ab\le4a^2+4b^2\le d(a)+d(b).
\]
These product-rule bounds use only unitarity of the remaining links
and include their noncommuting configurations. Each pivot meets six
plaquettes; adding the costs gives \cref{eq:active-debit}. The
independent-coordinate speed squared remains $1638$: solved pivots
are not extra coordinates in the product metric.

There are finitely many admissible patterns on the fixed collar.
Their unit action Hessians are at most $-\eta/1638$. Uniform continuity
on the finite compact input closure gives nearby chart curvature at
most $-\eta/(2\cdot1638)$. Unlike the inactive case, the inverse-pivot
amplitude is not constant. Its logarithmic second derivatives are
bounded independently of volume, as are the coordinate-Haar terms.
Absorbing these actual order-one terms at large $\beta$ gives the same
measured local coefficient $c_\eta=\eta/26208$. The fixed finite
cover, global resolvent, and marginal domination prove the floor.

Every centre has a positive inner source with $C_e^{-1}P_{e,j}=I$.
Shrink its chart so this source remains inside chordal radius $\rho$;
then it violates \cref{eq:inactive-margin} for every $\sigma>0$.
This does not delete its nonconstant Jacobian. \src{F13}{V17}
\end{proof}
Eleven reversals on distinct coarse edges give
$S''\le-312+11(1356/49)=-372/49$, while five reversals on one edge
give $S''\le-312+580/3=-356/3$. The first margin yields
$c_\eta=31/107016$. Occupancies six and seven fail this particular
sufficient cost test; that failure does not exclude other negative
directions. General transition-cutoff configurations, high occupancy
loads, and other noncentral patterns remain outside this certified
family. The chart radii and thresholds are existential, not certified
simulation parameters.

\section{A quantitative local-to-global interface criterion}
\label{sec:weighted-patch-criterion}
Local killed floors do not automatically imply the Poincar\'e
inequality of an unrestricted interface. A useful alternative
criterion keeps all interactions inside actual conditional patches
and makes their required strength explicit.

For the exact unrestricted interface law, let
$h_i=I-P_i$ and $\mathsf H=\sum_i h_i$ in the original rate-one clock.
Assume a commutation graph of degree at most $\Delta$, with owner
range $R_*$: nonadjacent $h_i,h_j$ commute, and adjacent owners differ
by at most $R_*$ in maximum norm. For each graph edge choose
$0\le c_{ij}\le1$ with
\[
 \{h_i,h_j\}\succeq-c_{ij}(h_i+h_j).
\]
The universal value one follows from $(h_i+h_j)^2\succeq0$ and
$h_i^2=h_i$. Smaller values require actual pair estimates.

For a nonnegative finitely supported owner weight $a$, put
\begin{align*}
 m_1&=\sum_y a(y),&m_2&=\sum_y a(y)^2,\\
 d_a(\delta)&=\tfrac12\sum_y|a(y)-a(y-\delta)|^2,
 & A_x&=\sum_i a(\operatorname{owner}(i)-x)h_i,\\
 D_a&=\max_i\sum_{j:\{i,j\}\in E}c_{ij}
               d_a(\operatorname{owner}(i)-\operatorname{owner}(j)).
\end{align*}
Periodize on a sufficiently large compatible torus, without a second
wraparound overlap. Only the deterministic weights are translated;
no translation invariance of the measure is assumed.

\begin{criterion}[Weighted conditional patch squares]
\label{crit:weighted-patches}
Suppose every $A_x$ has gap at least $\gamma_a>0$ in its actual
conditional patch law, uniformly over all admitted compact exterior
and retained coarse data. Then
\begin{equation}\label{eq:weighted-patches}
 \mathsf H^2\succeq\frac{m_1\gamma_a-D_a}{m_2}\mathsf H.
\end{equation}
If $m_1\gamma_a>D_a$, the coefficient is a volume-uniform lower bound
on $\gap\mathsf H$. For a rate-one patch gap $\gamma_\ell$ and
$\alpha_a=\min\{a(y):a(y)>0\}$, it suffices that
\[
 \gamma_\ell>\tau_a:=\frac{D_a}{\alpha_a m_1},\qquad
 \gap\mathsf H\ge\frac{\alpha_a m_1}{m_2}(\gamma_\ell-\tau_a).
\]
\end{criterion}
\begin{proof}
With $q_a(\delta)=\sum_y a(y)a(y-\delta)=m_2-d_a(\delta)$,
\[
 \sum_x A_x=m_1\mathsf H,\qquad
 \sum_x A_x^2=m_2\mathsf H+
       \sum_{i<j}q_a(\operatorname{owner}(i)-\operatorname{owner}(j))
                                                  \{h_i,h_j\}.
\]
Subtracting gives
$m_2\mathsf H^2-\sum_xA_x^2\succeq-D_a\mathsf H$;
nonedges contribute positively, and the edge row sum bounds the loss.
Positive compact conditional densities and positive rates on each
patch identify the kernel with its exterior-measurable functions.
The fibrewise gap and spectral theorem give
$A_x^2\succeq\gamma_a A_x$. Sum and divide by $m_2$.
The full generator has only constant functions in its kernel.
Finally $A_x\succeq\alpha_a\sum_{i\in B_x}h_i$ yields
$\gamma_a\ge\alpha_a\gamma_\ell$ by the conditional variational
principle. No operator inequality has been squared by a monotonicity
assumption. \src{F12}{V63}
\end{proof}

\begin{corollary}[An explicit four-dimensional finite-size threshold]
\label{cor:weighted-threshold}
Let $\ell$ be divisible by four, with
$\ell\ge4$, $\ell>2R_*$, $\ell\ge R_*^2$, and sufficiently large
ambient owner side at least $2\ell+2R_*$. Define
\[
 C_*=32\Delta(R_*^2+32R_*).
\]
If all the actual rate-one cube-patch gaps, with arbitrary compact
exterior, satisfy $\gamma_\ell>C_*\ell^{-3/2}$, then
\begin{equation}\label{eq:patch-explicit-floor}
 \gap\mathsf H\ge\ell^{-1/2}
                   (\gamma_\ell-C_*\ell^{-3/2})>0.
\end{equation}
\end{corollary}
\begin{proof}
On $\{1,\ldots,\ell\}^4$ use the floor-plus-pyramid weight
\[
 a_\ell(y)=\ell^{-1/2}
       +\ell^{-1}\min_{1\le b\le4}\{y_b,\ell+1-y_b\},
\]
and set it to zero elsewhere. Then $\alpha_a\ge\ell^{-1/2}$,
$m_2\le m_1$, and $m_1\ge\ell^4/64$. For
$\norm\delta_\infty\le R_*$, differences in the common interior
are at most $R_*/\ell$. The symmetric difference of the cubes has
at most $8R_*\ell^3$ sites, on which the nonzero weights are at most
$2\ell^{-1/2}$. Hence
\[
 d_a(\delta)\le(R_*^2/2+16R_*)\ell^2,
 \qquad D_a\le\Delta(R_*^2/2+16R_*)\ell^2.
\]
Insert these bounds in \cref{crit:weighted-patches} and use
$m_1/m_2\ge1$ after the positive difference is established.
\end{proof}
A bound $\gamma_\ell\ge c_\beta\ell^{-p}$ with $p<3/2$ would therefore
suffice at a sufficiently large fixed patch size and fixed $\beta$.
The available surface-order certificate
$\exp[-A(\beta+1)\ell^3]$ does not verify this threshold. It is a
lower certificate, not an upper estimate on the actual gap.

\subsection{The remaining coverage and physical-transport requirements}
The repair events of \cref{thm:nonwrapping-repair,thm:matrix-repair,thm:active-repair}
are actual visible events with all hidden fields unrestricted. They
supply more than a bare rarity estimate, but not every all-exterior
conditional patch gap required by \cref{crit:weighted-patches}.
A finite configuration-space cover on one fixed patch cannot be
replaced by a union of an increasing number of spatial repair events
without a new localization or assembly estimate. Such union-avoidance
inferences are explicitly obstructed in \src{F12}{V79}.

The remaining auxiliary task is therefore to control the complementary
source and curvature cases, or to prove a different complete
same-law patch inequality which crosses the displayed threshold.
Neither the negative hidden-diffusion result nor the sufficient patch
criterion excludes another heat-bath mechanism. Even an unrestricted
volume-uniform auxiliary gap would require cofinal coupling control
and a proved comparison to the physical reflected transfer before it
could enter \cref{hyp:contraction}. This is the distinction maintained
by the physical endpoint in the following chapter.

\chapter{The physical spectral endpoint}
\label{ch:spectral}
The auxiliary results of \cref{ch:compact-repair} remain upstream of
this chapter. The full hidden diffusion can be metastable even though
specific visible interface events have uniform killed heat-bath floors.
Neither generator is identified with the physical reflected transfer by
its name or its invariant density. A physical application requires the
same-law source identification and rate-calibrated form comparison
specified below. No such comparison follows solely from the pivot
convexity of \cref{prop:pivot-potential}.

\section{The local reflected Hilbert space}
Let $\mathfrak A_+^{\rm loc}$ be the bounded gauge-invariant positive-time cylinder algebra with bounded physical support. For a candidate Euclidean state $\omega$ and reflection $\Theta$, set
\[
(F,G)_{\OS}=\omega((\Theta F)G).
\]
Reflection positivity permits quotienting by the null space and completion. The remaining Euclidean hypotheses must supply a strongly continuous positive-time semigroup $e^{-tH}$, $H\ge0$, with vacuum $\Omega$. Denote the centered local vector by
\[
v(F)=[F]-\omega(F)\Omega.
\]
This construction is conditional on the candidate correlations and axioms. A finite positive measure or a finite positive kernel is not itself the continuum Euclidean theory.

For a finite family $\mathcal F=(F_1,\ldots,F_m)$, define $V_{\mathcal F}c=\sum_ic_iv(F_i)$ and
\begin{equation}\label{eq:localGrams}
G_{\mathcal F}^{0}=V_{\mathcal F}^*V_{\mathcal F},\qquad
G_{\mathcal F}^{\tau}=V_{\mathcal F}^*e^{-2\tau H}V_{\mathcal F}.
\end{equation}
The factor $2\tau$ is deliberate: the delayed Gram is the squared norm after evolution by time $\tau$.

\begin{theorem}[Complete local Gram characterization]\label{thm:local-Gram}
Suppose centered local vectors span a dense subspace of $\Omega^\perp$. For a fixed physical $\tau>0$ and $0<q<1$, the following are equivalent:
\begin{enumerate}[label=\textup{(\roman*)},leftmargin=2.5em]
\item $G_{\mathcal F}^{\tau}\preceq q^2G_{\mathcal F}^{0}$ for every finite local family;
\item $\norm{e^{-\tau H}|_{\Omega^\perp}}\le q$;
\item $H|_{\Omega^\perp}\succeq-\tau^{-1}\log q$.
\end{enumerate}
In particular, this packet implies uniqueness of the vacuum in the reconstructed local Hilbert space.
\end{theorem}
\begin{proof}
For every coefficient vector $c$,
\[
c^*G_{\mathcal F}^{\tau}c
=\norm{e^{-\tau H}V_{\mathcal F}c}^2,
\qquad c^*G_{\mathcal F}^{0}c=\norm{V_{\mathcal F}c}^2.
\]
Thus (i) is the contraction inequality on the dense centered local span and extends by continuity to (ii). The spectral theorem makes (ii) equivalent to (iii). A second zero-energy vector orthogonal to $\Omega$ would contradict the strict lower bound. The reverse implications are immediate. \src{F1}{V1}
\end{proof}

All mixed linear combinations are part of the theorem. Separate diagonal tests on named loops do not suffice: a matrix with diagonal entries $0.6$ and off-diagonal entries $0.3$ has largest eigenvalue $0.9$, despite both diagonal entries being below $0.7$. The complete Gram order is the relevant finite datum.

\section{Passing through changing cutoff null spaces}
At cutoff $j$, use the corresponding source family and write $G_j^0,G_j^{\tau_j}$. Suppose for each fixed family
\begin{equation}\label{eq:Gram-convergence}
G_j^0\to G^0,\qquad G_j^{\tau_j}\to G^\tau,
\qquad \tau_j\to\tau>0,
\end{equation}
and
\begin{equation}\label{eq:Gram-error}
G_j^{\tau_j}\preceq q^2G_j^0+E_j,\qquad \norm{E_j}\to0.
\end{equation}
Then $G^\tau\preceq q^2G^0$. This is simply passage to the limit in every quadratic form, with no inverse Gram matrix. It remains valid when the limiting null space grows.

A lower-frame estimate has a different purpose. A bound $G_{j,m}^0\succeq c_mG_m^{\rm phys}$ with $c_m>0$ prevents a predeclared physical coefficient direction from disappearing and stabilizes generalized eigenvalue comparisons. It is not needed merely to pass a common Loewner inequality to the surviving quotient. Some nonzero limiting local states and a suitable dense algebra are nevertheless indispensable for a nontrivial theory. An absolute error estimate can conceal normalized soft states whose entrance norm vanishes. \src{F1}{V1}; \cite{GT}

\section{The physical clock}
Let $T_j$ be a one-step positive transfer at lattice time spacing $a_j\downarrow0$.

\begin{proposition}[Fixed one-step contraction can collapse the limit]\label{prop:step-collapse}
Under local kernel convergence to a nonzero strongly continuous continuum sector, a regulator-independent bound $\norm{T_j}\le q<1$ on a reducing sector containing a noncollapsed local family is impossible. A sufficient rate-calibrated condition is instead
\begin{equation}\label{eq:physicalrate}
\liminf_{j\to\infty}\frac{-\log\rho_j}{a_j}\ge m>0,
\qquad \rho_j=\norm{T_j|_{\Omega_j^\perp}}.
\end{equation}
Under the same convergence and density hypotheses it yields $H_\infty|_{\Omega_\infty^\perp}\succeq mI$.
\end{proposition}
\begin{proof}
A fixed one-step $q$ gives $\norm{T_j^{\lfloor t/a_j\rfloor}}\le q^{\lfloor t/a_j\rfloor}\to0$ for every $t>0$. The limiting kernels vanish at every positive time, contradicting continuity at zero unless the entrance family collapses. Under \cref{eq:physicalrate}, for each $m'<m$, the powers are eventually bounded by $e^{-m'a_j\lfloor t/a_j\rfloor}$. Pass the corresponding finite Gram inequalities to the limit and let $m'\uparrow m$. \src{F1}{V1}
\end{proof}

For $\rho_j=1-\gamma_j\to1$, the physical rate is $(\gamma_j/a_j)(1+o(1))$, not the raw deficit. The fixed-delay criterion of \cref{thm:local-Gram} avoids imposing invariance on each finite local synthesis range.

\section{Separating the retained and conditional temporal forms}
\label{sec:two-level-transfer}
The conditional covariance identity \cref{eq:totalcov} has an exact
operator counterpart for a physical transfer. Let $r:X\to Z$ be a
retained field map for the physical law $\pi$, with
$\overline\pi=r_*\pi$. Write
\[
\mathsf Jm=m\circ r,\qquad
\mathsf Kf=\E_\pi[f\mid r],\qquad \mathsf K=\mathsf J^*.
\]
Then $\mathsf J$ is an isometry and $\mathsf K\mathsf J=I$. Every
$f\in L^2(\pi)$ decomposes uniquely as
\begin{equation}\label{eq:two-level-decomposition}
f=\mathsf Jm+g,\qquad m=\mathsf Kf,\quad \mathsf Kg=0,
\qquad \norm f^2=\norm m^2+\norm g^2.
\end{equation}
For centered $f$, the retained field $m$ is centered as well.

\begin{proposition}[Exact two-level comparison transfer]
\label{prop:two-level-transfer}
Let $Q$ be a positive self-adjoint Markov contraction on
$L^2(\overline\pi)$, and define $R=\mathsf JQ\mathsf K$.
Then $R$ is a positive self-adjoint Markov contraction and
\begin{equation}\label{eq:two-level-form}
\ip{f}{(I-R)f}=\norm g^2+\ip{m}{(I-Q)m}.
\end{equation}
Its centered deficit is exactly the retained deficit:
\begin{equation}\label{eq:two-level-gap}
1-\norm{R|_{\one^\perp}}
=1-\norm{Q|_{\one^\perp}}.
\end{equation}
\end{proposition}
\begin{proof}
Positivity, self-adjointness, and preservation of constants follow
from $\mathsf K=\mathsf J^*$ and the corresponding properties of $Q$.
Since $R(\mathsf Jm+g)=\mathsf JQm$, orthogonality proves
\cref{eq:two-level-form}. The restriction of $R$ to $\Ker\mathsf K$
is zero, while its restriction to $\Ran\mathsf J$ is unitarily
equivalent to $Q$. This proves \cref{eq:two-level-gap}.
\src{F6}{V56}
\end{proof}

The comparison transfer completely refreshes the conditional detail.
Nevertheless, its slowest nonconstant mode is retained. Thus even
perfect conditional refresh does not remove the need for a coarse
physical estimate.

\begin{criterion}[Two-level physical coercivity]
\label{crit:two-level-physical}
Let $P$ be the actual positive self-adjoint physical transfer on
$L^2(\pi)$. Suppose, for every centered decomposition in
\cref{eq:two-level-decomposition},
\begin{equation}\label{eq:two-level-domination}
\ip{f}{(I-P)f}\ge
 a\norm g^2+b\ip{m}{(I-Q)m},
\qquad a,b\ge0,
\end{equation}
and $\ip{m}{(I-Q)m}\ge\gamma_Q\norm m^2$ for centered $m$.
Then
\begin{equation}\label{eq:two-level-physical-gap}
\ip{f}{(I-P)f}\ge\gamma_P\norm f^2,
\qquad \gamma_P=\min\{a,b\gamma_Q\}.
\end{equation}
For a physical step $\delta$ and $0<\gamma_P<1$, the corresponding
spectral lower bound is
\begin{equation}\label{eq:two-level-physical-mass}
-\delta^{-1}\log\norm{P|_{\one^\perp}}
\ge-\delta^{-1}\log(1-\gamma_P)
\ge\gamma_P/\delta.
\end{equation}
\end{criterion}
\begin{proof}
Insert the coarse inequality in \cref{eq:two-level-domination} and
use the orthogonal norm identity \cref{eq:two-level-decomposition}.
The spectral statement follows from positivity of $P$ and
$-\log(1-x)\ge x$ for $0\le x<1$. \src{F6}{V56}
\end{proof}

This supplies a physical mass along a cofinal family only with the
same convergence, noncollapse, and time identification as in
\cref{prop:step-collapse}. It is enough that both $a_j$ and
$b_j\gamma_{Q,j}$ retain the required order-$\delta_j$ margin.
The actual-law form domination \cref{eq:two-level-domination} is an
input: the definition of $R$ does not identify it with $P$.

\begin{lemma}[Comparison accuracy at the physical deficit scale]
\label{lem:transfer-comparison-debit}
For positive self-adjoint transfers $P,R$ on the same centered carrier,
\begin{equation}\label{eq:transfer-comparison-debit}
1-\norm{P|_{\one^\perp}}
\ge1-\norm{R|_{\one^\perp}}
  -\norm{(P-R)|_{\one^\perp}}.
\end{equation}
Hence a reference deficit $m\delta_j+o(\delta_j)$ is preserved at that
coefficient by an $o(\delta_j)$ transfer error; an unscaled $o(1)$
error alone is insufficient.
\end{lemma}
\begin{proof}
Apply the triangle inequality to the centered operator norms. The
last assertion follows by division by $\delta_j$.
\src{F6}{V56}
\end{proof}

The exact half-slab form in \cref{prop:half-variance} provides one
same-law route to \cref{eq:two-level-domination}. The core--bad-set
criterion \cref{crit:core-smoothing} provides another possible upstream
functional inequality, but only after its energy form has been
identified with the physical transfer or generator. Neither route
licenses a change of law or of physical clock.

\section{Orbit coercivity provides a legitimate static-to-dynamic incidence}
For a finite synthesis $V$ into the form domain of the actual $H$, define
\[
G(s)=V^*e^{-2sH}V,\qquad
K(s)=V^*e^{-sH}He^{-sH}V.
\]
Spectral calculus gives
\begin{equation}\label{eq:orbitidentity}
G'(s)=-2K(s),\qquad
G(\tau)=G(0)-2\int_0^\tau K(s)\dd s.
\end{equation}

\begin{criterion}[Orbit coercivity with a residual]\label{thm:orbit}
Suppose $K(s)\succeq mG(s)-r(s)R$ for almost every $s\in[0,\tau]$, where $m>0$, $R\succeq0$, and $r\ge0$ is integrable. Then
\begin{equation}\label{eq:orbitbound}
G(\tau)\preceq e^{-2m\tau}G(0)
+2\left(\int_0^\tau e^{-2m(\tau-s)}r(s)\dd s\right)R.
\end{equation}
If $R\preceq CG(0)$ and the resulting scalar coefficient is below one, this supplies the delayed Gram contraction.
\end{criterion}
\begin{proof}
Test on a coefficient vector, use \cref{eq:orbitidentity}, multiply the scalar differential inequality by $e^{2ms}$, and integrate. \src{F1}{V2}
\end{proof}

An entrance energy bound alone is not orbit coercivity: the evolving normalized state can concentrate in a slow component. Alternatively, a static source form $\mathbb N_j$ can contribute only after a proved relation
\begin{equation}\label{eq:static-incidence}
G_j^0-G_j^\tau\succeq c_\tau R_j^*\mathbb N_jR_j-E_j
\end{equation}
for the same history law and physical delay. Combining $R_j^*\mathbb N_jR_j\succeq\delta G_j^0$ and $E_j\preceq\varepsilon_jG_j^0$ gives the contraction factor $1-c_\tau\delta+\varepsilon_j$. Whitening a static matrix does not construct \cref{eq:static-incidence}.

\section{Vacuum-corrected Feshbach coercivity}
Let
\[
H=\begin{pmatrix}A&B\\B^*&C\end{pmatrix},\qquad C\succeq\gamma I>0,
\]
and set
\begin{equation}\label{eq:Feshbach}
S=A-BC^{-1}B^*,\quad \mathcal R=C^{-1}B^*,\quad
Jx=(x,-\mathcal Rx),\quad M=J^*J=I+\mathcal R^*\mathcal R.
\end{equation}
Then $\Ker H=J\Ker S$. If $\Omega=Jx_0$ is a normalized vacuum, the correct head complement is $x_0^{\perp_M}$, not the complement of a guessed bare constant.

\begin{criterion}[Vacuum-corrected effective head]\label{thm:Feshbach}
If $\Ker S=\C x_0$ and
\[
\ip{x}{Sx}\ge\delta\ip{x}{Mx}\qquad(x\perp_Mx_0),
\qquad\delta>0,
\]
then
\begin{equation}\label{eq:Feshbachgap}
H|_{\Omega^\perp}\succeq\frac{\delta\gamma}{\delta+\gamma}I.
\end{equation}
The same estimate applies on an orbit set when the head inequality is required only for the actual centered head coordinates in that set.
\end{criterion}
\begin{proof}
Write a centered vector as a graph part plus its innovation, and remove the graph-vacuum component in the $M$ metric. If $a$ is the norm of the centered graph part and $b$ the innovation norm, Schur completion gives energy at least $\delta a^2+\gamma b^2$, while the centered norm is at most $a+b$. Weighted Cauchy--Schwarz yields
\[
(a+b)^2\le(\delta a^2+\gamma b^2)(\delta^{-1}+\gamma^{-1}),
\]
which proves \cref{eq:Feshbachgap}. \src{F1}{V3--V4}
\end{proof}

\subsection{A retained infrared obstruction survives fine coercivity}
The Feshbach criterion has a converse obstruction at the level of a
nonnegative transfer form. It explains why a large detail floor cannot
replace the head condition.

\begin{proposition}[Massless Schur edge]
\label{prop:massless-Schur-edge}
On a declared centered carrier, let the bounded nonnegative form
$L=I-P$ have blocks
\[
L=\begin{pmatrix}A&B\\B^*&C\end{pmatrix},\qquad C\succeq cI>0,
\qquad S_L=A-BC^{-1}B^*.
\]
If there are unit retained vectors $x_n$ with
$\ip{x_n}{S_Lx_n}\to0$, then $L$ has no positive lower spectral bound.
For a family, such a sequence rules out a common positive
dimensionless lower bound on its respective centered carriers.
\end{proposition}
\begin{proof}
Lift $x_n$ by the minimizing detail:
$\xi_n=(x_n,-C^{-1}B^*x_n)$. Schur completion gives
$\ip{\xi_n}{L\xi_n}=\ip{x_n}{S_Lx_n}$, while
$\norm{\xi_n}\ge1$. The Rayleigh quotients of the normalized lifts
therefore tend to zero. \src{F6}{V57}
\end{proof}

For a cofinal physical application, the vanishing Rayleigh quotient
must be compared with the step duration using
\cref{eq:two-level-physical-mass,eq:physicalrate}. A dimensionless
deficit tending to zero can still be of order $\delta_j$ and therefore
be compatible with a positive physical mass. The obstruction is not
used without this rate comparison.

If $C\succeq0$ has a null space, positivity first forces
$B n=0$ for every $n\in\Ker C$: test the quadratic form on
$(x,tn)$ for both signs of real $t$, and then on $(\ii x,tn)$.
Only when $C$ has closed range and a positive lower bound on its
orthogonal complement may the same completion use $C^+$.
Gauge directions must therefore be removed or retained in the declared
physical carrier before inversion. This restriction is the transfer-form
analogue of the transverse inverse convention in \cref{ch:tower}.

The Casimir analysis gives finite ultraviolet floors and exploits the finite representation bandwidth of a Wilson plaquette. Nested cuts shield an inner Casimir region from a remote router, leaving mixed terms in a finite collar. These are useful exact reductions. A total magnetic norm extensive in volume is not thereby converted into a cofinal infrared invariant; growing-head transport needs its collective tail and mixed-stock bounds. \src{F1}{V3--V6}

\section{Common-time source transport and signed return}
Let $J$ be an isometry from an old centered source range into a refined range, and let $T_X,T_Y$ be transfers at the same physical time. Set $P=JJ^*$ and
\[
M_1=J^*T_YJ,\quad M_2=J^*T_Y^2J,\quad
C=M_1-T_X,\quad L=(I-P)T_YJ.
\]
Then the exact source--detail identity is
\begin{equation}\label{eq:transport-defect}
T_YJ-JT_X=L+JC,\qquad
(T_YJ-JT_X)^*(T_YJ-JT_X)=\mathbb V+C^2,
\end{equation}
where $\mathbb V=L^*L=M_2-M_1^2\succeq0$.

If the refined tail block $D$ obeys $D\preceq dI$ with $d<q$, the Schur complement gives the sufficient threshold condition
\begin{equation}\label{eq:transport-threshold}
qI-M_1-\frac{\mathbb V}{q-d}\succeq0
\quad\Longrightarrow\quad T_Y\preceq qI.
\end{equation}
The compression mismatch $C$ is signed. A negative direction can increase the threshold reserve, and replacing it everywhere by $\norm C I$ loses that information. This is the precise common-time alternative to a generic norm-only cutoff transport. \cite{GT}

For a simpler positive block transfer $T=\left(\begin{smallmatrix}A&B\\B^*&D\end{smallmatrix}\right)$, $\norm A=a<1$, $\norm B=b$, and $\norm D\le1-\mu$, the condition $b^2<(1-a)\mu$ gives
\begin{equation}\label{eq:2x2-transfer}
\norm T\le\frac{a+1-\mu+\sqrt{(a-1+\mu)^2+4b^2}}2<1.
\end{equation}
Thus a source-head contraction needs an explicit complement reserve and mixing bound. A finite bank cannot represent the full gap by itself.

\section{A complete effective-action terminal criterion}
For a compact effective Gibbs interaction $\Phi=\{\Phi_X\}$, define the mixed rectangular oscillation $\partial_{ef}\Phi_X$ as the supremum of
\[
|\Phi_X(u_e,u_f)-\Phi_X(u'_e,u_f)-\Phi_X(u_e,u'_f)+\Phi_X(u'_e,u'_f)|
\]
over the displayed variables and the remaining exterior. Put
\[
\kappa_{ef}(\Phi)=\sum_{X\ni e,f}\partial_{ef}\Phi_X,
\qquad
\mathfrak D_\mu(\Phi)=\sup_e\sum_{f\ne e}e^{\mu d_E(e,f)}
\frac{e^{\kappa_{ef}(\Phi)}-1}{2}.
\]

\begin{criterion}[Weighted Dobrushin terminal ball]\label{thm:Dobrushin}
If $\mathfrak D_\mu(\Phi)=\alpha_\mu<1$ for some $\mu>0$, then the complete Gibbs specification has a unique infinite-volume state and
\begin{equation}\label{eq:Dobrushin-cov}
|\Cov(F,G)|\le
\frac{4|S_F|\norm F_\infty\norm G_\infty}{1-\alpha_\mu}
 e^{-\mu\dist(S_F,S_G)}
\end{equation}
for bounded cylinder observables. If that state is reflection positive at terminal physical spacing $a_*$ and the gauge-invariant cylinders generate its physical OS space, its Hamiltonian gap is at least $\mu/a_*$.
\end{criterion}
\begin{proof}[Proof structure]
A change in exterior link $f$ tilts the conditional density at $e$ by a function with oscillation at most $\kappa_{ef}$. The conditional influence is bounded by $(e^{\kappa_{ef}}-1)/2$. Its weighted row norm is below one, so the comparison matrix has convergent series $\sum_{n\ge0}C^n$ with weighted norm at most $(1-\alpha_\mu)^{-1}$. The coupling comparison gives \cref{eq:Dobrushin-cov}. Exponential decay under physical time translations, reflection positivity, and the dense local core imply the spectral bound. \src{F7}{V82--V83}
\end{proof}

The effective interaction here is complete, including its generated quasi-local terms and retained labels. A regular conditional fibre, an unevaluated reset reference, or a truncation without a tail estimate is not the $\Phi$ in this theorem. The rare-defect and bridge-resolvent developments give alternative stability criteria around already gapped regular kernels; they still require their selected-path, response, reflection, and physical-clock hypotheses. \src{F7}{V84--V91}

\section{An explicit strong-electric region}
For the finite-volume $\SU(3)$ heat-kernel/Wilson slab, let $\epsilon_\kappa=\norm{K_\kappa-1}_\infty$. The source proves that
\[
2\epsilon_\kappa+12\tanh\eta<1
\]
implies a volume-uniform neutral-sector contraction
\begin{equation}\label{eq:strong-electric}
q\le\frac{1-12\tanh\eta-
\sqrt{(1-12\tanh\eta)^2-4\epsilon_\kappa^2}}
{2\epsilon_\kappa}<1,
\end{equation}
with its continuous interpretation at $\epsilon_\kappa=0$. In particular, $\kappa\ge4$ and $0\le\eta\le1/18$ give $q\le9/13$. The temporal influence and the twelve same-slice magnetic neighbors provide the comparison row, and the Doob transform identifies the resulting boundary contraction with the physical slab transfer. \cite{GT}

This supplies a concrete terminal domain. It does not establish that the weak-coupling continuum trajectory enters that domain after exact renormalization, nor that all generated interactions satisfy the same bound. A proof using this terminal route must establish that bridge.
\sourcebasis{Local Gram, quotient passage, physical time, and orbit coercivity: \src{F1}{V1--V2}. Vacuum-corrected Feshbach and ultraviolet shielding: \src{F1}{V3--V6}. Common-time signed transport and strong-electric criterion: \cite{GT}. Complete-action terminal clustering: \src{F7}{V82--V91}.}

\chapter{The integrated conditional closure theorem}
\label{ch:closure}
\section{The sufficient packet}
The preceding results isolate a short logical endpoint. It is useful to state it without requiring any one historical sufficient mechanism for the final contraction.

\begin{hypothesis}[Common compact construction]\label{hyp:compact}
For the chosen gauge group and declared local-vacuum branch, there is a cofinal regulator family of actual finite renewal-loaded compact gauge laws with compatible source writers and exact blocking. All generated interactions, Haar factors, quotient data, calibration marks, and required sector labels are retained. The comparison objects used in estimates are connected to these laws by explicit proved identities or controlled defects.
\end{hypothesis}

\begin{hypothesis}[Renormalized local continuum theory]\label{hyp:continuum}
The family has a physically calibrated renormalized trajectory and convergent local gauge-invariant Euclidean correlations. The limiting correlations satisfy the required Euclidean symmetry, regularity, reflection positivity, and continuity properties for the stated Yang--Mills reconstruction. The coupling is fixed by a nonzero physical normalization at a declared scale, rather than by an uncontrolled over-weak bare limit. The resulting theory is nontrivial, and bounded local gauge-invariant observables generate its physical Hilbert space.
\end{hypothesis}

\begin{hypothesis}[Noncollapsed local probes]\label{hyp:noncollapse}
The centered local source banks have well-defined limiting entrance and delayed Gram matrices. At least the specified nontrivial local states have nonzero limiting norm; lower frames are supplied where physical coefficient directions are required to survive. The source identification and required non-Gaussian response packet establish the intended interacting Yang--Mills branch rather than a zero or purely auxiliary theory.
\end{hypothesis}

\begin{hypothesis}[Uniform physical contraction]\label{hyp:contraction}
There exist one physical $\tau>0$ and one $0<q<1$ such that for every finite centered bounded-support local source family,
\[
G_{j,\mathcal F}^{\tau_j}\preceq q^2G_{j,\mathcal F}^0+E_{j,\mathcal F},
\qquad \tau_j\to\tau,\qquad \norm{E_{j,\mathcal F}}\to0.
\]
The same $q$ applies to every family, although convergence rates may depend on the fixed family.
\end{hypothesis}

\begin{theorem}[Conditional local-vacuum Yang--Mills closure]\label{thm:master}
Under \cref{hyp:compact,hyp:continuum,hyp:noncollapse,hyp:contraction}, the reconstructed local Yang--Mills theory has a unique vacuum and
\begin{equation}\label{eq:master-gap}
\boxed{\displaystyle
H_\infty|_{\Omega_\infty^\perp}\succeq m_*I,
\qquad m_*=-\frac{\log q}{\tau}>0.}
\end{equation}
\end{theorem}
\begin{proof}
The first three hypotheses give the stated nontrivial local OS theory and its dense centered local source span. For each fixed family, pass the Loewner inequality of \cref{hyp:contraction} to the limiting Grams using \cref{eq:Gram-convergence,eq:Gram-error}. It holds with the same $q$ for every finite family. Apply \cref{thm:local-Gram}. No inverse entrance Gram, fixed source dimension, or independent auxiliary generator is needed. The upstream construction and renormalization assumptions identify the Yang--Mills application of this local closure implication. \src{F1}{V1}
\end{proof}

\section{Established compact inputs and their limits}
The explicit compact and local analytic acquisitions can be separated
from the hypotheses of \cref{thm:master}. The smooth root-path completion
has a positive finite-volume density and exact scores on its entire
compact carrier (\cref{cor:compact-product}). Its local germ has a
moving stationary section, a complete measured coefficient, and an
$O(\beta^{-1/2})$ volume-uniform one-loop Hessian remainder on the
original controlled box (\cref{thm:uniform-local-loop}). The actual
one-step transverse coefficient is an analytic differentiated integral,
while flat-background Haar cancellation is exact at every momentum
(\cref{thm:theta-integral,thm:flat-bloch}).

In the compact spectral problem, the reported finite stationary-orbit
certificates imply both the local tube inequality and the hidden
metastability bound (\cref{thm:actual-tube-gap,thm:hidden-metastability}).
These conclusions are compatible. The localized marginal comparison
in \cref{prop:well-removal} does not use a hidden gap. The same-law
global resolvent and marginal variance domination yield volume-uniform
local repair floors for the specified non-wrapping, noncentral, and
active-source events. None of those events is claimed to cover the
whole interface. \Cref{crit:weighted-patches} is an alternative complete
conditional-patch implication, whose all-exterior gap input has not
been supplied by the local cylinders.

Thus neither a new finite block nor a local repair floor changes the
terminal theorem into an unconditional conclusion. They reduce the
remaining obligations to compatibility, complete compact control,
cofinal construction, and physical contraction on the same branch.

\section{A conditional construction route through the Gaussian tower}
The explicit Gaussian construction contributes to
\cref{hyp:continuum} through the following sufficient packet:
\begin{enumerate}[label=\textup{(R\arabic*)},leftmargin=3em]
\item A compatible exact compact blocking realizes the tower as its
principal Gaussian family, with a common analytic domain, declared
quotient and sector transport, and control of the complete transported
vertices. The one-step smooth construction of \cref{cor:compact-product}
provides a concrete candidate carrier, but its common germ does not by
itself prove this all-level compatibility.
\item The fluctuation covariance is placed in the representative of the
vertex banks, with common bounds on the growing alias trace and its
external derivatives, including the treatment of the toron region.
\item The full-torus complete-principal response has the consistently
normalized upward coefficient $b^*=b_{\rm YM}$. It is evaluated either in
the harmonic background of the exact blocked action or in the fixed
background together with \cref{eq:background-Cesaro}. An equivalent
SCIC/ADMC comparison is sufficient; equality of corner cards is not.
\item Nonlinear and stable-coordinate estimates remain uniform along the
same family. The degree-faithful forest growth condition, complete
four-mark source activities, and local current jets control the measured
Jacobian response. Owner bounds, local sector control, and reset or
boundary terms satisfy their common regulator hypotheses. The principal margin used in remainder protection
belongs to that same physically matched response.
\item The resulting local correlations have the compactness, regularity,
reflection, continuity, and nontriviality required by
\cref{hyp:continuum,hyp:noncollapse}.
\end{enumerate}
Under (R1)--(R4), the fixed-point response criterion, the background
comparison where used, the owner remainder, the stable cone, and the
Ces\`aro recursion provide the intended weak-coupling marginal asymptotics
and error bounds. Condition (R5) remains the continuum-existence handoff.
The exact finite representative map supplies an algebraic part of (R2);
it does not supply its uniform analytic estimates or (R3). None of these
five requirements replaces the fixed-physical-time contraction in
\cref{hyp:contraction}.

A spectral route can then use the complete-action terminal ball, an
actual-law response/capacity bound with same-generator incidence, the
two-level form criterion, the orbit-Feshbach criterion, or direct local
Gram estimates. A finite source
head requires its complement and mixed return. An effective terminal
action requires a demonstrated flow into its contraction domain. These
are alternative sufficient mechanisms for the same physical endpoint.

\subsection{The marked local input to the compact continuation}
The analytic part of (R4) has a more explicit internal structure than a
single unspecified remainder estimate. Exact blocking supplies the
conditional differential and its retained unit sector
(\cref{thm:exact-block-differential}). Measured field-redefinition
tangents pass to the retained quotient by
\cref{thm:measured-redundancy}, through the conditional current of the
same law. The Jacobian Hessian spends the third jet of that current
(\cref{prop:Jacobian-third-jet}), and the normalized score hierarchy
identifies the connected cumulants which must be controlled.

Once the actual compact source representation, local jet bounds,
analytic radii, and contour activities in
\cref{crit:nonlinear-current-gate} are supplied uniformly, the
four-mark theorem provides those current and Jacobian estimates. The
Gaussian covariance bounds and degree-faithful forest census are
upstream inputs to this marked representation, not substitutes for
it. In particular, this chain does not infer a uniform conditional
four-point theorem from a quadratic vacuum Hessian or from pair decay.

On the spectral side, \cref{crit:two-level-physical} separates the
conditional detail estimate from the retained physical deficit, while
\cref{prop:massless-Schur-edge} records the obstruction when the retained
form remains infrared soft. The core--bad-set compiler can supply an
upstream energy bound only with its physical-law, tensorization, and
incidence assumptions. These refinements leave the terminal closure
packet unchanged: construction, noncollapse, and a common physical
contraction must ultimately hold for the same theory.

\section{Obstructions that determine the admissible route}
\subsection{The massless balanced baseline}
For the Gaussian bridge analyzed in the Grand Tensor, both the reflected head and its balanced reference can equal $2\beta/L$ despite the model being massless. A positive raw balanced ratio is therefore not a gap test. The meaningful quantity is the reference-subtracted normalized excess. In the scalar Mehler calibration,
\[
\delta=\sinh^2(x/2)=\frac{(1-e^{-x})^2}{4e^{-x}},
\qquad
e^{-x}=(\sqrt{1+\delta}-\sqrt\delta)^2.
\]
A strictly positive excess then has a calibrated meaning in that model. The identity is a baseline test, not a universal substitution for the Yang--Mills delayed Gram. \cite{GT}

\subsection{The Gaussian infrared channel}
For the transverse quadratic colour theory,
\[
\omega_a(k)=a^{-1}\operatorname{arcosh}(1+q(k)/2),\qquad
q(k)=4\sum_{\ell\ {\rm spatial}}\sin^2(k_\ell/2).
\]
At fixed physical torus size $L$, the lowest nonzero frequency is $2\pi/L+O(a^2L^{-3})$. As $L\to\infty$, its half-slab correlation approaches one. A norm-small perturbation of that infrared Gaussian channel cannot open an order-one uniform contraction below one. The ultraviolet Gaussian construction must therefore be connected to a genuinely interacting infrared mechanism; its Gaussian positivity is not enough. \cite{GT}

\subsection{The over-weak toron trajectory}
The Grand Tensor constructs a center-neutral wrapping soft sequence on a specified over-weak bare trajectory. In its parameters, normalized vectors satisfy an energy asymptotic of order
\[
\ip{v_j}{H_j^{\rm phys}v_j}
=\frac{\mu_8+o(1)}{LK_j^{1/3}},\qquad K_j\to\infty,\quad\mu_8>0.
\]
At every fixed physical delay, their squared evolved norms tend to one. This rules out a uniform full finite-torus gap on that trajectory. It does not rule out a correctly renormalized local infinite-volume theory. A replacement trajectory must enforce a nonzero physical coupling condition and its step scaling; a replacement algebra must be declared openly. Wrapping vectors cannot simply be deleted from a theorem that originally quantified over them. \cite{GT}; \src{F1}{V1}

\subsection{The actual compact hidden-well obstruction}
For the specified $\SU(2)$ parent, \cref{thm:hidden-metastability}
excludes a temperature-uniform full hidden diffusion gap. Consequently
that gap cannot be retained as an already acquired premise in a
low-temperature source-transport argument. The obstruction does not
exclude a same-law heat-bath route, a localized marginal comparison,
or a different complete physical criterion. It does exclude inferring
a global hidden gap from normal positivity at a single orbit or from
the small probability of the second well.

\subsection{Finite banks and continuum visibility}
A fixed bank can converge while normalized near-null directions outside it remain soft. A source-complete argument therefore requires either all fixed finite local families with one contraction constant, as in \cref{hyp:contraction}, or a collective tail mechanism in addition to a growing finite head. For example, the source graph-Rellich condition
\[
\norm{(I-\Pi_{j,m})v}^2
\le\eta_m\bigl(\ip{v}{H_jv}+\norm v^2\bigr),\qquad\eta_m\downarrow0,
\]
on the stated branch supplies a quantitative high-energy complement. It is a sufficient separate input, not a consequence of finite-dimensional diagonalization. One should not impose global compactness unnecessarily on a local infinite-volume route when the all-local Gram theorem already applies. \cite{GT}; \src{F1}{V1--V5}

\section{The exact remaining obligations}
\paragraph{Complete Gaussian response.}
The covariance tower, its positive shells, the full-edge vertex banks, and
the finite tree-gauge representative are explicit. The next quantitative
requirement is a higher-level, full-torus shell response with common
control of momentum jets, infrared charts, and normalization. Reuse of
the fixed background additionally requires the averaged comparison with
the harmonic background. Neither an all-level coefficient limit nor its
identification with \cref{eq:bYM} has been established by the level-zero
cards.

\paragraph{Compatible interacting continuation.}
A smooth compact nonabelian block and a volume-uniform local one-loop
remainder are available in \cref{ch:compact-response}. The remaining
requirement is that this compact blocking share the Gaussian principal family
and retain common analytic constants, degree-faithful covariance growth,
local sector and boundary control, and a physically normalized trajectory.
The exact marked polymer representation must control the local current
and score jets on those same charts; the four-mark implication alone
does not establish its activity hypothesis. Measured quotient invariance
is exact at finite regulator, whereas the locality and uniform bounds of
its conditional transport remain quantitative requirements.
The owner theorem protects a matched principal margin but does not create
it or pay for an unbounded background-comparison defect. A uniform tail
for an extensive global sector label is not substituted for the required
local or otherwise explicitly normalized sector control.

\paragraph{Complete compact interface control.}
The local killed floors leave transition-cutoff configurations, large
active-source loads, and other noncentral patterns uncovered. A proof
must either assemble a genuine same-law cover with controlled losses
or establish all-exterior conditional patch estimates strong enough
for \cref{crit:weighted-patches}. The explicit sufficient threshold is
$\gamma_\ell>C_*\ell^{-3/2}$; the available surface-order lower
certificate does not verify it. Neither a growing spatial union of
local repair events nor an unconditional rarity bound supplies this
missing form estimate. Temperature-uniform hidden diffusion is not a
valid shortcut for the parent in \cref{thm:hidden-metastability}.

\paragraph{Continuum existence and physical spectral control.}
The interacting local correlations must have the regularity, reflection,
nontriviality, and continuity needed for reconstruction. The physical
spectral requirement is a uniform contraction on every finite centered
local reflected family at one fixed physical delay, or a fully populated
sufficient mechanism proving it. These remain separate from anomaly
matching. An evaluation of the Gaussian scalar would resolve a major
renormalization interface, not by itself the mass-gap theorem.

\section{Interpretation}
The construction separates the sources of physical content. The common
finite law determines which source responses and conditional operations
belong to the theory. Exact blocking determines the action and its
normalization. The Gaussian tower supplies a controlled multiscale
principal object, while the representative map specifies how its shells
enter the finite Wilson response. The background comparison determines
whether that response is the marginal derivative of the intended
blocked action. The explicit nonlinear fibre then fixes the actual
constraint, moving minimum, and measure derivatives, while the local
one-loop estimate states exactly which uniform norm is controlled.
The compact spectral analysis retains the higher well rather than
confusing its small probability with fast mixing; the visible repair
mechanism works through the original marginal conditionals. Measured field changes remain redundant after blocking,
but their local analytic control is carried by the complete normalized
current, not by an omitted Jacobian. Finally, the two-level form keeps
the retained infrared channel visible, and the complete local Gram
criterion determines what physical spectral decay must mean.

This separation turns finite positivity, analyticity, and computational
identities into usable components without promoting them beyond their
scope. The conditional conclusion is quantitative: once the interacting
continuum law and a noncollapsed fixed-time local contraction are
established on the same calibrated branch, the Hamiltonian has the
strictly positive mass bound of \cref{thm:master}. The remaining work is
localized to the displayed analytic and physical interfaces, rather than
hidden in changes of law, representative, background, or clock.

\appendix
\chapter{Notation and interface ledger}
\label{app:notation}
\section{Principal notation}
\begin{longtable}{L{0.23\textwidth}L{0.70\textwidth}}
\toprule Symbol & Meaning and restriction \\ \midrule
\endfirsthead
\toprule Symbol & Meaning and restriction \\ \midrule
\endhead
$\mu,\pi,\omega$ & Respectively a declared finite law, the actual Perron ground law when so identified, and a candidate limiting Euclidean state. They are not interchangeable. \\
$a_j,\tau$ & Lattice time spacing and a fixed physical delay. The continuum gap uses $\tau$ or the rate $-a_j^{-1}\log\rho_j$. \\
$u_j,\beta_j$ & Physical coupling coordinate and its reciprocal, $u_j=\beta_j^{-1}$, in the declared calibration. A polymer activity denoted by $u$ elsewhere must be distinguished by context. \\
$G_{\mathcal F}^0,G_{\mathcal F}^\tau$ & Entrance and delayed local reflected Grams, with the delay matrix containing $e^{-2\tau H}$. \\
$\mathcal A,\mathcal V,\mathcal K,\mathcal M$ & Positive one-form comparison, defect synthesis, relative return channel, and physical landing metric in \cref{eq:woodbury-objects}. \\
$S,M$ & Feshbach head and its graph metric $M=I+\mathcal R^*\mathcal R$. Vacuum shortening is in the $M$ metric. \\
$\theta_{ij}$ & Normalized coordinate-mark overlap of \cref{eq:marks}; not a raw Fourier coefficient or a representation-dimension probability. \\
$z,K_{\rm HK}$ & Primitive compact analytic activity and the complete weighted forest contraction stock of \cref{eq:KHK}. \\
$\kappa_a,\Pi(I)$ & Normalized joint cumulants and labelled set partitions in \cref{eq:connected-score-hierarchy}; this $\Pi(I)$ is not the transverse momentum projection. \\
$\mathcal O_V,\mathcal J_V,\overline V$ & Measured field-redefinition tangent, pushed vector density, and normalized conditional current in \cref{thm:measured-redundancy}. \\
$\mathsf K_S,\mathsf J_P$ & Conditional expectation and retained pullback for the exact action differential, \cref{thm:exact-block-differential}. \\
$Q_{a,\mu}(u),\Delta$ & Rooted contour-activity majorant and actual contour degree; filling neighbourhoods require $\Delta=80$. \\
$\gamma_P,\gamma_Q,\delta$ & Full and retained dimensionless transfer deficits and physical step duration in \cref{crit:two-level-physical}. \\

$\mathfrak b$ & The calibrated transverse $F^2$ response extractor with its action, colour, and time-orientation convention fixed. \\
ADMC & Averaged determinant matching card, \cref{def:ADMC}. \\
SCIC & Sublinear compact--continuum intertwiner card, \cref{def:SCIC}. \\
$\Pi,\lambda,G^{(j)},S^{(j)}$ & Free transverse projection, lattice dispersion, level-$j$ tower covariance, and transverse inverse kernel. \\
$\Gamma_\mu^*,\Gamma_\ell$ & A diagonal component of the fixed-point lattice sum and, separately, a positive lifted covariance shell. The index and context distinguish them. \\
$\kappa_{\rm W}$ & Finite Wilson relative derivative ceiling below $4/3$, not a stable-flow contraction constant below one. \\
$T_{{\rm g},j},\Xi_j$ & Nested tree-gauge representative and the covariance transported into its full fine-edge vertex banks, \cref{eq:nested-tree-cov}. \\
$b^{\rm fix},b_j^{\rm harm}$ & Fixed fine router background and the harmonic extension of the level-$(j+1)$ coarse background. \\
$\delta_j^{\rm bg}$ & Difference of the consistently normalized harmonic- and fixed-background principal shell coefficients; ADMC requires its Ces\`aro mean to vanish when the fixed route is used. \\
$\widetilde K,B(Y,C),j$ & Smooth compact root-path block, its actual solved pivots, and the product inverse-Haar amplitude; \cref{cor:compact-product}. \\
$G(c),H_c,\widehat Q$ & Minimized local action, eliminated fibre Hessian, and normalized measured one-loop coefficient; the retained Schur form need not have a gap. \\
$d,d_{\rm ch}$ & Hilbert--Schmidt geodesic and chordal distances in the explicit $\SU(2)$ fibre; the cutoff uses $d$, while $d_{\rm ch}\le d$ gives sufficient inactive margins. \\
$\theta,\mathcal I_0$ & Actual one-step transverse $p^2$ coefficient and flat constant response integrand, \cref{thm:theta-integral,thm:flat-bloch}; neither is a physical mass. \\
$L,\mathsf H_\mu,\kappa_D$ & Auxiliary product diffusion, rate-one coordinate heat bath, and killed floor on a visible event; distinct from the physical Hamiltonian $H$. \\
$T,F,M_{\rm vis}$ & Fixed 819-link repair shell, its 10034-group visible collar, and the actual noncentral $3\times3$ curvature certificate. \\
$\mathcal L(R),c_\eta$ & Active-source curvature cost including pivot acceleration, and the measured local diffusion coefficient $c_\eta=\eta/26208$. \\
$\gamma_\ell,\tau_a$ & All-exterior conditional patch gap and sufficient weighted patch threshold, \cref{crit:weighted-patches}; local killed floors do not supply $\gamma_\ell$. \\
\bottomrule
\end{longtable}

\section{Which conclusions require which inputs}
\begin{longtable}{L{0.27\textwidth}L{0.34\textwidth}L{0.29\textwidth}}
\toprule Desired conclusion & Required interface & What does not replace it \\ \midrule
\endfirsthead
\toprule Desired conclusion & Required interface & What does not replace it \\ \midrule
\endhead
Physical Poincar\'e or Witten gap & Actual ground law, full signed sources, physical form coefficient, controlled response/capacity & An unweighted reduced fibre or an arbitrary spatial sampler \\
Strong all-order source analyticity & Exponential fixed-tree grade and a valid complete spatial compiler & A factorial-graded inequality at every arity \\
Compact forest convergence & One analytic radius budget, degree-faithful sum, $8zK_{\rm HK}<1$ & A degree-free estimate that discards a star factorial \\
Measured redundant quotient & Exact pushed current, its divergence, and the same pushed measure & Strict verticality or action transport without its Jacobian \\
Local pushed Jacobian Hessian & Three current jets, four-entry connected bounds, source radii, and corrected contour activity & Pair decay, a determinant expansion alone, or two current jets \\
Core-based physical coercivity & Conditional core gap, restricted smoothing, tensorization, and same-form incidence & A bare core floor or a small unconditional bad-set probability \\
Gaussian ADMC & Exact nested determinant, representative-matched shell response, physical scalar, and sublinear endpoints & A corner mean or a frozen-chart polynomial coefficient \\
Use of fixed-background banks & Complete normalization and a vanishing averaged defect relative to the harmonic prescription & Gauge equivalence of fluctuations or equality at finitely many corners \\
Stable ultraviolet trajectory & Calibrated marginal recursion and invariant stable cone & Strict contraction of the full retained action \\
Interacting continuum theory & Local correlation compactness, regularity, reflection, physical normalization, nontriviality & Convergence of Gaussian kernels alone \\
Two-level temporal gap & Physical form domination plus both conditional and retained deficits at the correct time scale & Complete detail refresh with no retained infrared estimate \\
Physical mass gap & Complete local reflected contraction at a fixed time, or a populated equivalent route & A finite eliminated Hessian floor or positivity on a selected bank \\
Volume-compatible sector control & A specified local, density, fixed-volume, or phase-resolved target & A uniform tail for an extensive total charge without entropy control \\
Smooth compact fibre & Actual globally solved pivots, inverse-Haar factors, and a declared completion & Independent pivot caps or equality of local Taylor germs \\
Local measured one-loop limit & Moving minimum, full density jets, fixed-domain wall comparison, and localized derivative norms & A frozen-background determinant or a bound on the compact complement \\
Local visible repair & Uniform all-hidden curvature, original-clock resolvent, and a visible-only support condition & A force bound restricted to a hidden posterior event \\
Unrestricted auxiliary interface gap & Complete same-law assembly or all-exterior weighted patch gaps above their threshold & One good tube, rare bad fields, or an uncontrolled spatial union \\
\bottomrule
\end{longtable}

\chapter{Corrections retained in the mathematical chain}
\label{app:corrections}
The following restrictions are part of the mathematical hypotheses and conclusions. Each records a distinction needed by the corresponding implication.

\begin{longtable}{L{0.24\textwidth}L{0.49\textwidth}L{0.17\textwidth}}
\toprule Earlier shortcut & Retained replacement & Main location \\ \midrule
\endfirsthead
\toprule Earlier shortcut & Retained replacement & Main location \\ \midrule
\endhead
Raw quartic numerator as normalized MTA & Correct first-connectivity and product-state coefficient-capacity argument, including globally summed rescued suffixes & \cref{ch:trees} \\
All-order factorial-graded MTA as strong MTA & Keep the valid $m!K^m$ bound, withdraw its uniform analytic closure inference; restore exact-Haar assembly & \cref{prop:factorial,thm:Haar} \\
Forest fixed-tree bound without degree factorials & Keep $\prod d_i!$ and use the exact Pr\"ufer sum & \cref{lem:degree-census,thm:forest} \\
A matching Wilson exponent as the Perron law & Prove the incidence equation or repair the actual-law density defect & \cref{prop:law-no-go} \\
Rare bad fields as a uniform conditional estimate & Establish complete-source incidence, restricted smoothing, and capacity & \cref{ch:source,ch:geometry} \\
Finite action-closed update block & Retain exterior plaquettes and control boundary response & \cref{prop:no-finite-block} \\
Strict verticality as the definition of redundancy & Push the measured tangent through its conditional current & \cref{thm:measured-redundancy} \\
Two current jets for the Haar-Jacobian Hessian & Retain the third jet and four-entry normalized cumulants & \cref{prop:Jacobian-third-jet,thm:four-mark-clustering} \\
Degree-eight entropy for filling contours & Use the actual degree-eighty contour graph and the $6401$ threshold & \cref{lem:polymer-entropy-gate} \\
Pointwise centered remainder as a local analytic norm & Supply source-uniform polymer control; the fibre range may be extensive & \cref{thm:exact-block-differential} \\

Conditional fast mixing as full mixing & Retain the slow covariance and the exact two-level temporal form & \cref{eq:totalcov,prop:two-level-transfer} \\
Unscaled small transfer error as physical gap transport & Compare the error to the physical step deficit & \cref{lem:transfer-comparison-debit} \\
Fine floor as retained infrared coercivity & Control the actual Schur head; a soft head has soft minimizing lifts & \cref{prop:massless-Schur-edge} \\
Finite grid signs as global certificates & Supply an analytic quadrature remainder or an actual all-torus proof & \cref{eq:quadrature,ch:oracle} \\
Direct transverse covariance in fixed tree-gauge banks & Push the covariance through the explicit representative map before contraction & \cref{sec:tree-gauge} \\
Fixed background as the harmonic blocked action & Retain the distinct prescriptions and prove the averaged background comparison where required & \cref{crit:background-ADMC} \\
Corner agreement as a tower coefficient & Separate finite algebraic checks, torus integration, and the uniform limiting response & \cref{sec:fixed-background-cards} \\
Frozen endpoint mean as the anomaly & Use an actual nested tower; a nonzero fixed reset mean obstructs that comparison & \cref{ch:anomaly,ch:tower} \\
Covariance nondegeneracy as a physical mass & Retain the massless axis behavior and distinguish inverse-kernel control & \cref{thm:tower-floor,eq:axis} \\
Finite one-step contraction as a continuum gap & Calibrate the physical rate or prove a fixed-delay Gram inequality & \cref{prop:step-collapse} \\
Bare head constant as the vacuum & Short in the exact Feshbach graph metric & \cref{thm:Feshbach} \\
Diagonal probe decay as a full gap & Control every finite mixed local Gram with one constant & \cref{thm:local-Gram} \\
Global sector tail from local action cost & Keep the tensorization no-go and use an explicitly different target & \cref{thm:sector-no-go} \\
Over-weak wrapping trajectory as the physical limit & Impose nonzero physical normalization or declare a different local branch & \cref{ch:closure} \\
Laurent representative at the toron & Use the rational punctured-torus map and its integrable weak derivative bounds & \cref{prop:toron-regularity} \\
Linear nonabelian averaging as a homomorphism & Retain endpoint path dressing and the actual constrained Hessian & \cref{prop:nonabelian-average,eq:actual-constrained-Hessian} \\
Identical local germs as identical compact laws & Distinguish completion choices and compare normalized pushforwards & \cref{scope:completion-germ} \\
Dropping pivot acceleration or Haar terms & Use the complete measured jets, including active-source response & \cref{eq:complete-fourth-jet,thm:active-repair} \\
Nonzero finite flat mean as a bulk constant & Identify the exact Fourier aliases after Haar cancellation & \cref{thm:flat-bloch} \\
Small higher-well mass as fast hidden mixing & Keep its Rayleigh quotient and distinguish marginal removal & \cref{thm:hidden-metastability,prop:well-removal} \\
Periodic shell transplanted into larger volume & Count the true outer plaquettes and embed the complete input closure & \cref{thm:nonwrapping-repair} \\
Fixed-patch cover as a growing spatial union & Require a proved global assembly or all-exterior patch criterion & \cref{sec:weighted-patch-criterion} \\
\bottomrule
\end{longtable}

\chapter{Result concordance and computational scope}
\label{app:concordance}
\section{Mathematical dependencies}
The bibliographic keys F1--F13 identify the successive Yang--Mills
series. F11 contains the nested Gaussian construction, its representative
and background analysis, and the nonlinear local fibre and measured
coefficient. F12 develops the compact completion, local one-loop
comparison, and conditional spectral machinery. F13 develops the
actual orbit geometry, metastability, and fixed-patch visible repair. PAR contains the independent owner-remainder
and sector arguments. The Grand Tensor reference supplies the inherited
finite construction and physical-transfer results. These are unpublished
programme manuscripts, not external literature used to discharge an
unproved input.

The following table records the mathematical result locations. Original
labels, version locations, and, where retained, one-based manuscript
line numbers identify the statements;
the exposition above standardizes notation and displays the implications
used in the conditional chain. Reported computational results remain
distinct from the algebraic consequences proved here.

\section{Principal result map}
\small
\begin{longtable}{L{0.32\textwidth}L{0.12\textwidth}L{0.477\textwidth}}
\toprule Result in this monograph & Source & Original label and line \\ \midrule
\endfirsthead
\toprule Result in this monograph & Source & Original label and line \\ \midrule
\endhead
Half-slab and physical Doob channel (\cref{thm:half}) & \cite{GT} & \nolinkurl{thm:YM-half-slab-Doob}\newline Line 62,008 \\
Actual Perron law incidence (\cref{prop:law-no-go}) & \cite{F1} & \nolinkurl{thm:V21-incidence-equation}\newline Line 13,157 \\
Sequential physical response matrix (\cref{thm:response-matrix}) & \cite{F1} & \nolinkurl{thm:V14-response-matrix}\newline Line 8,721 \\
Physical landing metric (\cref{thm:landing}) & \cite{F1} & \nolinkurl{thm:V17-landing-ledger}\newline Line 10,363 \\
Complete-source capacity factorization (\cref{eq:bad-capacity}) & \cite{F1} & \nolinkurl{thm:V20-incidence-capacity}\newline Line 12,803 \\
Corrected quartic weak-cut argument (\cref{ch:trees}) & \cite{F2} & \nolinkurl{thm:V1-quartic-weak-cut}\newline Line 292 \\
Completed normalized quartic theorem (\cref{thm:quartic}) & \cite{F3} & \nolinkurl{thm:V96-global-quartic-MTA}\newline Line 46,293 \\
Projected-chain degree filtration (\cref{ch:trees}) & \cite{F4} & \nolinkurl{lem:V100-filtered-chain}\newline Line 44,273 \\
Withdrawal of the overstrong all-order grade (\cref{prop:factorial}) & \cite{F5} & \nolinkurl{thm:newV12-correction}\newline Line 5,072 \\
Exact assembled Haar derivatives (\cref{thm:Haar}) & \cite{F5} & \nolinkurl{thm:newV12-exact-Haar-card}\newline Line 5,248 \\
Strong forest projection, not source-occurrence closure (\cref{ch:trees}) & \cite{F6} & \nolinkurl{thm:s6v7-VAFPC}\newline Line 3,135 \\
Connected-carrier activity criterion (\cref{eq:carrier-row}) & \cite{F6} & \nolinkurl{thm:s6v16-carrier}\newline Line 6,944 \\
Exact half-slab conditional-variance form (\cref{prop:half-variance}) & \cite{F6} & \nolinkurl{thm:s6v58-conditional-variance}\newline Line 25,626 \\
Same-carrier half-slab density comparison (\cref{lem:half-density-comparison}) & \cite{F6} & \nolinkurl{lem:s6v58-density-comparison}\newline Line 25,896 \\
Complete normalized score hierarchy (\cref{thm:connected-score-hierarchy}) & \cite{F8} & \nolinkurl{thm:e8v60-all-score}\newline Line 25,780 \\
Conditional core-mean gluing (\cref{lem:core-mean-gluing}) & \cite{F6} & \nolinkurl{lem:s6v97-mean}\newline Line 45,631 \\
Core gap plus restricted smoothing (\cref{crit:core-smoothing}) & \cite{F6} & \nolinkurl{thm:s6v97-bootstrap}\newline Line 45,675 \\
Exact retained-pullback sector (\cref{thm:exact-block-differential}) & \cite{F7} & \nolinkurl{thm:s7v99-pullback}\newline Line 44,998 \\
Centered exact nonlinear remainder (\cref{eq:exact-centered-remainder}) & \cite{F7} & \nolinkurl{thm:s7v99-exact-split}\newline Line 45,085 \\
Measured redundancy through pushforward (\cref{thm:measured-redundancy}) & \cite{F8} & \nolinkurl{thm:e8v58-pushforward-redundancy}\newline Line 25,050 \\
Gaussian measured EOM has no marginal feed (\cref{cor:Gaussian-EOM-no-feed}) & \cite{F8} & \nolinkurl{cor:e8v58-no-Gaussian-feed}\newline Line 25,187 \\
Third current jet in the Haar-Jacobian Hessian (\cref{prop:Jacobian-third-jet}) & \cite{F8} & \nolinkurl{prop:e8v60-third-jet-debit}\newline Line 25,736 \\
Volume-uniform four-mark clustering (\cref{thm:four-mark-clustering}) & \cite{F8} & \nolinkurl{thm:e8v61-four-mark-cluster}\newline Line 26,291 \\
Corrected nonlinear-current activity criterion (\cref{crit:nonlinear-current-gate}) & \cite{F8} & \nolinkurl{thm:e8v63-revised-gate}\newline Line 27,112 \\
Exact two-level comparison transfer (\cref{prop:two-level-transfer}) & \cite{F6} & \nolinkurl{thm:s6v56-two-level}\newline Line 24,778 \\
Two-level physical temporal coercivity (\cref{crit:two-level-physical}) & \cite{F6} & \nolinkurl{thm:s6v56-compiler}\newline Line 24,837 \\
Physical-scale transfer comparison (\cref{lem:transfer-comparison-debit}) & \cite{F6} & \nolinkurl{lem:s6v56-additive-debit}\newline Line 24,898 \\
Massless retained Schur obstruction (\cref{prop:massless-Schur-edge}) & \cite{F6} & \nolinkurl{thm:s6v57-massless-Schur}\newline Line 25,301 \\
Corrected filling-contour entropy gate (\cref{lem:polymer-entropy-gate}) & \cite{F8} & \nolinkurl{prop:e8v63-contour-entropy}\newline Line 27,057 \\
Single analytic radius budget (\cref{ch:forests}) & \cite{F9} & \nolinkurl{lem:n1v1-one-radius}\newline Line 478 \\
Correct regulator-dependent forest power (\cref{eq:forest-beta}) & \cite{F10} & \nolinkurl{thm:newv100-degree-scale}\newline Line 21,200 \\
Averaged determinant matching definition (\cref{def:ADMC}) & \cite{F9} & \nolinkurl{def:n1v14-ADMC}\newline Line 6,600 \\
ADMC marginal implication (\cref{ch:anomaly}) & \cite{F9} & \nolinkurl{thm:n1v14-ADMC}\newline Line 6,617 \\
Sublinear compact--continuum intertwiner (\cref{def:SCIC}) & \cite{F9} & \nolinkurl{def:n1v15-SCIC}\newline Line 7,097 \\
One-dimensional analytic alias cohomology (\cref{thm:alias}) & \cite{F9} & \nolinkurl{thm:n1v17-cohomology}\newline Line 7,755 \\
Exact remaining mean criterion (\cref{eq:edge-scalar}) & \cite{F9} & \nolinkurl{thm:n1v17-edge-SCIC}\newline Line 7,854 \\
Reported all-torus relative margin (\cref{eq:relative-floor}) & \cite{F10} & \nolinkurl{thm:newv97-full-torus}\newline Line 20,399 \\
Normalized covariance-derivative consequence (\cref{thm:covariance-margin}) & \cite{F10} & \nolinkurl{thm:newv100-normalized-margin}\newline Line 21,318 \\
Fixed-point Gaussian matching criterion (\cref{thm:fixedpoint-anomaly}) & \cite{F11} & \nolinkurl{thm:v34-card}\newline Line 12,168 \\
Explicit finite-level block-average covariance (\cref{thm:tower-explicit}) & \cite{F11} & \nolinkurl{thm:v37-explicit}\newline Line 13,356 \\
Positive lattice-sum fixed point (\cref{thm:tower-fixed}) & \cite{F11} & \nolinkurl{thm:v37-fixed-point}\newline Line 13,378 \\
Global transverse floor (\cref{thm:tower-floor}) & \cite{F11} & \nolinkurl{thm:v37-nondegenerate}\newline Line 13,419 \\
Composed harmonic extension (\cref{eq:tower-minimizer}) & \cite{F11} & \nolinkurl{thm:v38-harmonic}\newline Line 13,634 \\
Positive conditional covariance shells (\cref{thm:shells}) & \cite{F11} & \nolinkurl{thm:v38-shells}\newline Line 13,684 \\
Harmonic tree-level shell functional and fixed prescription (\cref{sec:background-prescription}) & \cite{F11} & \nolinkurl{def:v38-shell-functional}\newline Line 13,742; \nolinkurl{dec:v41-background}\newline Line 14,334 \\
Fluctuation-representative dependence (\cref{prop:representative-dependence}) & \cite{F11} & \nolinkurl{prop:v41-gauge-dependence}\newline Line 14,271 \\
Tree-gauge map and finite route agreement (\cref{sec:tree-gauge}) & \cite{F11} & \nolinkurl{def:v41-map}\newline Line 14,296; \nolinkurl{thm:v41-map}\newline Line 14,309 \\
Fixed-background level-zero cards (\cref{eq:corners-fixed}) & \cite{F11} & \nolinkurl{thm:v41-cards}\newline Line 14,357 \\
Averaged background compatibility (\cref{crit:background-ADMC}) & \cite{F9,F11} & SCIC specialization of \cref{def:SCIC}; the unproved background comparison is identified in \nolinkurl{rem:v41-why}. \\
Non-double-counted response identity (\cref{eq:owner-ledger}) & \cite{PAR} & \nolinkurl{thm:exact-ledger}\newline Line 179 \\
Physical principal-margin protection (\cref{thm:remainder}) & \cite{PAR} & \nolinkurl{thm:margin-protection}\newline Line 670 \\
Global sector tensorization obstruction (\cref{thm:sector-no-go}) & \cite{PAR} & \nolinkurl{thm:global-sector-no-go}\newline Line 1,167 \\
Label refinement data processing (\cref{eq:label-TV}) & \cite{PAR} & \nolinkurl{prop:refinement-defect}\newline Line 1,262 \\
Complete local OS Gram criterion (\cref{thm:local-Gram}) & \cite{F1} & \nolinkurl{thm:local-Gram-criterion}\newline Line 231 \\
Passage through changing null spaces (\cref{eq:Gram-error}) & \cite{F1} & \nolinkurl{thm:quotient-stable-passage}\newline Line 382 \\
Physical-time collapse obstruction (\cref{prop:step-collapse}) & \cite{F1} & \nolinkurl{thm:fixed-step-collapse}\newline Line 558 \\
Orbit-tube coercivity (\cref{thm:orbit}) & \cite{F1} & \nolinkurl{thm:V2-orbit-coercivity}\newline Line 988 \\
Exact vacuum metric and Feshbach gap (\cref{thm:Feshbach}) & \cite{F1} & \nolinkurl{thm:V4-vacuum-corrected-gap}\newline Line 2,344 \\
Common-time signed return packet (\cref{eq:transport-threshold}) & \cite{GT} & \nolinkurl{thm:YM-signed-threshold-packet}\newline Line 66,113 \\
Head--tail transfer compiler (\cref{eq:2x2-transfer}) & \cite{GT} & \nolinkurl{thm:YM-regulated-criterion}\newline Line 66,078 \\
Complete effective-action terminal ball (\cref{thm:Dobrushin}) & \cite{F7} & \nolinkurl{thm:s7v83-ball}\newline Line 36,610 \\
Physical mass from terminal clustering (\cref{thm:Dobrushin}) & \cite{F7} & \nolinkurl{cor:s7v83-gap}\newline Line 36,651 \\
Explicit volume-uniform terminal region (\cref{eq:strong-electric}) & \cite{GT} & \nolinkurl{thm:YM-strong-electric-gap}\newline Line 66,036 \\
Local-vacuum conditional closure (\cref{thm:master}) & \cite{F1} & \nolinkurl{thm:local-vacuum-closure}\newline Line 726 \\
Toron regularity (\cref{prop:toron-regularity}) & \cite{F11} & \nolinkurl{prop:v42-not-laurent}; \nolinkurl{lem:v42-derivatives}\newline V42 \\
Nonabelian average obstruction (\cref{prop:nonabelian-average}) & \cite{F11} & V56 \\
Global compact pivot diffeomorphism (\cref{thm:global-pivot}) & \cite{F12} & \nolinkurl{thm:f12v18-pivot}\newline V18 \\
Global product carrier and smooth density (\cref{cor:compact-product}) & \cite{F12} & \nolinkurl{thm:f12v18-global}; \nolinkurl{cor:f12v18-density}\newline V18 \\
Actual pivot residual and strong convexity (\cref{prop:pivot-potential}) & \cite{F12} & \nolinkurl{thm:f12v98-potential}\newline V98 \\
Uniform local one-loop remainder (\cref{thm:uniform-local-loop}) & \cite{F12} & \nolinkurl{thm:f12v12-local}\newline V12 \\
Differentiated transverse coefficient (\cref{thm:theta-integral}) & \cite{F12} & \nolinkurl{thm:f12v5-integral}; \nolinkurl{thm:f12v5-quadrature}\newline V5 \\
Flat identity and Fourier aliases (\cref{thm:flat-bloch}) & \cite{F12} & \nolinkurl{thm:f12v9-flat}; \nolinkurl{prop:f12v9-alias}\newline V9 \\
Actual local orbit-tube inequality (\cref{thm:actual-tube-gap}) & \cite{F13} & \nolinkurl{thm:tube}\newline V1 \\
Second well and hidden metastability (\cref{thm:hidden-metastability}) & \cite{F13} & \nolinkurl{thm:f13v8-minimum}; \nolinkurl{thm:f13v8-metastability}\newline V8 \\
Localized effective-action comparison (\cref{prop:well-removal}) & \cite{F13} & \nolinkurl{thm:f13v9-potential}\newline V9 \\
Same-law diffusion and heat-bath comparison (\cref{thm:global-resolvent}) & \cite{F12} & \nolinkurl{thm:f12v60-convert}\newline V60; resolvent development V77--V80 \\
Non-wrapping original-clock repair (\cref{thm:nonwrapping-repair}) & \cite{F13} & \nolinkurl{thm:f13v14-curvature}; \nolinkurl{thm:f13v14-floor}\newline V14 \\
Locally curved-background transport (\cref{sec:coarse-gauge-repair}) & \cite{F13} & \nolinkurl{thm:f13v15-floor}\newline V15 \\
Noncentral visible matrix criterion (\cref{thm:matrix-repair}) & \cite{F13} & \nolinkurl{thm:f13v16-repair}\newline V16 \\
Active-source pivot and repair (\cref{thm:active-repair}) & \cite{F13} & \nolinkurl{thm:f13v17-curvature}; \nolinkurl{thm:f13v17-repair}\newline V17 \\
Weighted conditional patch criterion (\cref{crit:weighted-patches}) & \cite{F12} & \nolinkurl{thm:f12v63-weighted}\newline V63 \\
Explicit finite-size threshold (\cref{cor:weighted-threshold}) & \cite{F12} & \nolinkurl{thm:f12v63-threshold}\newline V63 \\
\bottomrule
\end{longtable}
\normalsize
\section{Reproducibility and scope}
The normalized coefficient and connected-source arguments depend on
\cite{F1,F2,F3,F4,F5,F6}; the conditional geometry and exact compact
blocking on \cite{F6,F7,F8,F9}; and the finite response oracle on
\cite{F9,F10}. The explicit tower, full-edge vertices, and representative
and background analysis, together with the actual local nonlinear
fibre, are in \cite{F11}. The smooth compact completion, local one-loop
error, global resolvent, and weighted patch criterion are in \cite{F12};
the two-orbit geometry, normalized well comparison, and visible repair
families are in \cite{F13}. The nonlinear remainder and
sector conclusions are in \cite{PAR}. The local physical endpoint is
developed in \cite{F1} and in the Grand Tensor \cite{GT}.

The collar topology, all-torus finite Hessian margin, response banks,
corner identities, compact stationary-orbit Hessians, and non-wrapping
shell incidence data refer to finite computational certificates documented
in these references. Their auxiliary executables are not reproduced in
this monograph, and no independent rerun of the complete certificate
chain is asserted. In particular, the finite alias--slice agreement and
corner means do not constitute a uniform full-torus or higher-shell
calculation. The covariance-transport propositions are algebraic
consequences of the displayed representative constraints;
\cref{crit:background-ADMC} is the explicit SCIC specialization, not a
proof of background independence.

The higher score, measured pushforward, core-gluing, and two-level
transfer statements are proved on their displayed finite or
same-carrier domains. The four-mark result is an analytic implication
from an exact source-dependent activity representation; it is not
claimed as an independently verified bound for every retained Wilson
background. The existing Gaussian tower and response certificates are
retained with their original scope. The actual orbit implications use
the specified finite certificates, not an independently reconstructed
numerical critical point. The local repair proofs use the stated
finite geometric incidence data, together with the displayed analytic
localization and same-law operator comparisons. Their chart radii and
large-$\beta$ thresholds are existential constants unless explicitly
stated otherwise.

No unproved interacting construction, continuum measure, infrared
contraction, or anomaly identification is imported from another sector.
The conditional statements retain those requirements explicitly.

\backmatter

\begin{thebibliography}{F13}

\bibitem[GT]{GT}

A.~P\'elissier, \emph{The Renewal Grand Tensor: Relative Primitive Minimality, Operational Spectralization, and Domain-Specific SM--Spacetime and Arithmetic Loadings}, version V075. Unpublished manuscript, 2026.

\bibitem[F1]{F1}

A.~P\'elissier, \emph{Renewal Yang--Mills mass-gap programme}, series 1, version V135\_f1. Unpublished manuscript, 2026.

\bibitem[F2]{F2}

A.~P\'elissier, \emph{Renewal Yang--Mills mass-gap programme}, series 2, version V100\_f2. Unpublished manuscript, 2026.

\bibitem[F3]{F3}

A.~P\'elissier, \emph{Renewal Yang--Mills mass-gap programme}, series 3, version V100\_f3. Unpublished manuscript, 2026.

\bibitem[F4]{F4}

A.~P\'elissier, \emph{Renewal Yang--Mills mass-gap programme}, series 4, version V100\_f4. Unpublished manuscript, 2026.

\bibitem[F5]{F5}

A.~P\'elissier, \emph{Renewal Yang--Mills mass-gap programme}, series 5, version V100\_f5. Unpublished manuscript, 2026.

\bibitem[F6]{F6}

A.~P\'elissier, \emph{Renewal Yang--Mills mass-gap programme}, series 6, version V100\_f6. Unpublished manuscript, 2026.

\bibitem[F7]{F7}

A.~P\'elissier, \emph{Renewal Yang--Mills mass-gap programme}, series 7, version V100\_f7. Unpublished manuscript, 2026.

\bibitem[F8]{F8}

A.~P\'elissier, \emph{Renewal Yang--Mills mass-gap programme}, series 8, version V65\_f8. Unpublished manuscript, 2026.

\bibitem[F9]{F9}

A.~P\'elissier, \emph{Renewal Yang--Mills mass-gap programme}, series 9, version V100\_f9. Unpublished manuscript, 2026.

\bibitem[F10]{F10}

A.~P\'elissier, \emph{Renewal Yang--Mills mass-gap programme}, series 10, version V100\_f10. Unpublished manuscript, 2026.

\bibitem[F11]{F11}

A.~P\'elissier, \emph{Renewal Yang--Mills mass-gap programme}, series 11, version V100\_f11. Unpublished manuscript, 2026.

\bibitem[F12]{F12}
A.~P\'elissier, \emph{Renewal Yang--Mills mass-gap programme}, series 12,
version V100\_f12. Unpublished manuscript, 2026.

\bibitem[F13]{F13}
A.~P\'elissier, \emph{Renewal Yang--Mills mass-gap programme}, series 13,
version V17\_f13. Unpublished manuscript, 2026.

\bibitem[PAR]{PAR}

A.~P\'elissier, \emph{Renewal Yang--Mills theory: Nonperturbative margin protection, sector tightness, and topological transport}, version V5. Unpublished manuscript, 2026.

\end{thebibliography}

\end{document}

